
\documentclass[UTF8,10.5pt,openany]{ctexbook}
\usepackage[a4paper,left=2.1cm,right=2.1cm,top=2.5cm,bottom=2.2cm,headheight=23pt,headsep=9pt]{geometry}
\usepackage{fontspec}
\usepackage{amsmath}
\usepackage{amssymb}
\usepackage{xeCJK}
\setCJKmainfont{Noto Serif CJK SC}[BoldFont={Noto Serif CJK SC},BoldFeatures={FakeBold=2.0}]
\setCJKsansfont{Noto Sans CJK SC}
\setCJKmonofont{Noto Sans Mono CJK SC}
\setmainfont{Noto Serif}
\setsansfont{Noto Sans}
\setmonofont{Noto Sans Mono}[Scale=0.86]
\setCJKfamilyfont{zhkai}{FandolKai-Regular.otf}\newcommand{\kai}{\CJKfamily{zhkai}}
\usepackage{xcolor}
\definecolor{ink}{RGB}{17,17,17}
\definecolor{grayt}{RGB}{110,110,110}
\definecolor{linet}{RGB}{150,150,150}
\definecolor{accent}{RGB}{20,60,120}
\usepackage{titlesec}
\usepackage{titletoc}
\usepackage{fancyhdr}
\usepackage{enumitem}
\usepackage{longtable}
\usepackage{multicol}
\usepackage{booktabs}
\usepackage{array}
\usepackage{needspace}
\usepackage[most]{tcolorbox}
\usepackage[bookmarks=true,bookmarksnumbered=true,bookmarksopen=true,bookmarksopenlevel=1,hidelinks,unicode]{hyperref}
\linespread{1.30}
\setlength{\parindent}{0pt}
\setlength{\parskip}{2pt}
\pagestyle{fancy}
\fancyhf{}
\renewcommand{\headrulewidth}{0.4pt}
\fancyhead[L]{\footnotesize\sffamily\color{grayt}AI Infra 精读笔记}
\fancyhead[R]{\footnotesize\sffamily\color{grayt}\leftmark}
\fancyfoot[C]{\footnotesize\color{grayt}\thepage}
\renewcommand{\chaptermark}[1]{\markboth{#1}{}}
\renewcommand{\thechapter}{\arabic{chapter}}
\titleformat{\chapter}[display]
  {\normalfont\centering}
  {\footnotesize\sffamily\color{grayt}第\,\thechapter\,篇}
  {4pt}
  {\sffamily\bfseries\fontsize{18}{23}\selectfont}
  [\vspace{3pt}{\color{linet}\rule{0.42\textwidth}{0.5pt}}]
\titlespacing*{\chapter}{0pt}{-8pt}{14pt}
\titlecontents{section}[2.2em]{\addvspace{1pt}\sffamily\footnotesize}{}{}{\hfill\contentspage}
\titlecontents{chapter}[0pt]{\addvspace{5pt}\sffamily}
  {\thecontentslabel\quad}{}{\hfill\contentspage}
\newcommand{\lec}[2]{\needspace{8\baselineskip}\vspace{10pt}%
  {\sffamily\bfseries\fontsize{12.5}{16}\selectfont\color{ink}#1\quad#2}\par
  \vspace{-3pt}{\color{linet}\rule{\linewidth}{0.35pt}}\par\vspace{3pt}}
\newcommand{\onel}[1]{\vspace{1pt}{\kai\color{accent}\fontsize{10.5}{15.5}\selectfont #1}\par\vspace{3pt}}
\newcommand{\meta}[2]{\vspace{1pt}\begingroup\sloppy\raggedright{\sffamily\footnotesize\color{grayt}#1：#2}\par\endgroup}

\begin{document}

\thispagestyle{empty}
\begin{center}
\vspace*{3.2cm}
{\sffamily\fontsize{10}{14}\selectfont\color{grayt}精读笔记 \quad 153 讲速查}\par
\vspace{14pt}
{\sffamily\bfseries\fontsize{30}{40}\selectfont AI Infra}\par
\vspace{6pt}
{\sffamily\bfseries\fontsize{22}{30}\selectfont 精读笔记}\par
\vspace{18pt}
{\color{linet}\rule{0.5\textwidth}{0.6pt}}\par
\vspace{16pt}
{\kai\fontsize{12}{19}\selectfont
《AI Infra Handbook》全 20 篇 153 讲，1342 页\\[3pt]
每讲压成一张卡片：一句话结论、机制要点、关键数字、易错点\par}
\vfill
{\sffamily\footnotesize\color{grayt}
原书：孙晓辉《AI Infra Handbook》V1.0\quad 笔记整理：按篇序逐讲压缩\\[2pt]
本笔记只保留结论、公式、数字与坑，不复述入门推导\par}
\vspace{1.2cm}
\end{center}
\clearpage


\chapter*{怎么用这本笔记}
\addcontentsline{toc}{chapter}{怎么用这本笔记}
\markboth{怎么用这本笔记}{}

这本笔记是把原书 1342 页「拆掉讲稿、只留骨架」的结果。原书是逐讲展开的教学稿，每讲都从
「为什么要学」讲起，配可运行的代码和交互式演示；笔记假设你已经会这些，所以只保留四类东西：

\begin{itemize}[leftmargin=1.6em,itemsep=3pt,topsep=4pt]
\item \textbf{一句话结论}（蓝色斜体）：这一讲到底解决了什么问题、给出的做法是什么。
\item \textbf{机制要点}（3\textasciitilde 6 条）：数据怎么流、谁依赖谁、为什么这么设计，尽量带数字。
\item \textbf{关键数字}：显存、带宽、FLOPs、延迟、参数量这类可以拿来估算的量。
\item \textbf{易错点}：原文「容易混淆的地方」里最反直觉的那一句。
\end{itemize}

\vspace{4pt}
{\sffamily\bfseries 建议的用法}\par
\vspace{2pt}
\begin{itemize}[leftmargin=1.6em,itemsep=2pt,topsep=2pt]
\item 复习：只看每讲的蓝色结论和加粗数字，一页扫十几讲。
\item 动手前：翻到对应的「篇」，把机制要点过一遍，再回去看原书的代码。
\item 面试/评审：用「关键数字」那几行撑估算，用「易错点」那几行挡追问。
\end{itemize}

\vspace{6pt}
{\sffamily\bfseries 说明}\par
\vspace{2pt}
讲号与原书一致（第 1\textasciitilde 153 讲），篇的划分也照原书（20 篇）。原书标注「选读」的
几讲（三角函数、长上下文技术等）同样收录，因为对做推理的人反而常用。数字与公式均来自原书；
个别地方原文的排版把数字和正文拆到了不同行，整理时按语义拼了回去，若与原书有出入以原书为准。

\vspace{8pt}
{\footnotesize\color{grayt}原书 20 篇 153 讲、1342 页，本笔记覆盖全部 153 讲。}
\clearpage


\tableofcontents
\clearpage
\chapter{数学基础}
{\footnotesize\sffamily\color{grayt}收录第 1--11 讲}\par\vspace{-4pt}
\phantomsection\addcontentsline{toc}{section}{1\quad 向量}
\lec{1}{向量}
\onel{把事物变成一串有序的数，事物之间的关系就变成可计算的长度与距离。}
\begin{itemize}[leftmargin=1.35em,itemsep=2.2pt,topsep=1pt,parsep=0pt]
\item \begingroup\sloppy{\fontsize{10}{14.6}\selectfont n 个数按固定顺序排成一串就是 n 维向量，(7,4) 与 (4,7) 因顺序不同而是两种东西。\par}\endgroup
\item \begingroup\sloppy{\fontsize{10}{14.6}\selectfont 长度 |v|=$\surd$(v$_1$$^2$+…+vₙ$^2$)，苹果 (7,4) 是 $\surd$65$\approx$8.06；两点距离即可当作相似度。\par}\endgroup
\item \begingroup\sloppy{\fontsize{10}{14.6}\selectfont 基本运算只有加法（对应分量相加、维度须相同）与数乘（每分量乘 k，长度变 |k| 倍）。\par}\endgroup
\item \begingroup\sloppy{\fontsize{10}{14.6}\selectfont 反复用到的都是线性组合 k$_1$a+k$_2$b：第 26 讲注意力输出、第 29 讲残差连接即此。\par}\endgroup
\item \begingroup\sloppy{\fontsize{10}{14.6}\selectfont 维度再高规则不变；Llama-3-8B 每词 4096 维，BF16 下 4096$\times$2=8192 字节即 8 KiB。\par}\endgroup
\item \begingroup\sloppy{\fontsize{10}{14.6}\selectfont 1000 个 token 的向量约 7.8 MiB；词表 128256 个词的向量表约 0.98 GiB。\par}\endgroup
\end{itemize}
\meta{公式}{$\lVert v\rVert=\sqrt{v_1^2+v_2^2+\cdots+v_n^2}$；$k\cdot v=(kv_1,kv_2,\ldots,kv_n)$}
\meta{数字}{4096 维；BF16 每分量 2 字节；单向量 8192 字节 = 8 KiB；1000 token $\approx$7.8 MiB；词表 $\approx$0.98 GiB}
\meta{易错}{4096 维向量在 NumPy 里是形状 (4096,) 的一维数组：向量的「维」指有几个分量，与数组的 ndim 不是一回事。}
\phantomsection\addcontentsline{toc}{section}{2\quad 点积与余弦相似度}
\lec{2}{点积与余弦相似度}
\onel{点积等于两长度乘夹角余弦，除掉长度就是只看方向的余弦相似度。}
\begin{itemize}[leftmargin=1.35em,itemsep=2.2pt,topsep=1pt,parsep=0pt]
\item \begingroup\sloppy{\fontsize{10}{14.6}\selectfont 距离把长度差当内容差：(4,1) 到 (12,3) 距离 8.25，到 (1,5) 只有 5，判错。\par}\endgroup
\item \begingroup\sloppy{\fontsize{10}{14.6}\selectfont 点积是对应分量相乘再相加，结果是一个数：(4,1)$\cdot$(12,3)=51，(4,1)$\cdot$(1,5)=9。\par}\endgroup
\item \begingroup\sloppy{\fontsize{10}{14.6}\selectfont 由余弦定理推得 a$\cdot$b=|a||b|cos$\theta$，n 维同样成立；点积符号即 cos$\theta$ 符号，为 0 即垂直。\par}\endgroup
\item \begingroup\sloppy{\fontsize{10}{14.6}\selectfont 余弦相似度 = a$\cdot$b/(|a||b|)，落在 -1 到 1：短文与长文 =1，短文与美食文 $\approx$0.43。\par}\endgroup
\item \begingroup\sloppy{\fontsize{10}{14.6}\selectfont 点积便宜：d 维约 2d 次 FLOPs，4096 维约 8192 次，且能成批交给 GPU 矩阵乘硬件。\par}\endgroup
\item \begingroup\sloppy{\fontsize{10}{14.6}\selectfont 注意力里每词对每词做点积：s=1000、32 头、128 维、32 层约 2.6$\times$10$^1$$^1$ FLOPs，随 s$^2$ 涨。\par}\endgroup
\end{itemize}
\meta{公式}{$a\cdot b=\sum_{i}a_i b_i=\lVert a\rVert\,\lVert b\rVert\cos\theta$；$\cos\theta=\frac{a\cdot b}{\lVert a\rVert\,\lVert b\rVert}$}
\meta{数字}{距离 8.25 与 5；点积 51 与 9；余弦 1.0 与 0.43；4096 维点积约 8192 次；1000$^2$$\times$32$\times$2$\times$128$\times$32$\approx$2.6$\times$10$^1$$^1$ FLOPs}
\meta{易错}{余弦为 1 不等于向量相同（(4,1) 与 (12,3) 方向同但不相等）；为 0 是垂直，相反才是 -1。}
\phantomsection\addcontentsline{toc}{section}{3\quad 矩阵乘法}
\lec{3}{矩阵乘法}
\onel{矩阵乘法 C=XW 的每一格都是 X 的一行与 W 的一列的点积，整张 C 就是点积表。}
\begin{itemize}[leftmargin=1.35em,itemsep=2.2pt,topsep=1pt,parsep=0pt]
\item \begingroup\sloppy{\fontsize{10}{14.6}\selectfont 4 个词对 4 个方向共 4$\times$4=16 个点积，按词成行、方向成列排成 4$\times$4 结果表。\par}\endgroup
\item \begingroup\sloppy{\fontsize{10}{14.6}\selectfont 形状规则 [m,k]@[k,n]$\rightarrow$[m,n]：左边列数必须等于右边行数，中间的 k 相消。\par}\endgroup
\item \begingroup\sloppy{\fontsize{10}{14.6}\selectfont 换角度：W 是同一台变换机器，X 的每一行各自变换，行与行互不干扰。\par}\endgroup
\item \begingroup\sloppy{\fontsize{10}{14.6}\selectfont 大模型里是 [s,d]@[d,d\_out]$\rightarrow$[s,d\_out]：Q/K/V 投影、FFN 升降维、输出层同形。\par}\endgroup
\item \begingroup\sloppy{\fontsize{10}{14.6}\selectfont 计算量约 2mnk FLOPs：k=n=4096 时 s=1 为 3.4$\times$10⁷，s=1000 为 3.4$\times$10$^1$⁰。\par}\endgroup
\item \begingroup\sloppy{\fontsize{10}{14.6}\selectfont 算术强度 = FLOPs÷读入字节 = 2skn/(kn$\cdot$2) = s；s=1 时每读 1 字节只算 1 次。\par}\endgroup
\end{itemize}
\meta{公式}{$C_{ij}=\sum_{t}A_{it}B_{tj}$；一次矩阵乘法约 $2mnk$ FLOPs；$\text{算术强度}=\frac{2skn}{kn\cdot 2}=s$}
\meta{数字}{16 个点积；16$\times$4=64 次乘加=128 FLOPs；4096$^2$ 权重 BF16 占 32 MiB；s=1 时 3.4$\times$10⁷、s=1000 时 3.4$\times$10$^1$⁰ FLOPs}
\meta{易错}{XW$\neq$WX，[2,3]@[3,2] 得 2$\times$2 而反过来得 3$\times$3；NumPy 的 * 是逐元素乘，矩阵乘要用 @。}
\phantomsection\addcontentsline{toc}{section}{4\quad 线性变换}
\lec{4}{线性变换}
\onel{矩阵就是线性变换，W 的每一行给出一个坐标轴方向被送去的终点。}
\begin{itemize}[leftmargin=1.35em,itemsep=2.2pt,topsep=1pt,parsep=0pt]
\item \begingroup\sloppy{\fontsize{10}{14.6}\selectfont e1=(1,0) 变换后正是 W 的第 1 行，e2 是第 2 行；这两行就完全确定了整个空间。\par}\endgroup
\item \begingroup\sloppy{\fontsize{10}{14.6}\selectfont 任意 (x$_1$,x$_2$) 的像 = x$_1$$\cdot$(W 第 1 行) + x$_2$$\cdot$(W 第 2 行)，即按原分量加权相加。\par}\endgroup
\item \begingroup\sloppy{\fontsize{10}{14.6}\selectfont 线性变换让原点不动、直线仍直、平行与等距保持；拉伸、旋转、剪切、镜像、压扁都是特例。\par}\endgroup
\item \begingroup\sloppy{\fontsize{10}{14.6}\selectfont 行列式 det=ad-bc 是面积缩放倍数：默认 W 为 2，旋转与剪切为 1，压扁为 0。\par}\endgroup
\item \begingroup\sloppy{\fontsize{10}{14.6}\selectfont det=0 表示空间被压扁、信息丢失、不可逆；第 40 讲 MLA 就是有控制地压到低维。\par}\endgroup
\item \begingroup\sloppy{\fontsize{10}{14.6}\selectfont 两次变换可合成一次 (x$\cdot$A)$\cdot$B=x$\cdot$(A$\cdot$B)，但顺序不可交换；纯线性层叠仍是线性，故需激活函数。\par}\endgroup
\end{itemize}
\meta{公式}{$x\cdot W=x_1\cdot(W\text{ 第 }1\text{ 行})+x_2\cdot(W\text{ 第 }2\text{ 行})$；$(x\cdot A)\cdot B=x\cdot(A\cdot B)$；$\det W=ad-bc$}
\meta{数字}{det 为 2（默认 W）；旋转、剪切 det=1；压扁 det=0；镜像 det=-1；4096$\times$4096 矩阵 1680 万个数，BF16 占 32 MiB}
\meta{易错}{这里向量写成横行、变换写成 x$\cdot$W，所以去向在 W 的行；课本是列向量 W$\cdot$x，差一个转置，旋转矩阵看起来也「反了」。}
\phantomsection\addcontentsline{toc}{section}{5\quad 张量与形状}
\lec{5}{张量与形状}
\onel{张量是任意维度的数的阵列，读代码先看形状往往就知道运算在做什么。}
\begin{itemize}[leftmargin=1.35em,itemsep=2.2pt,topsep=1pt,parsep=0pt]
\item \begingroup\sloppy{\fontsize{10}{14.6}\selectfont 标量 []、向量 [4]、矩阵 [4,4] 都是张量特例；[b,s,d] 是 b 句 $\times$ s token $\times$ d 维。\par}\endgroup
\item \begingroup\sloppy{\fontsize{10}{14.6}\selectfont 矩阵乘法只作用最后两维：[…,m,k]@[…,k,n]$\rightarrow$[…,m,n]，前面的维度当批一一对应。\par}\endgroup
\item \begingroup\sloppy{\fontsize{10}{14.6}\selectfont reshape 只换看法不动数据，transpose 才交换维度；拆多头要两者配合。\par}\endgroup
\item \begingroup\sloppy{\fontsize{10}{14.6}\selectfont 逐元素运算形状不变；归约沿哪一维哪一维就消失，softmax 要指定沿哪一维做。\par}\endgroup
\item \begingroup\sloppy{\fontsize{10}{14.6}\selectfont 广播从最后一维对齐：两维相等或其一为 1，缺的前面补 1；[2,4,4]+[4] 即每 token 加同一偏置。\par}\endgroup
\item \begingroup\sloppy{\fontsize{10}{14.6}\selectfont 大小 = 形状各维乘积 $\times$ 每数字节数：[8,2048,4096] 是 6700 万个数、BF16 下 134 MB。\par}\endgroup
\end{itemize}
\meta{公式}{$[\ldots,m,k]\mathbin{@}[\ldots,k,n]\to[\ldots,m,n]$；$\text{字节数}=\prod_{i}\text{shape}_i\times\text{每数字节}$}
\meta{数字}{[8,2048,4096] 6700 万个数 134 MB；注意力打分 [8,32,2048,2048] 为 2.1 GB；TinyLM 由 [2,4] 到 [2,2,4,4]}
\meta{易错}{reshape 与 transpose 不可互换：把 [s,h,d\_h] 直接 reshape 成 [h,s,d\_h] 只是重新切分，数据会错位。}
\phantomsection\addcontentsline{toc}{section}{6\quad 指数、对数与softmax}
\lec{6}{指数、对数与softmax}
\onel{softmax 先取指数再除以总和，把任意一组打分变成和为 1 的概率。}
\begin{itemize}[leftmargin=1.35em,itemsep=2.2pt,topsep=1pt,parsep=0pt]
\item \begingroup\sloppy{\fontsize{10}{14.6}\selectfont eˣ 恒正、单调、每加 1 就乘约 2.718；分数差 1 概率差约 2.7 倍，差 3 差 20 倍。\par}\endgroup
\item \begingroup\sloppy{\fontsize{10}{14.6}\selectfont ln 把乘法变加法：0.4666、0.2、0.1、0.05 相乘得 0.000467，对数和为 -7.67。\par}\endgroup
\item \begingroup\sloppy{\fontsize{10}{14.6}\selectfont 1000 个 0.01 相乘真值 10$^-$$^2$⁰⁰⁰ 会下溢成 0，对数和是 -4605.2，故训练用对数概率。\par}\endgroup
\item \begingroup\sloppy{\fontsize{10}{14.6}\selectfont TinyLM 打分 1.64、0.83、-0.39… 取指数得 5.155、2.293…，总和 11.048。\par}\endgroup
\item \begingroup\sloppy{\fontsize{10}{14.6}\selectfont 除以温度 T：T=0.5 时最高项 0.749，T=2 时降到 0.302，T$\rightarrow$0 退化成只取最高分。\par}\endgroup
\item \begingroup\sloppy{\fontsize{10}{14.6}\selectfont 分数同加一数结果不变，故先减最大值再取指数防溢出：softmax([1000,999]) 正常。\par}\endgroup
\end{itemize}
\meta{公式}{$p_i=\frac{e^{z_i/T}}{\sum_{j}e^{z_j/T}}$，实现时先减去 $\max(z)$}
\meta{数字}{e\textasciicircum{}z 总和 11.048；概率 0.4666/0.2076/0.0613；T=0.5$\rightarrow$0.749，T=2$\rightarrow$0.302；熵 1.48，ln6=1.79；Llama-3 词表 128256}
\meta{易错}{softmax 只看差值：(1,2,3) 与 (101,102,103) 概率相同，但 (2,4,6) 不同，乘 2 等于温度减半。}
\phantomsection\addcontentsline{toc}{section}{7\quad 概率与采样}
\lec{7}{概率与采样}
\onel{累积概率把 [0,1) 切成长度等于概率的段，均匀随机数 u 落哪段就选哪个词。}
\begin{itemize}[leftmargin=1.35em,itemsep=2.2pt,topsep=1pt,parsep=0pt]
\item \begingroup\sloppy{\fontsize{10}{14.6}\selectfont 离散分布要求每个概率在 0 到 1 之间且总和为 1，模型给出的下一个词就是这种分布。\par}\endgroup
\item \begingroup\sloppy{\fontsize{10}{14.6}\selectfont 累积概率 0.4666、0.6742、0.7355、0.8119、0.869、1；u=0.7 落第三段选「学」。\par}\endgroup
\item \begingroup\sloppy{\fontsize{10}{14.6}\selectfont 大数定律：抽 10 次「我」0.3，100 次 0.42，10000 次 0.468，逼近概率 0.4666。\par}\endgroup
\item \begingroup\sloppy{\fontsize{10}{14.6}\selectfont 伪随机数由种子决定序列，固定种子结果可复现；换种子或换 GPU 并行顺序才会变。\par}\endgroup
\item \begingroup\sloppy{\fontsize{10}{14.6}\selectfont 贪心总取概率最大，确定但易重复；采样多样却可能抽到低概率词，折中是 Top-k、Top-p。\par}\endgroup
\item \begingroup\sloppy{\fontsize{10}{14.6}\selectfont 熵 H=-$\Sigma$p$\cdot$ln p 衡量分散度：此分布 1.479，6 词均匀时最大 ln6=1.792。\par}\endgroup
\end{itemize}
\meta{公式}{$\text{cdf}_i=\sum_{j\le i}p_j$，取 $u\sim U[0,1)$ 落段即取样；$H=-\sum_{i}p_i\ln p_i$}
\meta{数字}{累积 0.4666/0.6742/0.7355/0.8119/0.869/1；抽 10000 次 0.468；熵 1.479；ln6=1.792}
\meta{易错}{频率不等于概率：一次生成只是单次抽样，评测要同题多次生成再统计。}
\phantomsection\addcontentsline{toc}{section}{8\quad 导数与链式法则}
\lec{8}{导数与链式法则}
\onel{导数是局部放大倍数，多层嵌套靠链式法则把各段局部导数连乘。}
\begin{itemize}[leftmargin=1.35em,itemsep=2.2pt,topsep=1pt,parsep=0pt]
\item \begingroup\sloppy{\fontsize{10}{14.6}\selectfont 导数是 (f(x+h)-f(x))/h 在 h$\rightarrow$0 的极限；x$^2$ 在 3 处得 6.1、6.01、6.001$\rightarrow$6。\par}\endgroup
\item \begingroup\sloppy{\fontsize{10}{14.6}\selectfont 常用导数：x$^n$$\rightarrow$n$\cdot$x$^n$$^-$$^1$，eˣ$\rightarrow$eˣ，ln x$\rightarrow$1/x，sigmoid′ = σ(1-σ)，最大只有 0.25。\par}\endgroup
\item \begingroup\sloppy{\fontsize{10}{14.6}\selectfont 偏导只让一个变量动：L=(wx+b-y)$^2$ 在 w=b=1 处 $\partial$L/$\partial$w=-16、$\partial$L/$\partial$b=-8。\par}\endgroup
\item \begingroup\sloppy{\fontsize{10}{14.6}\selectfont 梯度 ($\partial$L/$\partial$w, $\partial$L/$\partial$b) 指向损失增长最快的方向，梯度下降沿其反方向走一小步。\par}\endgroup
\item \begingroup\sloppy{\fontsize{10}{14.6}\selectfont 链式法则：x$\rightarrow$u=3x+1$\rightarrow$v=u$^2$$\rightarrow$y=ln v 在 x=0.5 处，dy/dx=0.16$\times$5$\times$3=2.4。\par}\endgroup
\item \begingroup\sloppy{\fontsize{10}{14.6}\selectfont 多路径时路径上相乘、路径之间相加；连乘导致梯度消失或爆炸，激活函数、残差即为此。\par}\endgroup
\end{itemize}
\meta{公式}{$f'(x)=\lim_{h\to 0}\frac{f(x+h)-f(x)}{h}$；$\frac{dy}{dx}=\frac{dy}{dv}\cdot\frac{dv}{du}\cdot\frac{du}{dx}$}
\meta{数字}{割线 6.1/6.01/6.001$\rightarrow$导数 6；$\partial$L/$\partial$w=-16、$\partial$L/$\partial$b=-8；0.16$\times$5$\times$3=2.4；σ′(0)=0.25；Llama-3-8B 80 亿参数}
\meta{易错}{同一路径上相乘、多条路径之间相加，两者别混；导数是局部的，步子走远方向就不准。}
\phantomsection\addcontentsline{toc}{section}{9\quad 三角函数与旋转}
\lec{9}{三角函数与旋转}
\onel{旋转保长度、保夹角；各转 mφ、nφ 时点积只取决于 m-n，即 RoPE。}
\begin{itemize}[leftmargin=1.35em,itemsep=2.2pt,topsep=1pt,parsep=0pt]
\item \begingroup\sloppy{\fontsize{10}{14.6}\selectfont 单位圆上角度 $\theta$ 对应点 (cos $\theta$, sin $\theta$)，sin、cos 周期 2$\pi$。\par}\endgroup
\item \begingroup\sloppy{\fontsize{10}{14.6}\selectfont 旋转不改变长度；先转 20$^{\circ}$ 再转 40$^{\circ}$ 等于直接转 60$^{\circ}$（角度相加）。\par}\endgroup
\item \begingroup\sloppy{\fontsize{10}{14.6}\selectfont 两个向量同转 φ 点积不变；各转 mφ、nφ 时点积只取决于 m-n。\par}\endgroup
\item \begingroup\sloppy{\fontsize{10}{14.6}\selectfont 代码取 φ=10$^{\circ}$，相差 2 步的三种情况点积都是 0.0643。\par}\endgroup
\item \begingroup\sloppy{\fontsize{10}{14.6}\selectfont 各维两两分组，第 i 组每步转 $\theta$$_i$ = 10000\textasciicircum{}(-2i/d) 弧度。\par}\endgroup
\item \begingroup\sloppy{\fontsize{10}{14.6}\selectfont 快的组近处敏感但会绕回，慢的组能分远方，快慢搭配像秒针分针时针。\par}\endgroup
\end{itemize}
\meta{公式}{$x'=a\cos\varphi-b\sin\varphi$，$y'=a\sin\varphi+b\cos\varphi$；$\theta_i=10000^{-2i/d}$；转一圈步数 $=\frac{2\pi}{\theta_i}$}
\meta{数字}{特殊角 0$^{\circ}$/30$^{\circ}$/90$^{\circ}$/180$^{\circ}$；(1, 0.5) 转 30$^{\circ}$ $\rightarrow$ (0.616, 0.933)，长度 1.1180 不变；q$\cdot$k = -0.15 转 50$^{\circ}$ 后仍是 -0.15；d=4 第 0 组 $\theta$=1、6 步一圈；d=128 第 63 组 $\theta$$\approx$0.000115、54,410 步一圈}
\meta{易错}{点积不变只在转同样角度时成立；转不同角度点积会变，但只取决于角度差------这正是 RoPE 能编码相对位置的原因。}
\phantomsection\addcontentsline{toc}{section}{10\quad 浮点数与精度}
\lec{10}{浮点数与精度}
\onel{浮点数=符号+指数+尾数，尾数定精度、指数定范围；训练需另存 FP32 主权重。}
\begin{itemize}[leftmargin=1.35em,itemsep=2.2pt,topsep=1pt,parsep=0pt]
\item \begingroup\sloppy{\fontsize{10}{14.6}\selectfont 符号 1 位，指数存「指数+bias」，尾数存 1.xxxx 的小数部分。\par}\endgroup
\item \begingroup\sloppy{\fontsize{10}{14.6}\selectfont 尾数位定精度，机器精度 2\textasciicircum{}-m：FP32 1.19e-7、FP16 9.77e-4。\par}\endgroup
\item \begingroup\sloppy{\fontsize{10}{14.6}\selectfont 指数位定范围：FP16 5 位最大 65504，BF16 与 FP32 同 8 位同范围。\par}\endgroup
\item \begingroup\sloppy{\fontsize{10}{14.6}\selectfont 舍入累积、上溢成 $\infty$（70000 在 FP16）、下溢成 0（1e-8 在 FP16）。\par}\endgroup
\item \begingroup\sloppy{\fontsize{10}{14.6}\selectfont 0.01 加 10000 次：FP16 停在 32.0，FP32 得 100.003。\par}\endgroup
\item \begingroup\sloppy{\fontsize{10}{14.6}\selectfont 混合精度：BF16 算前后向，更新加到 FP32 主权重上，否则 0.0001 被舍掉。\par}\endgroup
\end{itemize}
\meta{公式}{$\text{值}=(-1)^{s}\times\left(1+\frac{\text{尾数}}{2^{m}}\right)\times 2^{e-\text{bias}}$；尾数 $m$ 位时 $1$ 附近的间隔是 $2^{-m}$，相对误差约为它的一半}
\meta{数字}{FP32 8/23 位、FP16 5/10 位、BF16 8/7 位、FP8 E4M3 4/3 位、FP4 E2M1 2/1 位；最大值 3.4e38 / 65504 / 3.38e38 / 448 / 6；0.1 在 FP16 存 0.0999755859、BF16 存 0.1000976562}
\meta{易错}{FP8、FP4 范围太小，必须先乘缩放因子；BF16 不是「更好的 FP16」，它用精度换范围，训练多用它是因为不溢出。}
\phantomsection\addcontentsline{toc}{section}{11\quad 计算量与 FLOPs}
\lec{11}{计算量与 FLOPs}
\onel{计算量 ÷ 算力 = 理想时间；推理每 token 2N，训练 6N。}
\begin{itemize}[leftmargin=1.35em,itemsep=2.2pt,topsep=1pt,parsep=0pt]
\item \begingroup\sloppy{\fontsize{10}{14.6}\selectfont FLOPs 是运算次数，FLOPS 是每秒运算次数，两者相除得理想时间。\par}\endgroup
\item \begingroup\sloppy{\fontsize{10}{14.6}\selectfont 矩阵乘 [m,k]@[k,n] $\approx$ 2mnk；4$\times$4 逐次 112、估算 128。\par}\endgroup
\item \begingroup\sloppy{\fontsize{10}{14.6}\selectfont 一个 token 过一个权重矩阵 $\approx$ 该矩阵参数量的 2 倍，整个前向 $\approx$ 2N。\par}\endgroup
\item \begingroup\sloppy{\fontsize{10}{14.6}\selectfont 8B 每 token 约 15 GFLOP；上下文 2048 时注意力再占 7\%。\par}\endgroup
\item \begingroup\sloppy{\fontsize{10}{14.6}\selectfont 训练每 token $\approx$ 6N（前向 2N + 反向 4N），即前向的 3 倍。\par}\endgroup
\item \begingroup\sloppy{\fontsize{10}{14.6}\selectfont 生成 1 token 满算力只需 15.2 $\mu$s，实际毫秒级，瓶颈是搬 16 GB 权重。\par}\endgroup
\end{itemize}
\meta{公式}{$\text{时间}\approx\frac{\text{FLOPs}}{\text{FLOPS}}$；$\text{MatMul}\approx 2mnk$；推理 $2N$/token，训练 $6N$/token，训练总量 $\approx 6NT$}
\meta{数字}{H100 BF16 989 TFLOPS；Llama-3-8B 80.3 亿参数、去掉 5.3 亿嵌入表后 75.0 亿；prefill 1000 token 15.2 ms；上下文 2048 时注意力每 token +1.07 GFLOP；1000 张 H100 按 40\% 利用率约 21 天；MoE DeepSeek-V3 总参数 6710 亿、每 token 激活约 370 亿}
\meta{易错}{FLOPs（计算量）和 FLOPS（算力）差一个「每秒」；有的规格按乘加 MAC 计数，数字是本讲的一半，比较时先看口径。}
\chapter{神经网络基础}
{\footnotesize\sffamily\color{grayt}收录第 12--20 讲}\par\vspace{-4pt}
\phantomsection\addcontentsline{toc}{section}{12\quad 感知机}
\lec{12}{感知机}
\onel{感知机=点积+偏置+阶跃函数，用一条直线分割平面，学不会线性不可分的异或。}
\begin{itemize}[leftmargin=1.35em,itemsep=2.2pt,topsep=1pt,parsep=0pt]
\item \begingroup\sloppy{\fontsize{10}{14.6}\selectfont z = w$\cdot$x + b 大过 0 就输出 1；w 垂直于分界线、指向输出 1 的一侧。\par}\endgroup
\item \begingroup\sloppy{\fontsize{10}{14.6}\selectfont 学习规则：错了 w$\leftarrow$w+(y-输出)x、b$\leftarrow$b+(y-输出)，对了什么都不改。\par}\endgroup
\item \begingroup\sloppy{\fontsize{10}{14.6}\selectfont 改一次让该样本的 z 至少涨 |x|$^2$，线性可分时有限步收敛（1962 年证明）。\par}\endgroup
\item \begingroup\sloppy{\fontsize{10}{14.6}\selectfont 分界线不唯一：w、b 同乘 2 不变；只保证分对，不保证找到最好的线。\par}\endgroup
\item \begingroup\sloppy{\fontsize{10}{14.6}\selectfont 异或的矛盾：要有 b$\leq$0 且 w1+w2+b>0，又要有 w1+w2+b$\leq$0。\par}\endgroup
\item \begingroup\sloppy{\fontsize{10}{14.6}\selectfont 两层可解：h1 算或、h2 算与非、第三层算与，等于先换一次坐标再分割。\par}\endgroup
\end{itemize}
\meta{公式}{$z=w\cdot x+b$，$y=1\ (z>0)$ 否则 $0$；错时 $w\leftarrow w+(y-\hat{y})x$，$b\leftarrow b+(y-\hat{y})$；分界线 $w\cdot x+b=0$，$w$ 垂直指向输出 $1$ 一侧}
\meta{数字}{水果 7 点：w=(1, -1)、b=-1 全对；代码 2 轮学到 w=(5, -5)、b=0；与第 6 轮 w=(2, 1)、b=-2；或第 4 轮；异或跑满 20 轮仍有错}
\meta{易错}{异或学不会是模型只能画一条直线的限制，换起点、跑更多轮都没用；分界线被压到时 z=0 就输出 0。}
\phantomsection\addcontentsline{toc}{section}{13\quad 神经元与线性层}
\lec{13}{神经元与线性层}
\onel{神经元=点积+偏置+激活；n 个神经元并排打分就是 y=x$\cdot$W+b。}
\begin{itemize}[leftmargin=1.35em,itemsep=2.2pt,topsep=1pt,parsep=0pt]
\item \begingroup\sloppy{\fontsize{10}{14.6}\selectfont 阶跃函数下参数小改动输出不变，拿掉它输出才对参数有反应。\par}\endgroup
\item \begingroup\sloppy{\fontsize{10}{14.6}\selectfont 权重 wi 动 Δ，输出就动 Δ$\times$xi；偏置相当于连着恒为 1 的输入。\par}\endgroup
\item \begingroup\sloppy{\fontsize{10}{14.6}\selectfont W 是 din$\times$dout、b 有 dout 个，参数 din$\times$dout+dout。\par}\endgroup
\item \begingroup\sloppy{\fontsize{10}{14.6}\selectfont 样本叠成矩阵 X，一次 X$\cdot$W 算完整批，b 由广播加到每一行。\par}\endgroup
\item \begingroup\sloppy{\fontsize{10}{14.6}\selectfont 几何上是先线性变换再平移（仿射），把输入送到新的打分空间。\par}\endgroup
\item \begingroup\sloppy{\fontsize{10}{14.6}\selectfont Llama-3-8B 每层 7 个线性层都不带偏置，全模型近九成参数在线性层里。\par}\endgroup
\end{itemize}
\meta{公式}{$y=x\cdot W+b$；参数个数 $=d_{\text{in}}d_{\text{out}}+d_{\text{out}}$；$(x,1)\cdot[W;b]=x\cdot W+b$；PyTorch：$y=x\cdot W^{\mathsf{T}}+b$，$W$ 形状 $[d_{\text{out}},d_{\text{in}}]$}
\meta{数字}{水果 2$\times$3+3=9 个参数；TinyLM 4$\times$8+8=40；up\_proj 4096$\rightarrow$14336 为 58,720,256；每层线性层合计 218,103,808，32 层 6,979,321,856，占 80.3 亿的 86.9\%，含输出层超 93\%}
\meta{易错}{PyTorch 的 nn.Linear 把权重存成 [dout, din]（每行一个神经元），计算时用 y = x$\cdot$Wᵀ + b，读权重文件先看清存法。}
\phantomsection\addcontentsline{toc}{section}{14\quad 激活函数}
\lec{14}{激活函数}
\onel{线性层叠多少层都等于一层；插入非线性激活才打破塌缩，两个 ReLU 能算异或。}
\begin{itemize}[leftmargin=1.35em,itemsep=2.2pt,topsep=1pt,parsep=0pt]
\item \begingroup\sloppy{\fontsize{10}{14.6}\selectfont (x$\cdot$W1+b1)$\cdot$W2+b2 展开仍是 x 乘一个矩阵加一个向量，叠十层等于一层。\par}\endgroup
\item \begingroup\sloppy{\fontsize{10}{14.6}\selectfont ReLU(z) = max(0, z) 逐元素作用、不改形状，负数一律截成 0。\par}\endgroup
\item \begingroup\sloppy{\fontsize{10}{14.6}\selectfont 两个 ReLU 算异或：h2 只在 (1,1) 亮，y = h1 - 2h2 拉回 0。\par}\endgroup
\item \begingroup\sloppy{\fontsize{10}{14.6}\selectfont 一层 N 个 ReLU 只有 N 个折点，折点越多折线越贴近曲线（万能近似）。\par}\endgroup
\item \begingroup\sloppy{\fontsize{10}{14.6}\selectfont GELU=z$\cdot$Φ(z)、SiLU=z$\cdot$σ(z) 都是 z 乘 0\textasciitilde{}1 软开关。\par}\endgroup
\item \begingroup\sloppy{\fontsize{10}{14.6}\selectfont ReLU 负侧斜率 0 导致死亡 ReLU；Sigmoid 斜率上限 0.25、梯度消失。\par}\endgroup
\end{itemize}
\meta{公式}{$\operatorname{ReLU}(z)=\max(0,z)$；$\operatorname{GELU}(z)=z\,\Phi(z)$；$\operatorname{SiLU}(z)=z\,\sigma(z)$，$\sigma(z)=\frac{1}{1+e^{-z}}$；$h_1=\operatorname{ReLU}(x_1+x_2)$、$h_2=\operatorname{ReLU}(x_1+x_2-1)$、$y=h_1-2h_2$}
\meta{数字}{ReLU 用 N=6 拟合 sin(1.5x) 误差 0.0067，不加激活是 0.48；z=-1 处 GELU -0.159、SiLU -0.269；最低点 GELU z$\approx$-0.75、SiLU z$\approx$-1.28；Llama-3-8B 每层 3$\times$4096$\times$14336$\approx$1.76 亿次乘加}
\meta{易错}{激活计算量极小但是访存密集型，推理时把它和前后算子融合（第 83 讲）以免多一次读写。}
\phantomsection\addcontentsline{toc}{section}{15\quad 多层感知机}
\lec{15}{多层感知机}
\onel{MLP 把线性层与激活层层串起来：深度让折线段数指数增长，宽度只能线性增长。}
\begin{itemize}[leftmargin=1.35em,itemsep=2.2pt,topsep=1pt,parsep=0pt]
\item \begingroup\sloppy{\fontsize{10}{14.6}\selectfont 层间 h=f(h$\cdot$W+b)，最后接一个不加激活的线性层输出。\par}\endgroup
\item \begingroup\sloppy{\fontsize{10}{14.6}\selectfont 1$\rightarrow$8$\rightarrow$8$\rightarrow$1 共 97 个参数（16+72+9），宽 8、深 2。\par}\endgroup
\item \begingroup\sloppy{\fontsize{10}{14.6}\selectfont 第 1 层每个神经元把 x 折一次，折点在 x = -bj/wj 处。\par}\endgroup
\item \begingroup\sloppy{\fontsize{10}{14.6}\selectfont 第 2 层加权后穿过 0 处再折一次，产生新拐角：8 个折点变成 17 段。\par}\endgroup
\item \begingroup\sloppy{\fontsize{10}{14.6}\selectfont 锯齿网络每层 2 个神经元、k 层折 2\textasciicircum{}k 段；单层要 64 段需 63 个神经元。\par}\endgroup
\item \begingroup\sloppy{\fontsize{10}{14.6}\selectfont 大模型 FFN 是 4096$\rightarrow$14336$\rightarrow$4096 的单隐层 MLP，占每层参数 81\%。\par}\endgroup
\end{itemize}
\meta{公式}{$h_1=f(x\cdot W_1+b_1)$；$h_2=f(h_1\cdot W_2+b_2)$；$y=h_2\cdot W_3+b_3$；$t(u)=2\operatorname{ReLU}(u)-4\operatorname{ReLU}(u-0.5)$，$k$ 层折出 $2^k$ 段}
\meta{数字}{1$\rightarrow$8$\rightarrow$8$\rightarrow$1 参数 97、输出 17 段、MSE 0.00017；x=1 时 y=0.998，sin(1.5)=0.997；第 1 层 8 个折点在 -0.74\textasciitilde{}1.41；第 2 层只有 6 个神经元有输出；FFN 三个矩阵 176,160,768 个参数}
\meta{易错}{「几层」的数法不统一：1$\rightarrow$8$\rightarrow$8$\rightarrow$1 可数成 2 层（隐藏层）、3 层（权重矩阵）或 4 层；大模型的「32 层」数的是 Transformer 块。}
\phantomsection\addcontentsline{toc}{section}{16\quad 损失函数}
\lec{16}{损失函数}
\onel{拟合数值用 MSE，预测概率用交叉熵 -ln p，它的指数就是困惑度。}
\begin{itemize}[leftmargin=1.35em,itemsep=2.2pt,topsep=1pt,parsep=0pt]
\item \begingroup\sloppy{\fontsize{10}{14.6}\selectfont MSE 先平方再平均：误差不会正负抵消、大错罚得重、处处光滑。\par}\endgroup
\item \begingroup\sloppy{\fontsize{10}{14.6}\selectfont logit 过 Sigmoid 得 p，损失 -ln p；p=0.01 时 4.605。\par}\endgroup
\item \begingroup\sloppy{\fontsize{10}{14.6}\selectfont 交叉熵配 Sigmoid，对 logit 斜率是 p-1；平方误差多乘 p(1-p)。\par}\endgroup
\item \begingroup\sloppy{\fontsize{10}{14.6}\selectfont 多分类用 softmax 把 V 个分变概率，损失仍是 -ln(正确词概率)。\par}\endgroup
\item \begingroup\sloppy{\fontsize{10}{14.6}\selectfont 一句话的损失是各位置交叉熵的平均；困惑度是 e\textasciicircum{}平均交叉熵。\par}\endgroup
\item \begingroup\sloppy{\fontsize{10}{14.6}\selectfont 均匀乱猜损失就是 ln V；Llama-3 词表 128256 $\rightarrow$ 11.76。\par}\endgroup
\end{itemize}
\meta{公式}{$\text{MSE}=\frac{1}{n}\sum_{i}(\hat{y}_i-y_i)^2$；$\text{CE}=-\ln p$；$\text{PPL}=e^{\text{平均交叉熵}}$；softmax：$p_i=\frac{e^{z_i}}{\sum_{j}e^{z_j}}$；修正信号 $\frac{\partial L}{\partial z_i}=p_i-y_i$}
\meta{数字}{MSE 例 0.125；p=0.99 时 CE 0.010，p=0.01 时 4.605；p=0.01 时信号 -0.99、平方误差只有 -0.0196；TinyLM 四个位置 1.572 / 2.323 / 1.601 / 2.421 $\rightarrow$ 平均 1.979、PPL 7.24；均匀猜 ln 6 = 1.792、PPL 6；ln 128256 $\approx$ 11.76；8192$\times$128256 的 logits 要 3.91 GiB}
\meta{易错}{困惑度要先平均交叉熵再取指数；把每个位置的困惑度先算出来再平均会得到另一个更大的数。}
\phantomsection\addcontentsline{toc}{section}{17\quad 梯度下降}
\lec{17}{梯度下降}
\onel{对每个参数求偏导排成梯度，沿反方向走 η 倍，就能把损失降到最低。}
\begin{itemize}[leftmargin=1.35em,itemsep=2.2pt,topsep=1pt,parsep=0pt]
\item \begingroup\sloppy{\fontsize{10}{14.6}\selectfont 损失从预测的函数变成参数的函数，换一组参数就得到一个损失，训练找最小的那组。\par}\endgroup
\item \begingroup\sloppy{\fontsize{10}{14.6}\selectfont 斜率是负就往右走，步长 = 学习率 $\times$ 斜率；越近谷底斜率越小，步子自动变小。\par}\endgroup
\item \begingroup\sloppy{\fontsize{10}{14.6}\selectfont 距离每步乘 (1-7η)：η<1/7 单侧降，1/7\textasciitilde{}2/7 震荡收敛，>2/7 发散。\par}\endgroup
\item \begingroup\sloppy{\fontsize{10}{14.6}\selectfont 两个参数的偏导数排成梯度 $\nabla$L，它指向损失上升最快方向且垂直等高线。\par}\endgroup
\item \begingroup\sloppy{\fontsize{10}{14.6}\selectfont 扁碗里先冲向谷底线再沿谷慢慢挪：η>0.12 震荡、>0.24 发散，比 2/7 更小。\par}\endgroup
\item \begingroup\sloppy{\fontsize{10}{14.6}\selectfont 训练要额外存一份同样大的梯度：8B 模型 BF16 权重与梯度各 16 GB。\par}\endgroup
\end{itemize}
\meta{公式}{$\theta\leftarrow\theta-\eta\nabla L$；单参数收敛条件 $\lvert 1-7\eta\rvert<1$，即 $\eta<\frac{2}{7}\approx 0.286$}
\meta{数字}{w=1.6、b=1.1、L=0.05；斜率 -11.2$\rightarrow$-7.28$\rightarrow$-4.732$\rightarrow$-3.076；损失 9.01$\rightarrow$3.84$\rightarrow$1.65$\rightarrow$0.73；η=0.3 时 20 步到 -9.164；上限 2/7$\approx$0.286}
\meta{易错}{梯度指向上坡，更新必须减去它，代码里写成加号会越训越差；求导对象是参数而不是输入。}
\phantomsection\addcontentsline{toc}{section}{18\quad 反向传播}
\lec{18}{反向传播}
\onel{一次前向加一次反向，用上游梯度乘局部导数，就能算出所有参数的梯度。}
\begin{itemize}[leftmargin=1.35em,itemsep=2.2pt,topsep=1pt,parsep=0pt]
\item \begingroup\sloppy{\fontsize{10}{14.6}\selectfont 计算图把网络拆成乘、加、ReLU、平方等小运算，前向从左往右算并留下中间值。\par}\endgroup
\item \begingroup\sloppy{\fontsize{10}{14.6}\selectfont 反向从 L 出发，每个节点做一次「上游梯度 $\times$ 局部导数」，把梯度传给自己的输入。\par}\endgroup
\item \begingroup\sloppy{\fontsize{10}{14.6}\selectfont 一个节点被用几次，梯度就加几次；PyTorch 的 backward() 是累加而非覆盖。\par}\endgroup
\item \begingroup\sloppy{\fontsize{10}{14.6}\selectfont 线性层 Y=XW 的反向是两次矩阵乘法，梯度形状永远和对应变量相同。\par}\endgroup
\item \begingroup\sloppy{\fontsize{10}{14.6}\selectfont 前向的中间结果必须留到反向用完：乘法要用另一个输入，ReLU 要知道正负。\par}\endgroup
\item \begingroup\sloppy{\fontsize{10}{14.6}\selectfont 每 token 约 6N 次浮点运算：前向 2N + 反向 4N，是推理的 3 倍。\par}\endgroup
\end{itemize}
\meta{公式}{$\frac{\partial L}{\partial W}=X^{\mathsf{T}}dY$；$\frac{\partial L}{\partial X}=dY\,W^{\mathsf{T}}$；训练 $\text{FLOPs}\approx 6N=2N+4N$}
\meta{数字}{小网络 L=0.04、梯度 (-0.9, -0.6, -0.4, -0.4)；Llama-3-8B 每 token 前向 1.61e10、反向 3.21e10、合计 4.82e10 FLOPs；FFN 14336$\times$2$\times$8192$\times$32$\approx$7.0 GiB}
\meta{易错}{反向传播只负责算梯度，更新参数是优化器的事；反向也不是把网络倒过来运行。}
\phantomsection\addcontentsline{toc}{section}{19\quad 优化器}
\lec{19}{优化器}
\onel{Momentum 累积速度压震荡，Adam 用梯度平方为每个参数定步长。}
\begin{itemize}[leftmargin=1.35em,itemsep=2.2pt,topsep=1pt,parsep=0pt]
\item \begingroup\sloppy{\fontsize{10}{14.6}\selectfont x 平移 2 后最陡与最平方向弯曲差 166 倍，学习率上限 0.238$\rightarrow$0.069。\par}\endgroup
\item \begingroup\sloppy{\fontsize{10}{14.6}\selectfont SGD 每步用一小批数据估梯度，噪声让它在窄谷里来回震荡、平方向进展慢。\par}\endgroup
\item \begingroup\sloppy{\fontsize{10}{14.6}\selectfont Momentum 把过去的梯度累成速度，$\beta$=0.9 使一致方向的速度放大到 10g。\par}\endgroup
\item \begingroup\sloppy{\fontsize{10}{14.6}\selectfont Adam 维护 m、v 两个滑动平均，用 m̂/$\surd$v̂ 定标，每个参数都走约 η 远。\par}\endgroup
\item \begingroup\sloppy{\fontsize{10}{14.6}\selectfont 偏差修正除以 1-0.9\textasciicircum{}t 和 1-0.999\textasciicircum{}t，ε$\approx$1e-8 防除零。\par}\endgroup
\item \begingroup\sloppy{\fontsize{10}{14.6}\selectfont 训练每参数 16 字节（权重/梯度/主权重/m/v），8B 要 128.5 GB。\par}\endgroup
\end{itemize}
\meta{公式}{Momentum：$v\leftarrow\beta v+g$，$\theta\leftarrow\theta-\eta v$；Adam：$\theta\leftarrow\theta-\eta\frac{\hat{m}}{\sqrt{\hat{v}}+\varepsilon}$，$\hat{m}=\frac{m}{1-0.9^{t}}$、$\hat{v}=\frac{v}{1-0.999^{t}}$}
\meta{数字}{弯曲 14$\rightarrow$166 倍；η 上限 0.238$\rightarrow$0.069；SGD 0.03 走 300 步 L=0.1267；Momentum 39 步、Adam 93 步；梯度 (-28, -6.5)$\rightarrow$SGD (0.84, 0.195)、Adam (0.3, 0.3)；128.5 GB vs 推理 16.1 GB}
\meta{易错}{不同优化器的学习率不能直接比：SGD 步长是 η$\times$梯度，Adam 步长约等于 η。}
\phantomsection\addcontentsline{toc}{section}{20\quad 训练循环}
\lec{20}{训练循环}
\onel{训练循环五步：取批、前向、损失、反向、更新；batch 与学习率调度决定收敛。}
\begin{itemize}[leftmargin=1.35em,itemsep=2.2pt,topsep=1pt,parsep=0pt]
\item \begingroup\sloppy{\fontsize{10}{14.6}\selectfont 内层五步：取批、前向、损失、反向、更新；更新前先清零，PyTorch 梯度是累加的。\par}\endgroup
\item \begingroup\sloppy{\fontsize{10}{14.6}\selectfont batch 大则每 epoch 更新少，B=1 噪声大停不下来，B=8 兼顾两者。\par}\endgroup
\item \begingroup\sloppy{\fontsize{10}{14.6}\selectfont 噪声决定损失地板：MSE 0.0101；B=8 降到 0.0067 已在记忆噪声。\par}\endgroup
\item \begingroup\sloppy{\fontsize{10}{14.6}\selectfont 预热 20 epoch 升到 0.01，再 400 epoch 余弦降到 0。\par}\endgroup
\item \begingroup\sloppy{\fontsize{10}{14.6}\selectfont 大模型按 token 计 batch，一步几百万 token，靠梯度累积凑足。\par}\endgroup
\item \begingroup\sloppy{\fontsize{10}{14.6}\selectfont 总计算量 $\approx$6ND（N 参数、D token），8B$\times$15T 约 7.23e23。\par}\endgroup
\end{itemize}
\meta{公式}{$\text{lr}=\text{base}\cdot\min\left(1,\frac{\text{epoch}+1}{\text{warmup}}\right)\cdot 0.5\cdot\left(1+\cos\frac{\pi\cdot\text{epoch}}{\text{total}}\right)$；总计算 $\approx 6ND$}
\meta{数字}{64 样本、97 参数、400 epoch、3200 步；初始损失 4.8043；噪声地板 0.0101；B=64 0.0102、B=8 0.0067、B=1 0.0288；989 TFLOPS、40\% $\rightarrow$ 5.07e5 GPU 小时}
\meta{易错}{损失低于噪声地板说明在记忆噪声，要看验证集；横轴是 step 还是 epoch 会得出相反结论。}
\chapter{从文本到向量}
{\footnotesize\sffamily\color{grayt}收录第 21--24 讲}\par\vspace{-4pt}
\phantomsection\addcontentsline{toc}{section}{21\quad 分词与BPE}
\lec{21}{分词与BPE}
\onel{BPE 反复数相邻 token 对的频次、合并最高的一对，从语料学出子词词表。}
\begin{itemize}[leftmargin=1.35em,itemsep=2.2pt,topsep=1pt,parsep=0pt]
\item \begingroup\sloppy{\fontsize{10}{14.6}\selectfont 按字切词表小、任何文本都能切，但序列长、算力与显存都随 token 数涨。\par}\endgroup
\item \begingroup\sloppy{\fontsize{10}{14.6}\selectfont 按词切序列短但词表爆炸，未登录词没法处理，中文还没有空格分词。\par}\endgroup
\item \begingroup\sloppy{\fontsize{10}{14.6}\selectfont BPE 每轮两步：数所有相邻对频次 $\rightarrow$ 合并最高频的一对，重复到无高频对。\par}\endgroup
\item \begingroup\sloppy{\fontsize{10}{14.6}\selectfont 合并规则按顺序记录，切新句子先拆单字再依次应用规则，保证训练推理一致。\par}\endgroup
\item \begingroup\sloppy{\fontsize{10}{14.6}\selectfont 真实分词器从 256 个 UTF-8 字节起步，任何文字最坏退回字节，永不 OOV。\par}\endgroup
\item \begingroup\sloppy{\fontsize{10}{14.6}\selectfont 词表 128256 时，嵌入表+输出层约 10.5 亿参数，占整个模型 13\%。\par}\endgroup
\end{itemize}
\meta{公式}{词表 $=256+$ 合并次数；嵌入表与输出层共 $2Vd$ 个参数}
\meta{数字}{27 token $\rightarrow$ 16 token；11 单字 + 4 合并 = 15；合并 4/3/2/2 次；词表 128256、隐藏 4096、嵌入+输出 1.05B}
\meta{易错}{token 不等于字或词，同一段话的 token 数取决于分词器；编号大小没有含义，只是查表下标。}
\phantomsection\addcontentsline{toc}{section}{22\quad 词嵌入}
\lec{22}{词嵌入}
\onel{编号查 V$\times$d 的嵌入表得到向量，等价于独热乘 E，即无偏置的线性层。}
\begin{itemize}[leftmargin=1.35em,itemsep=2.2pt,topsep=1pt,parsep=0pt]
\item \begingroup\sloppy{\fontsize{10}{14.6}\selectfont 编号自带大小和距离，直接当输入会把虚假关系当真；独热又太长且彼此等距。\par}\endgroup
\item \begingroup\sloppy{\fontsize{10}{14.6}\selectfont 嵌入表是 V 行 d 列，第 i 行即编号 i 的向量，一句话查出 s$\times$d 矩阵。\par}\endgroup
\item \begingroup\sloppy{\fontsize{10}{14.6}\selectfont 查表 = onehot(id)$\cdot$E，所以嵌入层是没有偏置的线性层，表里每个数都可训练。\par}\endgroup
\item \begingroup\sloppy{\fontsize{10}{14.6}\selectfont 反向时梯度只加回用到的行，一个 batch 里没出现的 token 这一步行不变。\par}\endgroup
\item \begingroup\sloppy{\fontsize{10}{14.6}\selectfont 训练后苹果最像香蕉 0.99、鱼 0.97；训练前是学 0.67 等乱序。\par}\endgroup
\item \begingroup\sloppy{\fontsize{10}{14.6}\selectfont Llama-3-8B 嵌入表约 5.25 亿参数，BF16 下 0.98 GiB。\par}\endgroup
\end{itemize}
\meta{公式}{$x=E[\text{id}]$；$\text{onehot}(\text{id})\cdot E= E$ 的第 $\text{id}$ 行}
\meta{数字}{TinyLM 6$\times$4；小语料 14 词$\times$8 维；苹果$\rightarrow$香蕉 0.99、鱼 0.97、AI 0.93；128256$\times$4096=525M 参数、0.98 GiB、8 KiB/token}
\meta{易错}{嵌入里没有位置信息，同一词在任何上下文都是同一向量；含义分散在很多维，单维没意义。}
\phantomsection\addcontentsline{toc}{section}{23\quad 预测下一个token}
\lec{23}{预测下一个token}
\onel{预训练只做条件概率 P(下一个 token | 前文)，每个位置一道交叉熵题。}
\begin{itemize}[leftmargin=1.35em,itemsep=2.2pt,topsep=1pt,parsep=0pt]
\item \begingroup\sloppy{\fontsize{10}{14.6}\selectfont 输出是词表上的概率分布（logit 经 softmax），不是某一个词。\par}\endgroup
\item \begingroup\sloppy{\fontsize{10}{14.6}\selectfont 一段文本有多少 token 就有多少道题，答案是文本自身，属自监督学习。\par}\endgroup
\item \begingroup\sloppy{\fontsize{10}{14.6}\selectfont 损失是各位置交叉熵的平均 -ln p，小语料下限 0.6681、困惑度 1.951。\par}\endgroup
\item \begingroup\sloppy{\fontsize{10}{14.6}\selectfont 小模型概率几乎等于语料频率：我吃$\rightarrow$苹果 0.501/香蕉 0.499，小猫吃$\rightarrow$鱼 1.0。\par}\endgroup
\item \begingroup\sloppy{\fontsize{10}{14.6}\selectfont 训练时一次前向同时算所有位置（靠因果掩码），推理只要最后一个位置的分布。\par}\endgroup
\item \begingroup\sloppy{\fontsize{10}{14.6}\selectfont 语料够大时，为了压低损失被迫学出语法、事实、算术与推理，规模是关键。\par}\endgroup
\end{itemize}
\meta{公式}{$\text{loss}=-\frac{1}{T}\sum_{t}\ln p(y_t)$；$\text{perplexity}=\exp(\text{loss})$}
\meta{数字}{96 道题；交叉熵 0.6681、下限 0.6681、困惑度 1.951；我吃$\rightarrow$苹果 0.501/香蕉 0.499；15T token、约 7e23 FLOPs}
\meta{易错}{输出是分布不是词，选词规则独立于模型；基座模型只会续写，不会回答问题。}
\phantomsection\addcontentsline{toc}{section}{24\quad 自回归生成}
\lec{24}{自回归生成}
\onel{每步前向得分布、选一个 token 接回输入，循环到结束标记。}
\begin{itemize}[leftmargin=1.35em,itemsep=2.2pt,topsep=1pt,parsep=0pt]
\item \begingroup\sloppy{\fontsize{10}{14.6}\selectfont 循环三步：前向得分布、选词、接到末尾，直到结束标记或最大长度。\par}\endgroup
\item \begingroup\sloppy{\fontsize{10}{14.6}\selectfont 贪心每步取最大概率，确定但单调：苹果 0.501 永远压过香蕉 0.499。\par}\endgroup
\item \begingroup\sloppy{\fontsize{10}{14.6}\selectfont 采样按概率抽，T<1 分布变尖趋近贪心，T>1 变平会出现病句。\par}\endgroup
\item \begingroup\sloppy{\fontsize{10}{14.6}\selectfont 第 k 个 token 依赖第 k-1 个，生成 n 个要做 n 次前向，无法并行。\par}\endgroup
\item \begingroup\sloppy{\fontsize{10}{14.6}\selectfont 无 KV Cache 要处理约 150 万个位置，有缓存只需 1999 个。\par}\endgroup
\item \begingroup\sloppy{\fontsize{10}{14.6}\selectfont decode 受带宽限，读一遍 16.1 GB 权重至少 4.8 ms。\par}\endgroup
\end{itemize}
\meta{公式}{$p_i=\frac{\exp(z_i/T)}{\sum_{j}\exp(z_j/T)}$}
\meta{数字}{T=0.3 我 0.782；T=1 我 0.416；T=3 我 0.29；不缓存 1,499,500 位置 vs 缓存 1,999；16.1 GB、3.35 TB/s、4.8 ms、209 token/s}
\meta{易错}{温度设 0 按贪心处理，不是除以 0；答案不同是采样造成的，权重并没有变。}
\chapter{Transformer核心}
{\footnotesize\sffamily\color{grayt}收录第 25--35 讲}\par\vspace{-4pt}
\phantomsection\addcontentsline{toc}{section}{25\quad 位置编码}
\lec{25}{位置编码}
\onel{注意力本身不看顺序，靠给每个位置一个向量加到词嵌入上来补回顺序。}
\begin{itemize}[leftmargin=1.35em,itemsep=2.2pt,topsep=1pt,parsep=0pt]
\item \begingroup\sloppy{\fontsize{10}{14.6}\selectfont 自注意力按内容比较词与词，打乱输入只会让输出跟着打乱，数值完全不变。\par}\endgroup
\item \begingroup\sloppy{\fontsize{10}{14.6}\selectfont 没有位置嵌入时复读正确率只有 16.4\%，加上位置嵌入恢复到 100\%。\par}\endgroup
\item \begingroup\sloppy{\fontsize{10}{14.6}\selectfont 可学习位置嵌入（GPT-2）长度被表锁死；正弦余弦编码不需训练，位置任意长。\par}\endgroup
\item \begingroup\sloppy{\fontsize{10}{14.6}\selectfont 分母 10000\textasciicircum{}(2i/d) 越大该维变化越慢，像快慢不同的钟表指针合成唯一读数。\par}\endgroup
\item \begingroup\sloppy{\fontsize{10}{14.6}\selectfont 位置编码是相加不是拼接，维度不变，词义和位置混在同一个向量里交给后面的层。\par}\endgroup
\item \begingroup\sloppy{\fontsize{10}{14.6}\selectfont 相邻位置编码更相似：自点积 32，隔 1 为 30.92，隔 30 为 17.07。\par}\endgroup
\end{itemize}
\meta{公式}{$x_p=E[\text{token}]+P[p]$；$P(p,2i)=\sin\left(\frac{p}{10000^{2i/d}}\right)$，$P(p,2i+1)=\cos\left(\frac{p}{10000^{2i/d}}\right)$}
\meta{数字}{复读 100\% $\rightarrow$ 16.4\%；GPT-2 位置表 1024 行；d=4 时分母 1 与 100；点积 32/30.92/23.5/17.07；Llama-3 rope\_theta 10000$\rightarrow$500000}
\meta{易错}{位置编码是加到词嵌入上而非拼接；注意力公式没变，变的是输入已带位置。}
\phantomsection\addcontentsline{toc}{section}{26\quad 自注意力}
\lec{26}{自注意力}
\onel{每个词用 q 与所有 k 点积打分，softmax 成权重后对 v 加权求和。}
\begin{itemize}[leftmargin=1.35em,itemsep=2.2pt,topsep=1pt,parsep=0pt]
\item \begingroup\sloppy{\fontsize{10}{14.6}\selectfont Q=X$\cdot$Wq、K=X$\cdot$Wk、V=X$\cdot$Wv，让「我在找什么」与「我能提供什么」分开学。\par}\endgroup
\item \begingroup\sloppy{\fontsize{10}{14.6}\selectfont 打分 S=Q$\cdot$Kᵀ，一次矩阵乘法就算完所有词两两之间的分数。\par}\endgroup
\item \begingroup\sloppy{\fontsize{10}{14.6}\selectfont 除以 $\surd$d\_k 防 softmax 饱和，128 维最大权重 0.92 $\rightarrow$ 0.38。\par}\endgroup
\item \begingroup\sloppy{\fontsize{10}{14.6}\selectfont 输出是所有 v 的加权平均，每行权重和为 1，不是从句子里挑出一个词。\par}\endgroup
\item \begingroup\sloppy{\fontsize{10}{14.6}\selectfont 计算量和显存随序列长度平方增长，这是长上下文昂贵的根本原因。\par}\endgroup
\item \begingroup\sloppy{\fontsize{10}{14.6}\selectfont 未训练的 TinyLM 每行权重都在 0.2\textasciitilde{}0.3 间近乎平均，看哪里是训练出来的。\par}\endgroup
\end{itemize}
\meta{公式}{$\operatorname{Attention}(Q,K,V)=\operatorname{softmax}\left(\frac{QK^{\mathsf{T}}}{\sqrt{d_k}}\right)V$；$q$、$k$ 各分量方差为 $1$ 时 $\operatorname{Var}(q\cdot k)\approx d_k$}
\meta{数字}{点积 0.0986/0.142/-0.0606/-0.1876；除以 $\surd$4=2；点积标准差 d=4/64/128 $\rightarrow$ 2/8/11.23；注意力矩阵 32 头$\times$8192$^2$$\times$2B = 4 GiB；投影 0.69 对打分 1.10 TFLOPs}
\meta{易错}{可学的只有 Wq、Wk、Wv 三个矩阵；“看哪里”是运行时点积临时算出来的。}
\phantomsection\addcontentsline{toc}{section}{27\quad 因果掩码}
\lec{27}{因果掩码}
\onel{softmax 前把 j>i 的分数设成 -$\infty$，每个位置只能看自己和前面的词。}
\begin{itemize}[leftmargin=1.35em,itemsep=2.2pt,topsep=1pt,parsep=0pt]
\item \begingroup\sloppy{\fontsize{10}{14.6}\selectfont 不加掩码时第 i 个位置会看到第 i+1 个词，也就是它要预测的答案，直接抄就行。\par}\endgroup
\item \begingroup\sloppy{\fontsize{10}{14.6}\selectfont e\textasciicircum{}\{-$\infty$\}=0，被遮位置权重恰为 0，其余重新归一化仍和为 1，矩阵只剩含对角下三角。\par}\endgroup
\item \begingroup\sloppy{\fontsize{10}{14.6}\selectfont 改第 4 个词时前 3 个位置输出最多变 0.3804，加掩码后变化为 0。\par}\endgroup
\item \begingroup\sloppy{\fontsize{10}{14.6}\selectfont 第 i 个位置只依赖前 i 个词，前面词的 k、v 永不再变，KV Cache 因此成立。\par}\endgroup
\item \begingroup\sloppy{\fontsize{10}{14.6}\selectfont decode 时新词只算自己的 q、k、v，对整份缓存打成 1$\times$s 一行，无需掩码。\par}\endgroup
\item \begingroup\sloppy{\fontsize{10}{14.6}\selectfont 没有掩码每生成一个词都要整段重算，生成 1000 词要多处理上百万 token 位置。\par}\endgroup
\end{itemize}
\meta{公式}{$A=\operatorname{softmax}\left(\frac{QK^{\mathsf{T}}}{\sqrt{d_k}}+M\right)$，$M_{ij}=0\ (j\le i)$、$-\infty\ (j>i)$}
\meta{数字}{Llama-3-8B 每 token 2$\times$32$\times$8$\times$128$\times$2 = 128 KiB；8192 token 1 GiB；50 个请求 50 GiB，权重仅 16 GB}
\meta{易错}{掩码必须加在 softmax 之前，否则剩下的权重和不到 1，且后面的词已混进分母。}
\phantomsection\addcontentsline{toc}{section}{28\quad 多头注意力}
\lec{28}{多头注意力}
\onel{把 Q、K、V 按列切成 n\_h 段，各头独立做注意力，输出拼接后再乘 W\_o。}
\begin{itemize}[leftmargin=1.35em,itemsep=2.2pt,topsep=1pt,parsep=0pt]
\item \begingroup\sloppy{\fontsize{10}{14.6}\selectfont 切分不增加参数：四个矩阵仍是 d$\times$d，计算量几乎不变，只多几遍 softmax。\par}\endgroup
\item \begingroup\sloppy{\fontsize{10}{14.6}\selectfont 每头只拿 d/n\_h 维，头越多每头越窄，点积分辨力下降，头数并非越多越好。\par}\endgroup
\item \begingroup\sloppy{\fontsize{10}{14.6}\selectfont 复读模型头 1 学「回头看 3 格」，关掉它复读正确率从 100\% 跌到 8.1\%。\par}\endgroup
\item \begingroup\sloppy{\fontsize{10}{14.6}\selectfont 头 2 第 4--6 行几乎都盯着句号，关掉仍是 100\%，是冗余头，大模型里也常见。\par}\endgroup
\item \begingroup\sloppy{\fontsize{10}{14.6}\selectfont 头的分工由训练自己选出，从没规定头 1 负责复读；换随机种子可能换头。\par}\endgroup
\item \begingroup\sloppy{\fontsize{10}{14.6}\selectfont Llama-3 32 头 $\times$ 128 维 = 4096；GQA 让查询头共享 k、v。\par}\endgroup
\end{itemize}
\meta{公式}{输出 $=[\text{head}_1,\ldots,\text{head}_{n}]\cdot W_o$；$\text{head}_h=\operatorname{softmax}\left(\frac{Q_hK_h^{\mathsf{T}}}{\sqrt{d_h}}+M\right)V_h$，$d_h=\frac{d}{n_h}$}
\meta{数字}{TinyLM d=4 拆 2 头各 2 维，四个矩阵各 16 参数；关头 1 $\rightarrow$ 8.1\%，关头 2 $\rightarrow$ 100\%；Llama-3 32 头$\times$128 维，GQA 使 KV Cache 降到 1/4}
\meta{易错}{头的维度是切出来的不是复制出来的；W\_o 不是可有可无，它才真正混合各头。}
\phantomsection\addcontentsline{toc}{section}{29\quad 残差连接}
\lec{29}{残差连接}
\onel{每层输出改成 x+F(x)，导数多一个 1，梯度沿捷径原样传回前面的层。}
\begin{itemize}[leftmargin=1.35em,itemsep=2.2pt,topsep=1pt,parsep=0pt]
\item \begingroup\sloppy{\fontsize{10}{14.6}\selectfont 每层乘 0.5 连乘 30 次得 9.31$\times$10$^-$$^1$⁰，这就是梯度消失。\par}\endgroup
\item \begingroup\sloppy{\fontsize{10}{14.6}\selectfont 局部导数变成 1+F′，F′ 再小也还剩 1；展开后总有一条只走捷径的短路径。\par}\endgroup
\item \begingroup\sloppy{\fontsize{10}{14.6}\selectfont tanh 下普通网络第 0 层只剩 0.0506，加残差是 19.57。\par}\endgroup
\item \begingroup\sloppy{\fontsize{10}{14.6}\selectfont 残差让数值累积变大：ReLU 残差网络第 0 层梯度 3704.18，光有残差不够。\par}\endgroup
\item \begingroup\sloppy{\fontsize{10}{14.6}\selectfont Transformer 每层 2 个残差，32 层共 64 次相加。\par}\endgroup
\item \begingroup\sloppy{\fontsize{10}{14.6}\selectfont 贯穿全网的 x 叫残差流，宽度就是隐藏维度 4096，子模块输出必须同形状才能相加。\par}\endgroup
\end{itemize}
\meta{公式}{$x_{l+1}=x_l+F(x_l)$；$\frac{\partial x_{l+1}}{\partial x_l}=1+F'$}
\meta{数字}{0.5$^3$⁰=9.31$\times$10$^-$$^1$⁰；1.5$^3$⁰=1.92$\times$10⁵；tanh 普通网络 0.0506 对残差 19.57；ReLU 普通网络 0 对残差 3704.18；64 次残差相加；残差流 4096 维}
\meta{易错}{残差是相加不是拼接，拼接会让维度每层翻倍；它管梯度不管数值膨胀。}
\phantomsection\addcontentsline{toc}{section}{30\quad LayerNorm}
\lec{30}{LayerNorm}
\onel{对每个词的向量减均值、除标准差，再乘 γ 加 $\beta$，统一到均值 0 方差 1。}
\begin{itemize}[leftmargin=1.35em,itemsep=2.2pt,topsep=1pt,parsep=0pt]
\item \begingroup\sloppy{\fontsize{10}{14.6}\selectfont 均值方差在一个词的 d 个数之间算，每个词各算各的，与 batch 无关。\par}\endgroup
\item \begingroup\sloppy{\fontsize{10}{14.6}\selectfont 归一化后向量长度约 $\surd$d：TinyLM 是 2，16 维网络是 4。\par}\endgroup
\item \begingroup\sloppy{\fontsize{10}{14.6}\selectfont 不归一化时输出长度一路涨到 82441.6，加了 LayerNorm 恒为 4.0。\par}\endgroup
\item \begingroup\sloppy{\fontsize{10}{14.6}\selectfont γ、$\beta$ 各 d 维，初始为全 1 和全 0，只占 2d 个参数。\par}\endgroup
\item \begingroup\sloppy{\fontsize{10}{14.6}\selectfont 对整体缩放和平移不敏感：整行乘 3 或加 5，归一化结果都不变。\par}\endgroup
\item \begingroup\sloppy{\fontsize{10}{14.6}\selectfont 原论文放在残差相加之后（Post-Norm）；现代模型改为子模块之前并用 RMSNorm。\par}\endgroup
\end{itemize}
\meta{公式}{$\operatorname{LN}(x)=\gamma\odot\frac{x-\mu}{\sqrt{\sigma^2+\varepsilon}}+\beta$，$\mu$ 和 $\sigma^2$ 在每行的 $d$ 个维度上算}
\meta{数字}{AI 行均值 -0.0434、方差 0.4165、标准差 0.6454；ε=1e-5；归一化前 3.4/42.9/2006.4/82441.6，归一化后全 4.0；Llama-3 共 65 个归一化 266240 参数}
\meta{易错}{归一化后不是所有数都在 -1\textasciitilde{}1，均值 0 方差 1 只是统计量，单个数可到 1.3123。}
\phantomsection\addcontentsline{toc}{section}{31\quad 前馈网络}
\lec{31}{前馈网络}
\onel{FFN 对每个词独立做升维$\rightarrow$ReLU$\rightarrow$降维，词间不交流，是参数与计算的大头。}
\begin{itemize}[leftmargin=1.35em,itemsep=2.2pt,topsep=1pt,parsep=0pt]
\item \begingroup\sloppy{\fontsize{10}{14.6}\selectfont W1 从 d 升到 d\_ff（约 4 倍），W2 降回 d 维才能与残差相加。\par}\endgroup
\item \begingroup\sloppy{\fontsize{10}{14.6}\selectfont 每个词独立算，整句一起算与逐个词算结果相同，原论文叫 position-wise。\par}\endgroup
\item \begingroup\sloppy{\fontsize{10}{14.6}\selectfont 可看成键---值表：W1 每列是一个待匹配模式，W2 对应行是匹配后写回的内容。\par}\endgroup
\item \begingroup\sloppy{\fontsize{10}{14.6}\selectfont TinyLM 的 d\_ff=8，8 个神经元只亮 4 个，按亮度加权 W2 的行即得输出。\par}\endgroup
\item \begingroup\sloppy{\fontsize{10}{14.6}\selectfont d\_ff=14336（3.5 倍），SwiGLU 三矩阵约占每层参数 81\%。\par}\endgroup
\item \begingroup\sloppy{\fontsize{10}{14.6}\selectfont 每 token 过一层 FFN 约 3.52 亿次浮点运算，序列不长时远超注意力。\par}\endgroup
\end{itemize}
\meta{公式}{$\operatorname{FFN}(x)=\operatorname{ReLU}(x\cdot W_1)\cdot W_2=\sum_{j}\operatorname{ReLU}(x\cdot W_1[:,j])\,W_2[j,:]$}
\meta{数字}{TinyLM FFN 64 参数 = 注意力 64；Llama-3 FFN 176.2M 参数、每 token 352M 次运算；d\_ff/d = 14336/4096 = 3.5}
\meta{易错}{d\_ff 不一定 4 倍：SwiGLU 为保持参数量常用 8/3 倍，Llama-3 是 3.5 倍。}
\phantomsection\addcontentsline{toc}{section}{32\quad Transformer Block}
\lec{32}{Transformer Block}
\onel{一层 = 注意力子层 + FFN 子层，各带残差与归一化，形状不变可叠 N 层。}
\begin{itemize}[leftmargin=1.35em,itemsep=2.2pt,topsep=1pt,parsep=0pt]
\item \begingroup\sloppy{\fontsize{10}{14.6}\selectfont 两个子层顺序固定：先注意力后 FFN，各自「子层计算$\rightarrow$加残差$\rightarrow$归一化」。\par}\endgroup
\item \begingroup\sloppy{\fontsize{10}{14.6}\selectfont 输入输出都是 [s,d]，所以层能一层接一层叠；结构相同但每层权重不共享。\par}\endgroup
\item \begingroup\sloppy{\fontsize{10}{14.6}\selectfont 数据流是嵌入加位置编码，过 N 个块，再经 lm\_head 与 softmax 出概率。\par}\endgroup
\item \begingroup\sloppy{\fontsize{10}{14.6}\selectfont TinyLM 304 个参数：嵌入 24、每层 128、输出层 24。\par}\endgroup
\item \begingroup\sloppy{\fontsize{10}{14.6}\selectfont Llama-3-8B 参数 8,030,261,248，与官方 8B 对上。\par}\endgroup
\item \begingroup\sloppy{\fontsize{10}{14.6}\selectfont 现代模型换 Pre-Norm、RMSNorm、RoPE、SwiGLU、GQA。\par}\endgroup
\end{itemize}
\meta{公式}{Post-Norm：$x\leftarrow\operatorname{LN}(x+\operatorname{Attn}(x))$，$x\leftarrow\operatorname{LN}(x+\operatorname{FFN}(x))$；Pre-Norm：$x\leftarrow x+\operatorname{Attn}(\operatorname{Norm}(x))$，$x\leftarrow x+\operatorname{FFN}(\operatorname{Norm}(x))$}
\meta{数字}{TinyLM 304；Llama-3-8B 每层 218,112,000，32 层 6,979,584,000，合计 8,030,261,248；70B 是 80 层宽 8192}
\meta{易错}{“32 层”指 32 个 Block，每块 7 个线性层加 2 个归一化；结构同、权重是 32 份。}
\phantomsection\addcontentsline{toc}{section}{33\quad 输出层}
\lec{33}{输出层}
\onel{输出层用 d$\times$V 的 lm\_head 加 softmax，给出词表概率。}
\begin{itemize}[leftmargin=1.35em,itemsep=2.2pt,topsep=1pt,parsep=0pt]
\item \begingroup\sloppy{\fontsize{10}{14.6}\selectfont lm\_head 是 d$\times$V 矩阵，一列即一个词的输出向量，logit 是点积。\par}\endgroup
\item \begingroup\sloppy{\fontsize{10}{14.6}\selectfont 整句算：[s,d] 乘 [d,V] 得 [s,V]，每位置一行 V 个打分。\par}\endgroup
\item \begingroup\sloppy{\fontsize{10}{14.6}\selectfont 推理只算最后一行，训练时每个位置都算，logits 需数 GiB 显存。\par}\endgroup
\item \begingroup\sloppy{\fontsize{10}{14.6}\selectfont TinyLM 的 Wlm 4$\times$6，“AI”6 个打分最高的是“推理”1.184。\par}\endgroup
\item \begingroup\sloppy{\fontsize{10}{14.6}\selectfont 权重绑定令 Wlm = Eᵀ 省一份参数，GPT-2 这样做，Llama-3 不绑定。\par}\endgroup
\item \begingroup\sloppy{\fontsize{10}{14.6}\selectfont Llama-3-8B 的 lm\_head 有 5.25 亿参数，比任何单个线性层都大。\par}\endgroup
\end{itemize}
\meta{公式}{$\text{logits}=x\cdot W_{\text{lm}}$；$p_i=\frac{e^{\text{logit}_i}}{\sum_{j}e^{\text{logit}_j}}$}
\meta{数字}{TinyLM Wlm 4$\times$6；softmax 后“推理”0.441；Llama-3-8B 4096$\times$128256 = 525M 参数；每 token 1.05 GFLOPs，占 2N 的 6.5\%；一个位置 logits 用 FP32 存约 501 KiB}
\meta{易错}{logits 不是概率：可正可负、加起来不为 1；Pre-Norm 模型在 lm\_head 前还要补一个最终 RMSNorm。}
\phantomsection\addcontentsline{toc}{section}{34\quad 前向传播}
\lec{34}{前向传播}
\onel{一次前向=查表、位置编码、N 个 Block、lm\_head、softmax。}
\begin{itemize}[leftmargin=1.35em,itemsep=2.2pt,topsep=1pt,parsep=0pt]
\item \begingroup\sloppy{\fontsize{10}{14.6}\selectfont 24 站含：查嵌入、加位置编码、每层 QKV、注意力、拼头、残差、LN、FFN。\par}\endgroup
\item \begingroup\sloppy{\fontsize{10}{14.6}\selectfont 第 2 到 22 站形状始终是 [4,4]，只有 FFN 中间变宽到 8 维。\par}\endgroup
\item \begingroup\sloppy{\fontsize{10}{14.6}\selectfont 因果掩码使“我爱”最后一行等于整句第 2 行，前缀可缓存。\par}\endgroup
\item \begingroup\sloppy{\fontsize{10}{14.6}\selectfont TinyLM 预测“推理”0.441，贪心生成得“我推理我推理”。\par}\endgroup
\item \begingroup\sloppy{\fontsize{10}{14.6}\selectfont Llama-3-8B 同一路径，每 token 约 160 亿 FLOPs。\par}\endgroup
\item \begingroup\sloppy{\fontsize{10}{14.6}\selectfont KV Cache 省 K/V 重算，FlashAttention 省掉注意力矩阵写显存。\par}\endgroup
\end{itemize}
\meta{公式}{$x=E[\text{ids}]+P$；每层 $x=\operatorname{LN}(x+\operatorname{Attn}(x))$，$x=\operatorname{LN}(x+\operatorname{FFN}(x))$；$p=\operatorname{softmax}(x\cdot W_{\text{lm}})$}
\meta{数字}{24 站；TinyLM d=4、2 层、2 头、FFN 8 维、词表 6；Llama-3-8B 每 token 约 2$\times$80 亿 = 160 亿 FLOPs}
\meta{易错}{一次前向不等于推理：推理是把前向放进循环（prefill 一次、decode 多次）；训练在前向之后还有反向。}
\phantomsection\addcontentsline{toc}{section}{35\quad 参数量与计算量}
\lec{35}{参数量与计算量}
\onel{每 token 推理$\approx$2N、训练$\approx$6N FLOPs，参数按结构逐项相加。}
\begin{itemize}[leftmargin=1.35em,itemsep=2.2pt,topsep=1pt,parsep=0pt]
\item \begingroup\sloppy{\fontsize{10}{14.6}\selectfont 每层注意力 2d$^2$+2d$\cdot$kv；FFN 用 SwiGLU 是 3d$\cdot$d\_ff。\par}\endgroup
\item \begingroup\sloppy{\fontsize{10}{14.6}\selectfont N 按嵌入表、每层注意力与 FFN、输出层逐项相加，结果与官方一致。\par}\endgroup
\item \begingroup\sloppy{\fontsize{10}{14.6}\selectfont 参数里 FFN 占 70.2\%、注意力 16.7\%，嵌入与输出层 13.1\%。\par}\endgroup
\item \begingroup\sloppy{\fontsize{10}{14.6}\selectfont 每个参与矩阵乘法的参数对每 token 贡献 2 FLOPs，前向$\approx$2(N-V$\cdot$d)。\par}\endgroup
\item \begingroup\sloppy{\fontsize{10}{14.6}\selectfont 上下文 8192 时注意力打分占 22.2\%，131072 时占 82.1\%。\par}\endgroup
\item \begingroup\sloppy{\fontsize{10}{14.6}\selectfont 2N、6N 忽略嵌入、打分、归一化与 softmax，只在上下文不长时准。\par}\endgroup
\end{itemize}
\meta{公式}{$N=Vd+L(2d^2+2d\cdot\text{kv}+3d\,d_{\text{ff}}+2d)+d+Vd$；每 token 前向 $\approx 2(N-Vd)+4Lsd$；训练 $\approx 6N$}
\meta{数字}{Llama-3-8B 8,030,261,248；Llama-3-70B 70,553,706,496；8B 每 token 15.01(s=0)/19.30(s=8192) GFLOPs，训练 48.18；BF16 权重 14.96 GiB$\approx$16.1 GB；KV Cache 每 token 128 KiB，8192$\times$16 请求 16 GiB}
\meta{易错}{FLOPs 是运算次数、FLOPS 是每秒次数；按 MACs 统计会差一半；MoE 的 2N 要用激活参数（V3 总 6710 亿、激活约 370 亿）。}
\chapter{现代LLM架构}
{\footnotesize\sffamily\color{grayt}收录第 36--44 讲}\par\vspace{-4pt}
\phantomsection\addcontentsline{toc}{section}{36\quad Pre-Norm 与 RMSNorm}
\lec{36}{Pre-Norm 与 RMSNorm}
\onel{归一化挪到子层入口成 Pre-Norm，并换成 RMSNorm。}
\begin{itemize}[leftmargin=1.35em,itemsep=2.2pt,topsep=1pt,parsep=0pt]
\item \begingroup\sloppy{\fontsize{10}{14.6}\selectfont Post-Norm 是 x$\leftarrow$LN(x+F(x))，每层捷径后都接归一化，梯度被逐层缩放。\par}\endgroup
\item \begingroup\sloppy{\fontsize{10}{14.6}\selectfont Pre-Norm 是 x$\leftarrow$x+F(Norm(x))，残差流成纯加法链，梯度直达输入。\par}\endgroup
\item \begingroup\sloppy{\fontsize{10}{14.6}\selectfont RMSNorm 不减均值、没有 $\beta$，只用 γ 按 $\surd$(mean(x$^2$)+ϵ) 缩放。\par}\endgroup
\item \begingroup\sloppy{\fontsize{10}{14.6}\selectfont 同一组数两者结果接近但不相同，差别来自 -0.0434 的均值。\par}\endgroup
\item \begingroup\sloppy{\fontsize{10}{14.6}\selectfont Fused Add-RMSNorm 省掉一半以上 HBM 往返，残差与归一化同内核完成。\par}\endgroup
\item \begingroup\sloppy{\fontsize{10}{14.6}\selectfont FP16 上限 65504，求平方和会溢出，须在寄存器内 upcast 到 FP32。\par}\endgroup
\end{itemize}
\meta{公式}{Post-Norm: $x\leftarrow\operatorname{LN}(x+F(x))$；Pre-Norm: $x\leftarrow x+F(\operatorname{Norm}(x))$；$\operatorname{RMSNorm}(x)=\frac{\gamma\odot x}{\sqrt{\operatorname{mean}(x^2)+\epsilon}}$}
\meta{数字}{96 层回传梯度 Post-Norm 5508.33、Pre-Norm 21.374；同一行 LayerNorm 得 (-0.5327, -1.3091, 0.5295, 1.3123)，RMSNorm 得 (-0.5986, -1.3732, 0.4612, 1.2422)；Llama-3-8B RMSNorm 262,144 参数(1.0 MiB)，LayerNorm 524,288(2.0 MiB)}
\meta{易错}{RMSNorm 不对整体平移不变（平移会改变输出）；Gemma 用 x*(1+gamma) 且 gamma 初始化为 0，按全 1 加载会错乱。}
\phantomsection\addcontentsline{toc}{section}{37\quad RoPE 旋转位置编码}
\lec{37}{RoPE 旋转位置编码}
\onel{RoPE 在注意力里按位置旋转 q、k，使 q$\cdot$k 只取决于位置差。}
\begin{itemize}[leftmargin=1.35em,itemsep=2.2pt,topsep=1pt,parsep=0pt]
\item \begingroup\sloppy{\fontsize{10}{14.6}\selectfont 旋转不改变向量长度，只改变方向；二维旋转用 cos、sin 的矩阵。\par}\endgroup
\item \begingroup\sloppy{\fontsize{10}{14.6}\selectfont q 转 m$\theta$、k 转 n$\theta$，夹角只变 (m-n)$\theta$，点积只取决于位置差。\par}\endgroup
\item \begingroup\sloppy{\fontsize{10}{14.6}\selectfont d 维配成 d/2 个平面，第 i 对转速 base\textasciicircum{}(-2i/d)，与正弦编码同源。\par}\endgroup
\item \begingroup\sloppy{\fontsize{10}{14.6}\selectfont 位置差 2 的三组点积都是 0.0069，差 3 的 (3,0) 是 -0.1178。\par}\endgroup
\item \begingroup\sloppy{\fontsize{10}{14.6}\selectfont 只旋转 q、k，不碰 v；KV Cache 存旋转后的 k，cos/sin 预算成查表。\par}\endgroup
\item \begingroup\sloppy{\fontsize{10}{14.6}\selectfont 远距离衰减：q=k=全 1 的 128 维，距离 1000 时点积降到 20.4。\par}\endgroup
\end{itemize}
\meta{公式}{旋转 $\varphi$：$(a,b)\to(a\cos\varphi-b\sin\varphi,\ a\sin\varphi+b\cos\varphi)$；$\theta_i=\text{base}^{-2i/d}$}
\meta{数字}{TinyLM 4 维 2 对，$\theta$0=1、$\theta$1=0.01；base 10000 时最慢一对 54410 个位置转一圈；Llama-3 用 500000，最慢约 2,559,196 个位置；最快一对 6.3 个位置一圈；Llama-3 每头 128 维即 64 对；衰减 128.0/124.2/85.6/61.1/20.4/-3.6 对应距离 0/1/10/100/1000/10000}
\meta{易错}{两两配对有交替配对与半半拆分两种，数学等价但旋转基准不同，搞反会直接输出乱码；Hugging Face 的 Llama 用半半拆分。}
\phantomsection\addcontentsline{toc}{section}{38\quad SwiGLU 与门控 FFN}
\lec{38}{SwiGLU 与门控 FFN}
\onel{门控 FFN 把升维拆成门 g 与内容 u，SiLU(g)$\odot$u 后降维。}
\begin{itemize}[leftmargin=1.35em,itemsep=2.2pt,topsep=1pt,parsep=0pt]
\item \begingroup\sloppy{\fontsize{10}{14.6}\selectfont 两路升维 g=x$\cdot$W\_gate、u=x$\cdot$W\_up，逐元素相乘后乘 W\_down 降维。\par}\endgroup
\item \begingroup\sloppy{\fontsize{10}{14.6}\selectfont SiLU(z)=z$\cdot$σ(z) 给出连续开度，可正可负，最低约 -0.28。\par}\endgroup
\item \begingroup\sloppy{\fontsize{10}{14.6}\selectfont 门开 0.886 的神经元几乎全放行 1.455，门近 0 的输出接近 0。\par}\endgroup
\item \begingroup\sloppy{\fontsize{10}{14.6}\selectfont Shazeer 2020 年论文：等参数下带门的版本都优于 ReLU、GELU。\par}\endgroup
\item \begingroup\sloppy{\fontsize{10}{14.6}\selectfont 等参数要 d\_ff$\approx$8d/3；d=4096 时取 10922，参数 134.2M。\par}\endgroup
\item \begingroup\sloppy{\fontsize{10}{14.6}\selectfont W\_gate 与 W\_up 拼成 [d,2d\_ff] 一次 GEMM，激活在片上完成。\par}\endgroup
\end{itemize}
\meta{公式}{$g=x\cdot W_{\text{gate}}$，$u=x\cdot W_{\text{up}}$，$h=\operatorname{SiLU}(g)\odot u$，$y=h\cdot W_{\text{down}}$；等参数时 $d_{\text{ff}}\approx\frac{8d}{3}$}
\meta{数字}{TinyLM 8 个神经元，y=(-0.347, 1.904, 0.364, -1.628)；门偏移 -3/+3 时 |h| 之和 1.271 对 23.745；Llama-3-8B d\_ff=14336(3.5d)，每层 176.2M 参数，权重拼成 4096$\times$28672}
\meta{易错}{门不是 0/1 开关；GLU 是两路相乘的结构、GELU 是激活函数，门用 GELU 叫 GEGLU，用 SiLU 才叫 SwiGLU；内容 u 不过激活。}
\phantomsection\addcontentsline{toc}{section}{39\quad MQA 与 GQA}
\lec{39}{MQA 与 GQA}
\onel{MQA 全部共用一组 K、V，GQA 分组共用，缓存随组数缩小。}
\begin{itemize}[leftmargin=1.35em,itemsep=2.2pt,topsep=1pt,parsep=0pt]
\item \begingroup\sloppy{\fontsize{10}{14.6}\selectfont MHA 每 token 512 KiB，16 个长请求就是 64 GiB。\par}\endgroup
\item \begingroup\sloppy{\fontsize{10}{14.6}\selectfont MQA 只留一组 K、V（2019，Shazeer），缓存降到 1/32，效果有损。\par}\endgroup
\item \begingroup\sloppy{\fontsize{10}{14.6}\selectfont GQA 把 32 个查询头分 8 组，每组 4 头共用，每 token 128 KiB。\par}\endgroup
\item \begingroup\sloppy{\fontsize{10}{14.6}\selectfont 查询头与注意力权重不变，变少的只是被缓存的 K、V；GQA 介于两者。\par}\endgroup
\item \begingroup\sloppy{\fontsize{10}{14.6}\selectfont 底层不复制 KV：只留 n\_kv 份，线程按组号索引在 SRAM 复用。\par}\endgroup
\item \begingroup\sloppy{\fontsize{10}{14.6}\selectfont 算术强度从 1 升到 4 FLOP/Byte：读一个 k 服务同组 4 个头。\par}\endgroup
\end{itemize}
\meta{公式}{$\text{kv\_group}=\frac{h}{n_h/n_{kv}}$；每 token KV Cache $=2\,L\,n_{kv}\,d_{\text{head}}\cdot\text{字节数}$}
\meta{数字}{Llama-3-8B：MHA 每 token 512 KiB、GQA(8) 128 KiB、MQA 16 KiB；8192 token$\times$16 请求 64/16/2 GiB；W\_k 与 W\_v 每层 33.6M/8.4M/1.0M 参数}
\meta{易错}{把 MHA 的头直接取平均合并成 GQA 会明显掉效果，要再训练约原训练量 5\%；KV 头数还要能被张量并行度整除，8 头切不了 16 卡。}
\phantomsection\addcontentsline{toc}{section}{40\quad MLA 多头潜在注意力}
\lec{40}{MLA 多头潜在注意力}
\onel{MLA 不缓存 K、V，每 token 只存 512+64=576 个数。}
\begin{itemize}[leftmargin=1.35em,itemsep=2.2pt,topsep=1pt,parsep=0pt]
\item \begingroup\sloppy{\fontsize{10}{14.6}\selectfont 先压缩再还原：c=x$\cdot$W\_DKV，k=c$\cdot$W\_UK，v=c$\cdot$W\_UV，只缓存 c。\par}\endgroup
\item \begingroup\sloppy{\fontsize{10}{14.6}\selectfont 直接缓存 K、V 时每层每 token 是 32768 个数，而 c 只有 512 维。\par}\endgroup
\item \begingroup\sloppy{\fontsize{10}{14.6}\selectfont 矩阵吸收：q$\cdot$k=(q$\cdot$W\_UKᵀ)$\cdot$c，Decode 时历史 K、V 从不展开。\par}\endgroup
\item \begingroup\sloppy{\fontsize{10}{14.6}\selectfont 解耦 RoPE：128 维内容部分不旋转可吸收，另留 64 维 k\_R 做位置。\par}\endgroup
\item \begingroup\sloppy{\fontsize{10}{14.6}\selectfont Prefill 展开成多头调 FlashAttention，Decode 走吸收省带宽。\par}\endgroup
\item \begingroup\sloppy{\fontsize{10}{14.6}\selectfont 只有 576 个数，比 8 组 GQA 小 3.56 倍，效果反超同尺寸 MHA。\par}\endgroup
\end{itemize}
\meta{公式}{$c=x\cdot W_{\text{DKV}}$，$k=c\cdot W_{\text{UK}}$，$v=c\cdot W_{\text{UV}}$；$q\cdot k=(q\cdot W_{\text{UK}}^{\mathsf{T}})\cdot c+q_R\cdot k_R$}
\meta{数字}{DeepSeek-V3：d=7168、61 层、128 头、每头 k/v 各 128 维；c 512 维、k\_R 64 维；每层每 token 576 个数，61 层 68.6 KiB；MHA 3904 KiB(56.89$\times$)、GQA 244 KiB(3.56$\times$)、MQA 30.5 KiB(0.44$\times$)；q\_lora\_rank 1536}
\meta{易错}{MLA 不减少头数，每个头仍有独立 K、V，只是约束在低秩子空间；吸收是纯代数等价，不改变模型。}
\phantomsection\addcontentsline{toc}{section}{41\quad MoE 混合专家}
\lec{41}{MoE 混合专家}
\onel{FFN 拆成 E 个专家，每 token 只挑 k 个算，总参数与计算量解耦。}
\begin{itemize}[leftmargin=1.35em,itemsep=2.2pt,topsep=1pt,parsep=0pt]
\item \begingroup\sloppy{\fontsize{10}{14.6}\selectfont 路由器 d$\times$E 线性层出 E 个分数，取 Top-K 后仅对 k 个归一化。\par}\endgroup
\item \begingroup\sloppy{\fontsize{10}{14.6}\selectfont 未选中的专家对该 token 完全旁路，不产生任何浮点运算。\par}\endgroup
\item \begingroup\sloppy{\fontsize{10}{14.6}\selectfont 底层三步：重排 token、Grouped GEMM、Scatter-Add 写回。\par}\endgroup
\item \begingroup\sloppy{\fontsize{10}{14.6}\selectfont 共享专家不参与路由、所有 token 必经，承担通用特征，省下路由专家容量。\par}\endgroup
\item \begingroup\sloppy{\fontsize{10}{14.6}\selectfont Mixtral 8x7B 总 467 亿、激活 129 亿，算力等同 130 亿稠密模型。\par}\endgroup
\item \begingroup\sloppy{\fontsize{10}{14.6}\selectfont 偏置只进 Top-K 挑选、不进 Softmax 权重，最忙/平均 1.81$\rightarrow$1.12。\par}\endgroup
\end{itemize}
\meta{公式}{$\operatorname{MoE}(x)=\operatorname{FFN}_{\text{shared}}(x)+\sum_{e\in\text{TopK}}w_{e}\operatorname{FFN}_{e}(x)$，$h(x)=x\cdot W_{r}$}
\meta{数字}{DeepSeek-V3 总 6710 亿/激活 370 亿（5.5\%），每层 1 共享+256 路由选 8；Mixtral 8x7B 467 亿/129 亿（28\%）；权重 FP8 约 671 GB}
\meta{易错}{2N 里的 N 要用激活参数，路由器、注意力打分、All-to-All 通信都不在内，故比同激活参数的稠密模型略慢。}
\phantomsection\addcontentsline{toc}{section}{42\quad 长上下文技术}
\lec{42}{长上下文技术}
\onel{位置靠 YaRN 分治；显存计算靠滑动窗口加 Attention Sink。}
\begin{itemize}[leftmargin=1.35em,itemsep=2.2pt,topsep=1pt,parsep=0pt]
\item \begingroup\sloppy{\fontsize{10}{14.6}\selectfont 128K 时打分占 82\%，1M 时占 97\%，KV 到 128 GiB。\par}\endgroup
\item \begingroup\sloppy{\fontsize{10}{14.6}\selectfont 直接外推把慢转维度的转角推出训练范围；位置插值乘 1/4 拉回。\par}\endgroup
\item \begingroup\sloppy{\fontsize{10}{14.6}\selectfont YaRN：高频不插值、低频完全插值、中频平滑过渡，另加温度补偿。\par}\endgroup
\item \begingroup\sloppy{\fontsize{10}{14.6}\selectfont 滑动窗口只看最近 W 个，KV 上限 W；叠 L 层信息最远传 L$\times$W。\par}\endgroup
\item \begingroup\sloppy{\fontsize{10}{14.6}\selectfont softmax 权重和必须为 1，多余注意力被倒到开头 4 个 token。\par}\endgroup
\item \begingroup\sloppy{\fontsize{10}{14.6}\selectfont 丢掉开头几个 token，权重泛滥到窗口内正常词上，生成立刻崩溃复读。\par}\endgroup
\end{itemize}
\meta{公式}{位置插值 $\text{pos}\times\frac{L_{\text{train}}}{L_{\text{target}}}$；最慢维转角 $0.0201\to0.0805$（外推）$\to0.0201$（插值）；base 改 $2{,}044{,}497$}
\meta{数字}{8,192$\rightarrow$1.0 GiB/19.3 GFLOPs/22\%；131,072$\rightarrow$16.0 GiB/83.7/82\%；1,048,576$\rightarrow$128.0 GiB/564.8/97\%；窗口 4096 只需 0.5 GiB，32 层接力最远 131,072}
\meta{易错}{位置插值后还要用少量长文本继续训练；直接压缩就上线效果会明显掉。}
\phantomsection\addcontentsline{toc}{section}{43\quad 架构演进对照}
\lec{43}{架构演进对照}
\onel{骨架没变，两条主线：每层更好用、推理更省，KV 压到 68.6 KiB。}
\begin{itemize}[leftmargin=1.35em,itemsep=2.2pt,topsep=1pt,parsep=0pt]
\item \begingroup\sloppy{\fontsize{10}{14.6}\selectfont 骨架五年不变：仅解码器、Pre-Norm、残差、预测下一 token。\par}\endgroup
\item \begingroup\sloppy{\fontsize{10}{14.6}\selectfont GPT-2 XL：MHA 25 头、GELU 4 倍、可学习位置嵌入、线性层带偏置。\par}\endgroup
\item \begingroup\sloppy{\fontsize{10}{14.6}\selectfont Llama-3：RMSNorm、RoPE、GQA、SwiGLU、去偏置。\par}\endgroup
\item \begingroup\sloppy{\fontsize{10}{14.6}\selectfont DeepSeek-V3：MLA 缓存 512+64 维、MoE 256 选 8、FP8。\par}\endgroup
\item \begingroup\sloppy{\fontsize{10}{14.6}\selectfont KV Cache 300$\rightarrow$128$\rightarrow$68.6 KiB，参数涨 400 倍缓存反降。\par}\endgroup
\item \begingroup\sloppy{\fontsize{10}{14.6}\selectfont 比推理成本看激活参数与 KV Cache，不看总参数。\par}\endgroup
\end{itemize}
\meta{数字}{15.6 亿/80.3 亿/6710 亿（激活 370 亿）；48/32/61 层；词表 50257/128256/129280；KV 300/128/68.6 KiB；Llama-2-7B 为 MHA、70B 才 GQA}
\meta{易错}{读新模型先看 config.json：num\_key\_value\_heads 定 GQA、kv\_lora\_rank 定 MLA、num\_experts 定 MoE、rope\_theta 与 max\_position\_embeddings 定上下文。}
\phantomsection\addcontentsline{toc}{section}{44\quad 混合架构与线性注意力}
\lec{44}{混合架构与线性注意力}
\onel{线性注意力把历史压进 d$\times$d 状态、每步与 n 无关，混合架构补精确召回。}
\begin{itemize}[leftmargin=1.35em,itemsep=2.2pt,topsep=1pt,parsep=0pt]
\item \begingroup\sloppy{\fontsize{10}{14.6}\selectfont 去 softmax 后结合律生效：(QK\textasciicircum{}T)V = Q(K\textasciicircum{}T V)，后者只算 d$\times$d。\par}\endgroup
\item \begingroup\sloppy{\fontsize{10}{14.6}\selectfont 状态 S 每步累加 φ(k\_t)\textasciicircum{}T v\_t，输出用 φ(q\_t) 读一次，代价 d$^2$。\par}\endgroup
\item \begingroup\sloppy{\fontsize{10}{14.6}\selectfont 总计算普通约 n$^2$d、线性约 nd$^2$；n=131072 时差约一千倍。\par}\endgroup
\item \begingroup\sloppy{\fontsize{10}{14.6}\selectfont 16$\times$16 状态存 4/16/64 对，取回误差 0.67/0.95/1.83。\par}\endgroup
\item \begingroup\sloppy{\fontsize{10}{14.6}\selectfont Mamba：Δ\_t 随输入变可跳过无用 token，状态留在 SRAM 不写 HBM。\par}\endgroup
\item \begingroup\sloppy{\fontsize{10}{14.6}\selectfont 7 线性 + 1 全注意力混合，131K 下 KV 32$\rightarrow$4.0 MiB/层，省 88\%。\par}\endgroup
\end{itemize}
\meta{公式}{$S_{t}=S_{t-1}+\varphi(k_{t})^{\top}v_{t}$，$o_{t}\propto\varphi(q_{t})\cdot S_{t}$；$(QK^{\top})V=Q(K^{\top}V)$}
\meta{数字}{n=131,072、d=128：普通 134.2 MFLOPs、KV 32.0 MiB/层；线性 0.1 MFLOPs、状态 16,384 个数；混合 16.8 MFLOPs、4.0 MiB；n=131072 总量 2.2$\times$10$^1$$^2$ 对 2.15$\times$10⁹}
\meta{易错}{线性注意力不是把注意力做成线性的，它仍在做软寻址，只是靠核函数让矩阵乘法结合律成立。}
\chapter{训练全流程}
{\footnotesize\sffamily\color{grayt}收录第 45--51 讲}\par\vspace{-4pt}
\phantomsection\addcontentsline{toc}{section}{45\quad 预训练与数据}
\lec{45}{预训练与数据}
\onel{网页为主按配比混合，过抽取去重质量安全过滤，数据质量决定能力。}
\begin{itemize}[leftmargin=1.35em,itemsep=2.2pt,topsep=1pt,parsep=0pt]
\item \begingroup\sloppy{\fontsize{10}{14.6}\selectfont 来源以 Common Crawl 网页为主，加书籍百科、代码、数学。\par}\endgroup
\item \begingroup\sloppy{\fontsize{10}{14.6}\selectfont Llama-3 配比：约半通用、1/4 数学推理、17\% 代码、8\% 多语言。\par}\endgroup
\item \begingroup\sloppy{\fontsize{10}{14.6}\selectfont 清洗：抽取$\rightarrow$语言识别$\rightarrow$完全与近似去重$\rightarrow$质量过滤$\rightarrow$安全$\rightarrow$配比分词。\par}\endgroup
\item \begingroup\sloppy{\fontsize{10}{14.6}\selectfont 近似去重用 3 字片段 Jaccard、阈值 0.7；十亿级靠 MinHash 指纹。\par}\endgroup
\item \begingroup\sloppy{\fontsize{10}{14.6}\selectfont 去重不足会让模型死记重复文本、白烧算力，去重好则更少 token 更优。\par}\endgroup
\item \begingroup\sloppy{\fontsize{10}{14.6}\selectfont 9 篇示意文档过流水线只剩 4 篇；15T token 约 60 TB。\par}\endgroup
\end{itemize}
\meta{数字}{15T token $\approx$ 60 TB；9$\rightarrow$4 篇（3 中文 1 英文）；Jaccard 0.79 > 0.7；广告重复片段占 61\%；每秒 1000 token 读 476 年}
\meta{易错}{预训练数据里没有答案，只教接着往下写；且知识有截止时间，之后的事只能靠检索或联网补。}
\phantomsection\addcontentsline{toc}{section}{46\quad Scaling Laws}
\lec{46}{Scaling Laws}
\onel{损失随 N、D 按幂律降；C$\approx$6ND 时两者一起放大，最优在 U 形底。}
\begin{itemize}[leftmargin=1.35em,itemsep=2.2pt,topsep=1pt,parsep=0pt]
\item \begingroup\sloppy{\fontsize{10}{14.6}\selectfont 1.69 是语言本身的损失下限，N、D 两项各按 0.34、0.28 次幂衰减。\par}\endgroup
\item \begingroup\sloppy{\fontsize{10}{14.6}\selectfont 固定 C=10$^2$$^2$：最优 5.15B 配 323B token，每参数 62.8 个。\par}\endgroup
\item \begingroup\sloppy{\fontsize{10}{14.6}\selectfont 算力 10$^2$⁰$\rightarrow$10$^2$$^4$，最优 0.64B$\rightarrow$41.35B、26B$\rightarrow$4030B，一起放大。\par}\endgroup
\item \begingroup\sloppy{\fontsize{10}{14.6}\selectfont Chinchilla 结论：模型与数据一起放大，偏任一边损失都变大。\par}\endgroup
\item \begingroup\sloppy{\fontsize{10}{14.6}\selectfont Llama-3-8B 每参数 1868 token，远超同算力最优点，是有意过度训练。\par}\endgroup
\item \begingroup\sloppy{\fontsize{10}{14.6}\selectfont 训练一次推理千万次，每 token 推理 2N，故宁多花训练算力喂小模型。\par}\endgroup
\end{itemize}
\meta{公式}{$L(N,D)=1.69+\frac{406.4}{N^{0.34}}+\frac{410.7}{D^{0.28}}$；$C\approx6ND$}
\meta{数字}{C=10$^2$$^2$$\rightarrow$5.15B/323B/2.139；C=10$^2$$^4$$\rightarrow$41.35B/4030B/1.911；Llama-3-8B C=7.23$\times$10$^2$$^3$、1868 token/参数、损失 1.948；同算力最优 36B/3.4T、损失 1.923、推理贵 4.4 倍；每参数 token 40.2/62.8/97.5}
\meta{易错}{常数 E、A、B、$\alpha$、$\beta$ 依赖数据、分词器与结构，换数据要重拟合；能推广的只是幂律形式和一起放大的结论。}
\phantomsection\addcontentsline{toc}{section}{47\quad 监督微调SFT}
\lec{47}{监督微调SFT}
\onel{用提示---回答对继续预测下一 token，模板标角色，掩码只对回答算损失。}
\begin{itemize}[leftmargin=1.35em,itemsep=2.2pt,topsep=1pt,parsep=0pt]
\item \begingroup\sloppy{\fontsize{10}{14.6}\selectfont 与预训练同法，只差对话模板和只对回答算损失两点。\par}\endgroup
\item \begingroup\sloppy{\fontsize{10}{14.6}\selectfont 数据几万到几百万条，几千条高质量即可，质量与多样性重于数量。\par}\endgroup
\item \begingroup\sloppy{\fontsize{10}{14.6}\selectfont 模板用特殊 token 标角色与轮次（<|eot\_id|> 结束一轮），推理须一致。\par}\endgroup
\item \begingroup\sloppy{\fontsize{10}{14.6}\selectfont 损失掩码按位置给 0/1，只在回答算交叉熵；提示照常前向。\par}\endgroup
\item \begingroup\sloppy{\fontsize{10}{14.6}\selectfont ToyLM 只训输出层 50 步：你学 0.67$\rightarrow$AI 1.00，未教的我学也变。\par}\endgroup
\item \begingroup\sloppy{\fontsize{10}{14.6}\selectfont SFT 教格式不教知识；全参数更新，显存与预训练同级。\par}\endgroup
\end{itemize}
\meta{数字}{数据几万\textasciitilde{}几百万条；ToyLM 学习率 0.5、50 步；推理 0.67 $\rightarrow$ AI 1.00}
\meta{易错}{掩码遮的是损失不是注意力，提示 token 仍在因果掩码内参与前向，并可被回答的 token 看到。}
\phantomsection\addcontentsline{toc}{section}{48\quad RLHF 与 DPO}
\lec{48}{RLHF 与 DPO}
\onel{二选一偏好数据：RLHF 训奖励模型再 PPO，DPO 合成相对参考模型的损失。}
\begin{itemize}[leftmargin=1.35em,itemsep=2.2pt,topsep=1pt,parsep=0pt]
\item \begingroup\sloppy{\fontsize{10}{14.6}\selectfont 偏好数据三元组：提示、两个回答、人的判断；比较比打分校准。\par}\endgroup
\item \begingroup\sloppy{\fontsize{10}{14.6}\selectfont 奖励模型输出单个标量，Bradley--Terry：P(A$\succ$B)=σ(r\_A-r\_B)。\par}\endgroup
\item \begingroup\sloppy{\fontsize{10}{14.6}\selectfont 分差 0/1/3 对应概率 0.5/0.731/0.953。\par}\endgroup
\item \begingroup\sloppy{\fontsize{10}{14.6}\selectfont PPO 目标：奖励 - $\beta$$\times$KL(参考模型)；要装策略、参考、奖励、价值四个模型。\par}\endgroup
\item \begingroup\sloppy{\fontsize{10}{14.6}\selectfont 隐式奖励是相对参考模型的对数概率提升乘 $\beta$，再套 σ 损失。\par}\endgroup
\item \begingroup\sloppy{\fontsize{10}{14.6}\selectfont ToyLM 30 步把苹果概率从 0.501 推到 0.981，两个回答被分开。\par}\endgroup
\end{itemize}
\meta{公式}{$P(A\succ B)=\sigma(r_{A}-r_{B})$；DPO 损失 $=-\ln\sigma\bigl(\beta[(\ln\pi(y_{w})-\ln\pi_{\text{ref}}(y_{w}))-(\ln\pi(y_{l})-\ln\pi_{\text{ref}}(y_{l}))]\bigr)$}
\meta{数字}{$\beta$=0.5；30 步 0.501$\rightarrow$0.884$\rightarrow$0.981；损失 0.693$\rightarrow$0.309$\rightarrow$0.129；分差 1$\rightarrow$0.731、3$\rightarrow$0.953；RLHF 四个模型 vs DPO 两个模型}
\meta{易错}{DPO 不是强化学习但目标相同，它用离线偏好数据直接训练，训练中不生成新回答，省掉奖励模型。}
\phantomsection\addcontentsline{toc}{section}{49\quad 推理模型与强化学习}
\lec{49}{推理模型与强化学习}
\onel{奖励来自可自动验证的结果，用 RL 让模型自己学会先想再答。}
\begin{itemize}[leftmargin=1.35em,itemsep=2.2pt,topsep=1pt,parsep=0pt]
\item \begingroup\sloppy{\fontsize{10}{14.6}\selectfont 思考 token 每多一个就多一次前向，还能看到前面写下的中间结果。\par}\endgroup
\item \begingroup\sloppy{\fontsize{10}{14.6}\selectfont 多次采样投票：单次 0.6 的正确率，5 次升到 0.771，15 次到 0.981。\par}\endgroup
\item \begingroup\sloppy{\fontsize{10}{14.6}\selectfont RLVR 只看结果：答案是否一致、代码能否过测试用例，不管过程。\par}\endgroup
\item \begingroup\sloppy{\fontsize{10}{14.6}\selectfont GRPO 每题采样一组回答，优势 =（奖励 - 组内均值）/ 组内标准差。\par}\endgroup
\item \begingroup\sloppy{\fontsize{10}{14.6}\selectfont 8 个答对 3 个：均值 0.375、标准差 0.484，对的 +1.291。\par}\endgroup
\item \begingroup\sloppy{\fontsize{10}{14.6}\selectfont 先想 8000 再答 500 个 token，生成量 17 倍。\par}\endgroup
\end{itemize}
\meta{公式}{$\text{优势}=\frac{\text{奖励}-\text{组内平均}}{\text{组内标准差}}$；$\text{pass@k}=1-\frac{C(n-c,\,k)}{C(n,\,k)}$}
\meta{数字}{8 采样答对 3：均值 0.375、标准差 0.484，优势 +1.291 / -0.775；多数投票 0.6 $\rightarrow$ 0.771（5 次）$\rightarrow$ 0.981（15 次）；KV Cache 62.5 $\rightarrow$ 1062 MiB}
\meta{易错}{「推理模型」的推理是 reasoning 不是 inference；写出的思考过程不保证忠实反映内部计算；全对或全错的题优势全为 0，不给信号。}
\phantomsection\addcontentsline{toc}{section}{50\quad LoRA 参数高效微调}
\lec{50}{LoRA 参数高效微调}
\onel{冻结 W，只训低秩旁路 A$\cdot$B，可训练参数约 0.52\%。}
\begin{itemize}[leftmargin=1.35em,itemsep=2.2pt,topsep=1pt,parsep=0pt]
\item \begingroup\sloppy{\fontsize{10}{14.6}\selectfont 冻结 W，只训旁路 ΔW = A$\cdot$B，A 为 d\_in$\times$r、B 为 r$\times$d\_out。\par}\endgroup
\item \begingroup\sloppy{\fontsize{10}{14.6}\selectfont r 远小于 d，常用 8、16、64；B 初始为 0，起步时模型与原来一致。\par}\endgroup
\item \begingroup\sloppy{\fontsize{10}{14.6}\selectfont 全模型 4194 万个 LoRA 参数占 0.52\%，r=64 时是 2.09\%。\par}\endgroup
\item \begingroup\sloppy{\fontsize{10}{14.6}\selectfont 显存：全量 128.5 GB；LoRA 为冻结权重 16.1 GB 加 0.67 GB。\par}\endgroup
\item \begingroup\sloppy{\fontsize{10}{14.6}\selectfont QLoRA 把冻结权重压到 4 位约 4 GB；合并 A$\cdot$B 后推理零开销。\par}\endgroup
\item \begingroup\sloppy{\fontsize{10}{14.6}\selectfont 不合并则一个基座挂多个几十 MB 适配器，S-LoRA、Punica 处理混批。\par}\endgroup
\end{itemize}
\meta{公式}{$\Delta W=A\cdot B$，$A\colon d_{\text{in}}\times r$，$B\colon r\times d_{\text{out}}$；$y=xW+(xA)B$；参数 $r\times(d_{\text{in}}+d_{\text{out}})$}
\meta{数字}{r=8 2100 万 / 0.26\%；r=16 4194 万 / 0.52\%；r=64 1.68 亿 / 2.09\%；q\_proj r=16：1678 万 $\rightarrow$ 13.1 万；128.5 GB $\rightarrow$ 16.1 + 0.67 GB；QLoRA 4.0 GB}
\meta{易错}{合并 A$\cdot$B 后推理计算量不变，省的是训练显存与存储；秩越大越易过拟合，也学不进大量新知识。}
\phantomsection\addcontentsline{toc}{section}{51\quad 模型评测}
\lec{51}{模型评测}
\onel{困惑度、基准测试、开放式对战三类尺子各管一段；评测本身是大批量推理。}
\begin{itemize}[leftmargin=1.35em,itemsep=2.2pt,topsep=1pt,parsep=0pt]
\item \begingroup\sloppy{\fontsize{10}{14.6}\selectfont 困惑度即 e\textasciicircum{}loss，只衡量预测下一个 token，跨分词器不能直接比。\par}\endgroup
\item \begingroup\sloppy{\fontsize{10}{14.6}\selectfont 选择题（MMLU，57 学科）比各选项概率，不生成 token，只需一次 prefill。\par}\endgroup
\item \begingroup\sloppy{\fontsize{10}{14.6}\selectfont 答案匹配（GSM8K）提取最终答案比对；代码测试（HumanEval）报 pass@k。\par}\endgroup
\item \begingroup\sloppy{\fontsize{10}{14.6}\selectfont 数据污染：题库公开加网上爬的预训练数据，模型背答案，分数虚高。\par}\endgroup
\item \begingroup\sloppy{\fontsize{10}{14.6}\selectfont Elo 按 R $\leftarrow$ R + K$\times$(实际 - 预期) 更新，K=32，分差定预期胜率。\par}\endgroup
\item \begingroup\sloppy{\fontsize{10}{14.6}\selectfont A 赢只加 7.7 分，爆冷输了 A 掉到 1175.7、B 升到 1024.3。\par}\endgroup
\end{itemize}
\meta{公式}{$P(\text{A 赢})=\frac{1}{1+10^{(R_{B}-R_{A})/400}}$；$R\leftarrow R+K\times(\text{实际}-\text{预期})$}
\meta{数字}{MMLU 57 个学科；ToyLM 4/5 = 80\%；K = 32；1200 vs 1000 $\rightarrow$ 0.760；赢加 7.7 分；1175.7 / 1024.3}
\meta{易错}{答案歧义、few-shot、答案提取规则都能让同一模型差几个百分点；基准分高不等于好用，Elo 只在同一竞技场内可比。}
\chapter{从零实现一个GPT模型}
{\footnotesize\sffamily\color{grayt}收录第 52--59 讲}\par\vspace{-4pt}
\phantomsection\addcontentsline{toc}{section}{52\quad g0 $\cdot$ 语料与分词器}
\lec{52}{g0 $\cdot$ 语料与分词器}
\onel{讲义清洗成 499,220 字符，训一个从字符出发、可退回字节的 BPE。}
\begin{itemize}[leftmargin=1.35em,itemsep=2.2pt,topsep=1pt,parsep=0pt]
\item \begingroup\sloppy{\fontsize{10}{14.6}\selectfont prepare.py 依次删元数据、代码块、公式、组件、HTML 标签，组件要整块先删。\par}\endgroup
\item \begingroup\sloppy{\fontsize{10}{14.6}\selectfont 不用 GPT-2 的 50,257 词表，因为两张表就是 2570 万参数。\par}\endgroup
\item \begingroup\sloppy{\fontsize{10}{14.6}\selectfont 编号 0--255 是字节、256 是分隔符、257 起是字符、再往后是合并。\par}\endgroup
\item \begingroup\sloppy{\fontsize{10}{14.6}\selectfont 编码前先切块，数字一位一切，2048 永远是 2、0、4、8 四个 token。\par}\endgroup
\item \begingroup\sloppy{\fontsize{10}{14.6}\selectfont 合并按学到的先后（编号最小）挑，不按块内频次；每块结果用 cache。\par}\endgroup
\item \begingroup\sloppy{\fontsize{10}{14.6}\selectfont 4096 = 256 字节 + 1356 字符 + 2483 合并 + 1 分隔符。\par}\endgroup
\end{itemize}
\meta{数字}{499,220 字符（约 1.1 MB）；4096 = 256+1+1356+2483；训练 22.3 s；41 字符 $\rightarrow$ 23 token；1.51 字符/token；上下文 256 token 约 390 字符}
\meta{易错}{训练集与验证集必须在分词之前切开，且只用前 90\% 训分词器，否则验证 loss 偏乐观。}
\phantomsection\addcontentsline{toc}{section}{53\quad g1 $\cdot$ 数据加载器与嵌入}
\lec{53}{g1 $\cdot$ 数据加载器与嵌入}
\onel{长序列按 stride 开窗口，x 前 T 个、y 错一位；两次查表相加。}
\begin{itemize}[leftmargin=1.35em,itemsep=2.2pt,topsep=1pt,parsep=0pt]
\item \begingroup\sloppy{\fontsize{10}{14.6}\selectfont 一个样本含 T 道题：位置 t 看 x[0..t] 猜 y[t] = x[t+1]。\par}\endgroup
\item \begingroup\sloppy{\fontsize{10}{14.6}\selectfont stride=T 时窗口首尾相接、token 不重叠，切出 37,809 个样本。\par}\endgroup
\item \begingroup\sloppy{\fontsize{10}{14.6}\selectfont stride < T 样本重叠、易背；stride > T 有 token 用不上。\par}\endgroup
\item \begingroup\sloppy{\fontsize{10}{14.6}\selectfont 训练用 stride=T/2=128，验证集用 stride=T 不重叠。\par}\endgroup
\item \begingroup\sloppy{\fontsize{10}{14.6}\selectfont 只存窗口起点 range(0, len-T, stride)，省下内存。\par}\endgroup
\item \begingroup\sloppy{\fontsize{10}{14.6}\selectfont h = E\_tok[x] + E\_pos[t]，两张表共 1,114,112 个参数。\par}\endgroup
\end{itemize}
\meta{公式}{$h[b,t]=E_{\text{tok}}[x[b,t]]+E_{\text{pos}}[t]$；起点 $\in\operatorname{range}(0,\ \text{len}-T,\ \text{stride})$；$y[t]=x[t+1]$}
\meta{数字}{训练集 302,475 / 验证集 33,048 token；stride 8/4/1 $\rightarrow$ 37,809 / 75,617 / 302,467 个样本；位置嵌入 256$\times$256；token 嵌入占 94\%}
\meta{易错}{y 是 x 往后错一位，写反了配上因果掩码 loss 降不下去；窗口起点写成 len-T+1 会让最后一个样本少一个 token。}
\phantomsection\addcontentsline{toc}{section}{54\quad g2 $\cdot$ 因果多头注意力}
\lec{54}{g2 $\cdot$ 因果多头注意力}
\onel{四个 d$\times$d 投影，拆头成 [B,H,T,h\_d]，批量矩阵乘一次算完所有头。}
\begin{itemize}[leftmargin=1.35em,itemsep=2.2pt,topsep=1pt,parsep=0pt]
\item \begingroup\sloppy{\fontsize{10}{14.6}\selectfont 四个 d$\times$d 线性层，所有头的投影并排成一个矩阵，不用 for 循环。\par}\endgroup
\item \begingroup\sloppy{\fontsize{10}{14.6}\selectfont 先 view 成 [B,T,H,h\_d] 再转置，直接 view 会混位置。\par}\endgroup
\item \begingroup\sloppy{\fontsize{10}{14.6}\selectfont 打分 [B,H,T,T] 除以 $\surd$h\_d，j > t 填 -$\infty$ 后 softmax。\par}\endgroup
\item \begingroup\sloppy{\fontsize{10}{14.6}\selectfont 参数 4d$^2$+d = 262,400，与头数无关；掩码是 buffer 不算参数。\par}\endgroup
\item \begingroup\sloppy{\fontsize{10}{14.6}\selectfont 和 F.scaled\_dot\_product\_attention 最大差 3$\times$10$^-$⁸。\par}\endgroup
\item \begingroup\sloppy{\fontsize{10}{14.6}\selectfont 注意力权重 419 万个数、FP32 16 MiB，与 T$^2$ 成正比。\par}\endgroup
\end{itemize}
\meta{公式}{$\text{Attn}=\operatorname{softmax}\left(\frac{QK^{\top}}{\sqrt{h_{d}}}+\text{mask}\right)V$，$j>t$ 处 $\text{mask}=-\infty$；$h_{d}=\frac{d}{H}$}
\meta{数字}{d=256、H=4、h\_d=64；参数 262,400；与 PyTorch 差 3$\times$10$^-$⁸；权重 [16,4,256,256] 419 万个数、16 MiB；改第 3 个 token 只有位置 3 变}
\meta{易错}{缩放用 $\surd$h\_d 不是 $\surd$d；掩码方向写成 tril 等于偷看答案，训练 loss 掉得飞快、生成却全崩。}
\phantomsection\addcontentsline{toc}{section}{55\quad g3 $\cdot$ 组装GPT}
\lec{55}{g3 $\cdot$ 组装GPT}
\onel{按 Pre-Norm 补齐 LN、FFN、残差，组装 530 万参数的 GPT。}
\begin{itemize}[leftmargin=1.35em,itemsep=2.2pt,topsep=1pt,parsep=0pt]
\item \begingroup\sloppy{\fontsize{10}{14.6}\selectfont Block 是 Pre-Norm：先 LN 再进子层，再加残差。\par}\endgroup
\item \begingroup\sloppy{\fontsize{10}{14.6}\selectfont FFN 是 d $\rightarrow$ 4d $\rightarrow$ GELU $\rightarrow$ d：注意力搬信息，FFN 在每个位置上加工。\par}\endgroup
\item \begingroup\sloppy{\fontsize{10}{14.6}\selectfont 每层注意力 262,400、FFN 525,568、两个 LN 1,024。\par}\endgroup
\item \begingroup\sloppy{\fontsize{10}{14.6}\selectfont 5,319,168 个参数里，两张词表大小的表占 39\%。\par}\endgroup
\item \begingroup\sloppy{\fontsize{10}{14.6}\selectfont GPT-2 small 参数 163,037,184，共用输出层后 124M。\par}\endgroup
\item \begingroup\sloppy{\fontsize{10}{14.6}\selectfont 每步把整段序列送进模型，取最后位置 logits 做 argmax。\par}\endgroup
\end{itemize}
\meta{公式}{$x\leftarrow x+\operatorname{Attn}(\operatorname{LN}(x))$；$x\leftarrow x+\operatorname{FFN}(\operatorname{LN}(x))$；每层 $\approx12d^{2}$；总参数 $\approx L\cdot12d^{2}+2Vd$}
\meta{数字}{5,319,168 参数（FP32 20.3 MiB）；一层 788,992；GPT-2 small 163,037,184 $\rightarrow$ 共用 124,439,808；输出层 d$\times$V 无偏置、此处未共用}
\meta{易错}{Post-Norm 与 Pre-Norm 混用都能跑，但装 GPT-2 权重时结构必须一致，否则输出完全不同；生成不截断到 T\_max 会越界。}
\phantomsection\addcontentsline{toc}{section}{56\quad g4 $\cdot$ 预训练}
\lec{56}{g4 $\cdot$ 预训练}
\onel{loss = -ln p，起点约 ln V；数据太少时验证 loss 先降后升。}
\begin{itemize}[leftmargin=1.35em,itemsep=2.2pt,topsep=1pt,parsep=0pt]
\item \begingroup\sloppy{\fontsize{10}{14.6}\selectfont loss = -ln p，起点约 ln V = 8.318；困惑度 = e\textasciicircum{}loss。\par}\endgroup
\item \begingroup\sloppy{\fontsize{10}{14.6}\selectfont 3000 步共 12.3 M token，等于把训练集过了 41 遍。\par}\endgroup
\item \begingroup\sloppy{\fontsize{10}{14.6}\selectfont Chinchilla 每参数约 20 token：需 106 M，只有 0.30 M。\par}\endgroup
\item \begingroup\sloppy{\fontsize{10}{14.6}\selectfont 训练 loss 8.45$\rightarrow$0.98，验证最低 4.843，末了 5.857。\par}\endgroup
\item \begingroup\sloppy{\fontsize{10}{14.6}\selectfont 6N$\times$token = 3.92e14 FLOPs，平均 0.69 TFLOPS。\par}\endgroup
\item \begingroup\sloppy{\fontsize{10}{14.6}\selectfont 同一 batch 反复训能否把 loss 压到近 0，是查训练代码最有效的一招。\par}\endgroup
\end{itemize}
\meta{公式}{$\text{loss}=-\ln p$；困惑度 $=e^{\text{loss}}$；训练算力 $\approx6N\times\text{token 数}$}
\meta{数字}{ln 4096 = 8.318；8.45 $\rightarrow$ 0.98；验证最低 4.843（第 900 步）$\rightarrow$ 5.857；41 遍；差 352 倍；569 s、0.69 TFLOPS、利用率 19\%}
\meta{易错}{验证 loss 回升不是 bug 而是过拟合；要翻转结论只能加数据或缩小模型，dropout 只是推迟。}
\phantomsection\addcontentsline{toc}{section}{57\quad g5 $\cdot$ 解码与存档}
\lec{57}{g5 $\cdot$ 解码与存档}
\onel{存档存配置+参数+优化器状态并留验证最优；解码加温度、Top-k、Top-p。}
\begin{itemize}[leftmargin=1.35em,itemsep=2.2pt,topsep=1pt,parsep=0pt]
\item \begingroup\sloppy{\fontsize{10}{14.6}\selectfont 存档含配置、参数、优化器状态、步数与日志，带配置才自描述。\par}\endgroup
\item \begingroup\sloppy{\fontsize{10}{14.6}\selectfont AdamW 每参数两份滑动平均，优化器状态是参数的两倍大。\par}\endgroup
\item \begingroup\sloppy{\fontsize{10}{14.6}\selectfont 验证 loss 创新低就存 gpt\_best.pt，等同早停：留最好那份。\par}\endgroup
\item \begingroup\sloppy{\fontsize{10}{14.6}\selectfont 数据只有计算最优的 1/350，第 900 步已看 12 遍，正则化无用。\par}\endgroup
\item \begingroup\sloppy{\fontsize{10}{14.6}\selectfont 0.2 的 dropout 让末步验证 loss 5.25，最低点仍 4.876。\par}\endgroup
\item \begingroup\sloppy{\fontsize{10}{14.6}\selectfont 温度调尖平（T$\rightarrow$0 即贪心），Top-k 留 k 个，Top-p 累到 p。\par}\endgroup
\end{itemize}
\meta{公式}{$l\leftarrow l/T$；Top-$k$ 取前 $k$ 大；Top-$p$ 按“前面的累计概率”截断；其余置 $-\infty$ 后 softmax 采样}
\meta{数字}{最低验证 loss 4.876（第 900 步）；存档 20.6$\rightarrow$61.2 MiB；T=1 凑 90\% 要 573 token}
\meta{易错}{Top-p 要按“前面的累计概率”判断截断，用“包含自身”会把凑够 p 的那个也删掉，分布极尖时全删、softmax 出 NaN。}
\phantomsection\addcontentsline{toc}{section}{58\quad g6 $\cdot$ 装进 GPT-2 的权重}
\lec{58}{g6 $\cdot$ 装进 GPT-2 的权重}
\onel{GPT-2 权重改名、切 qkv、转置后装进同一个 GPT 类。}
\begin{itemize}[leftmargin=1.35em,itemsep=2.2pt,topsep=1pt,parsep=0pt]
\item \begingroup\sloppy{\fontsize{10}{14.6}\selectfont GPT-2 small 有 149 个张量、1.24 亿参数，d=768、12 层。\par}\endgroup
\item \begingroup\sloppy{\fontsize{10}{14.6}\selectfont c\_attn 为 [768, 2304]，按列切出 q、k、v 三份再转置。\par}\endgroup
\item \begingroup\sloppy{\fontsize{10}{14.6}\selectfont Conv1D 存 [输入, 输出]，nn.Linear 存 [输出, 输入]，互为转置。\par}\endgroup
\item \begingroup\sloppy{\fontsize{10}{14.6}\selectfont wte 复制进输出层；装完 1.63 亿参数，多出 3860 万。\par}\endgroup
\item \begingroup\sloppy{\fontsize{10}{14.6}\selectfont 随机初始化即可验证：最大误差 5$\times$10$^-$⁶，argmax 全一致。\par}\endgroup
\item \begingroup\sloppy{\fontsize{10}{14.6}\selectfont 方阵忘转置：形状对、不报错，误差 3.29、argmax 全错。\par}\endgroup
\end{itemize}
\meta{公式}{$c_{\text{attn}}[768,2304]\to$ 按列切 $q/k/v$ 各 $[768,768]$；$W_{\text{ours}}=W_{\text{conv1d}}^{\top}$；$\text{out\_head}\leftarrow w_{\text{te}}$}
\meta{数字}{149 张量 / 1.24 亿参数；装后 1.63 亿（+3860 万）；误差 5$\times$10$^-$⁶；忘转置 3.29}
\meta{易错}{c\_attn 必须按列（输出维）切，split(dim=0) 也能切出三个方阵但结果是错的；分词器也得换成 GPT-2 自己的。}
\phantomsection\addcontentsline{toc}{section}{59\quad g7 $\cdot$ 指令微调}
\lec{59}{g7 $\cdot$ 指令微调}
\onel{SFT 只把指令位置的 y 设 -100，只在回答与结束符上算 loss。}
\begin{itemize}[leftmargin=1.35em,itemsep=2.2pt,topsep=1pt,parsep=0pt]
\item \begingroup\sloppy{\fontsize{10}{14.6}\selectfont 数据取自课程目录两栏，每讲两题，112 讲出 224 道。\par}\endgroup
\item \begingroup\sloppy{\fontsize{10}{14.6}\selectfont 每 10 讲留 1 讲：200 条训练、24 条测试，同讲题不跨界。\par}\endgroup
\item \begingroup\sloppy{\fontsize{10}{14.6}\selectfont 一条样本 34 token：前 28 个模板与指令，后 6 个回答加结束符。\par}\endgroup
\item \begingroup\sloppy{\fontsize{10}{14.6}\selectfont 指令与补齐位置的 y 设 -100，cross\_entropy 会跳过。\par}\endgroup
\item \begingroup\sloppy{\fontsize{10}{14.6}\selectfont 从第 900 步存档起训 20 轮，loss 5.23$\rightarrow$0.11。\par}\endgroup
\item \begingroup\sloppy{\fontsize{10}{14.6}\selectfont 训练集 198/200 全对、测试集 0/24，但回答格式全对。\par}\endgroup
\end{itemize}
\meta{公式}{$y[i]=-100$（指令与补齐位置）；$L=\operatorname{CE}(\text{logits}[\text{回答位置}],\ y[\text{回答位置}])$}
\meta{数字}{112 讲$\rightarrow$224 题；200 训练 / 24 测试；34 token（28+6）；loss 5.227$\rightarrow$0.112}
\meta{易错}{屏蔽差一位不报错，只是 loss 算在错的位置上；回答后不接结束符，模型学不会停。}
\chapter{硬件与性能基础}
{\footnotesize\sffamily\color{grayt}收录第 60--67 讲}\par\vspace{-4pt}
\phantomsection\addcontentsline{toc}{section}{60\quad GPU 架构}
\lec{60}{GPU 架构}
\onel{算力来自 132 个 SM 各 4 个 Tensor Core 的矩阵乘加。}
\begin{itemize}[leftmargin=1.35em,itemsep=2.2pt,topsep=1pt,parsep=0pt]
\item \begingroup\sloppy{\fontsize{10}{14.6}\selectfont CPU 少数强核擅复杂逻辑；GPU 上千简单核擅大量同种独立运算。\par}\endgroup
\item \begingroup\sloppy{\fontsize{10}{14.6}\selectfont H100 每个 SM 有 4 个 Tensor Core、128 个 FP32 核。\par}\endgroup
\item \begingroup\sloppy{\fontsize{10}{14.6}\selectfont CUDA 核峰值 62 TFLOPS，只有 Tensor Core 的 6\%。\par}\endgroup
\item \begingroup\sloppy{\fontsize{10}{14.6}\selectfont 64 核 CPU 约 10 TFLOPS；4096$^3$ 矩阵乘 13 ms。\par}\endgroup
\item \begingroup\sloppy{\fontsize{10}{14.6}\selectfont 全卡约 27 万线程驻留，用多线程掩盖数百周期访存延迟。\par}\endgroup
\item \begingroup\sloppy{\fontsize{10}{14.6}\selectfont 要达峰值需数据供得上、SM 都有活，实测 60\%\textasciitilde{}80\% 已很好。\par}\endgroup
\end{itemize}
\meta{公式}{$P=\text{SM}\times\frac{\text{TC}}{\text{SM}}\times\frac{\text{FMA}}{\text{TC}}\times2\times f=132\times4\times512\times2\times1.83\ \text{GHz}\approx989\ \text{TFLOPS}$}
\meta{数字}{H100：132 SM、16,896 个 CUDA 核、3.35 TB/s、每 SM 2048 线程；4096$^3$=137 GFLOP}
\meta{易错}{规格表的“CUDA 核数”不代表算力，大模型走 Tensor Core；比峰值要统一精度与稠密/稀疏（稀疏是稠密两倍）。}
\phantomsection\addcontentsline{toc}{section}{61\quad 存储层级}
\lec{61}{存储层级}
\onel{H100 从 HBM 每读 1 字节能做 295 次运算，所以要在片上复用数据。}
\begin{itemize}[leftmargin=1.35em,itemsep=2.2pt,topsep=1pt,parsep=0pt]
\item \begingroup\sloppy{\fontsize{10}{14.6}\selectfont 寄存器、共享内存各约 35 MB，L2 约 50 MB，HBM 80 GB。\par}\endgroup
\item \begingroup\sloppy{\fontsize{10}{14.6}\selectfont 989 TFLOPS ÷ 3.35 TB/s $\approx$ 295 次运算/字节。\par}\endgroup
\item \begingroup\sloppy{\fontsize{10}{14.6}\selectfont 1 GB 从 HBM、PCIe 读分别要 0.30、15.6 ms。\par}\endgroup
\item \begingroup\sloppy{\fontsize{10}{14.6}\selectfont n$\times$n 矩阵乘每字节约 n/3 次：256 喂不饱，4096 喂得饱。\par}\endgroup
\item \begingroup\sloppy{\fontsize{10}{14.6}\selectfont 逐元素每字节 0.25 次，读写 1 GB 要 0.60 ms，算 0.5 $\mu$s。\par}\endgroup
\item \begingroup\sloppy{\fontsize{10}{14.6}\selectfont SRAM 单元 6 个晶体管、面积数倍于 DRAM，片上只能放几十 MB。\par}\endgroup
\end{itemize}
\meta{公式}{$n\times n$ BF16 矩阵乘法：$2n^{3}$ 次运算 $\div6n^{2}$ 字节 $=\frac{n}{3}$ 次运算每字节；喂饱比例 $=\frac{\text{峰值算力}}{\text{带宽}}$}
\meta{数字}{共享内存约 31 TB/s；L2 5.5 TB/s；HBM 3.35 TB/s；A100 150、H100 295、B200 280 次/字节}
\meta{易错}{“显存”指 HBM、“内存”指 CPU 那侧的 DRAM；带宽与延迟是两回事，大模型运算通常卡带宽。}
\phantomsection\addcontentsline{toc}{section}{62\quad Roofline 模型}
\lec{62}{Roofline 模型}
\onel{速度上限 min(P, I$\times$B)；H100 屋脊点 295 FLOPs/字节。}
\begin{itemize}[leftmargin=1.35em,itemsep=2.2pt,topsep=1pt,parsep=0pt]
\item \begingroup\sloppy{\fontsize{10}{14.6}\selectfont 算术强度 I=W÷Q，速度 min(P, I$\times$B)，屋脊点 I*=P÷B。\par}\endgroup
\item \begingroup\sloppy{\fontsize{10}{14.6}\selectfont 时间 max(W÷P, Q÷B)；低于屋脊点卡带宽，高于卡算力。\par}\endgroup
\item \begingroup\sloppy{\fontsize{10}{14.6}\selectfont 线性层强度约等于一起处理的 token 数，1 个只到峰值 0.3\%。\par}\endgroup
\item \begingroup\sloppy{\fontsize{10}{14.6}\selectfont 32 个 token 强度 32$\rightarrow$106；4096 个强度 1792 卡算力。\par}\endgroup
\item \begingroup\sloppy{\fontsize{10}{14.6}\selectfont decode 注意力强度 1\textasciitilde{}4，GQA$\times$4 也只到 13.4 TFLOPS。\par}\endgroup
\item \begingroup\sloppy{\fontsize{10}{14.6}\selectfont A100 屋脊点 153、H100 295、B200 281；FP8 时翻倍到 590。\par}\endgroup
\end{itemize}
\meta{公式}{速度 $=\min(P,\ I\cdot B)$；$I=W/Q$；屋脊点 $I^{*}=P/B$；时间 $=\max(W/P,\ Q/B)$}
\meta{数字}{H100 989 TFLOPS ÷ 3.35 TB/s=295；256 token 强度 237$\rightarrow$794 TFLOPS（80\%）；逐元素 0.25$\rightarrow$0.8}
\meta{易错}{带宽区“利用率不到 1\%”不是 bug，是运算性质决定的；判断实现好坏要看它离自己那条屋顶有多远。}
\phantomsection\addcontentsline{toc}{section}{63\quad 低精度计算}
\lec{63}{低精度计算}
\onel{精度减半则算力翻倍：BF16 989、FP8 1979 TFLOPS。}
\begin{itemize}[leftmargin=1.35em,itemsep=2.2pt,topsep=1pt,parsep=0pt]
\item \begingroup\sloppy{\fontsize{10}{14.6}\selectfont H100 稠密算力 TF32 495、BF16 989、FP8 1979 TFLOPS。\par}\endgroup
\item \begingroup\sloppy{\fontsize{10}{14.6}\selectfont 乘法器面积约正比于尾数位数平方：FP8 是 BF16 的 0.25。\par}\endgroup
\item \begingroup\sloppy{\fontsize{10}{14.6}\selectfont 数字节减半，同带宽与片上存储能放两倍的数，能耗更低。\par}\endgroup
\item \begingroup\sloppy{\fontsize{10}{14.6}\selectfont decode 省带宽：70B 读权重 BF16 5.3 ms、FP8 2.6 ms。\par}\endgroup
\item \begingroup\sloppy{\fontsize{10}{14.6}\selectfont prefill 省算力：8192 token 的 70B，146 ms$\rightarrow$73 ms。\par}\endgroup
\item \begingroup\sloppy{\fontsize{10}{14.6}\selectfont FP8 乘法仍用 FP16/FP32 累加器，否则误差累积很大。\par}\endgroup
\end{itemize}
\meta{数字}{H100 80 GB 放不下 70B 的 BF16 141 GB；FP8 权重 71 GB；H200 141 GB 放得下}
\meta{易错}{存成 FP8 不等于用 FP8 算：读进来转 BF16 只省带宽；FP16 与 BF16 算力相同，但 BF16 范围与 FP32 一样。}
\phantomsection\addcontentsline{toc}{section}{64\quad 互联}
\lec{64}{互联}
\onel{机内 NVLink 450 GB/s、机间 50 GB/s，机内快约 9 倍。}
\begin{itemize}[leftmargin=1.35em,itemsep=2.2pt,topsep=1pt,parsep=0pt]
\item \begingroup\sloppy{\fontsize{10}{14.6}\selectfont PCIe 5.0 x16 每方向约 64 GB/s，比 HBM 慢 50 倍。\par}\endgroup
\item \begingroup\sloppy{\fontsize{10}{14.6}\selectfont NVLink 每卡单向：A100 300、H100 450、B200 900 GB/s。\par}\endgroup
\item \begingroup\sloppy{\fontsize{10}{14.6}\selectfont 机间用 InfiniBand 或 RoCE，均基于 RDMA，50 GB/s 量级。\par}\endgroup
\item \begingroup\sloppy{\fontsize{10}{14.6}\selectfont 路径带宽取最慢一段：同机 1 GB 约 2 ms，跨机约 20 ms。\par}\endgroup
\item \begingroup\sloppy{\fontsize{10}{14.6}\selectfont 小消息几乎全是固定开销 $\alpha$：NVLink 3、PCIe 5、IB 6 $\mu$s。\par}\endgroup
\item \begingroup\sloppy{\fontsize{10}{14.6}\selectfont 整机 NVLink 3600 GB/s、网卡 400 GB/s，差 9 倍。\par}\endgroup
\end{itemize}
\meta{公式}{$t=\alpha+\frac{m}{\beta}$；$\alpha$ 取 NVLink $3\ \mu\text{s}$、PCIe $5\ \mu\text{s}$、IB $6\ \mu\text{s}$、RoCE $8\ \mu\text{s}$；瓶颈带宽 $=\min(\text{路径各段})$}
\meta{数字}{PCIe 4.0 x16 32 GB/s、5.0 x16 64 GB/s；NVLink 4 双向 900 GB/s；IB NDR 400 Gb/s=50 GB/s}
\meta{易错}{“H100 NVLink 900 GB/s”是双向合计，估算一次传输要用单向 450 GB/s。}
\phantomsection\addcontentsline{toc}{section}{65\quad 超节点}
\lec{65}{超节点}
\onel{超节点把 scale-up 域扩到几十上百卡，All-to-All 受益最大}
\begin{itemize}[leftmargin=1.35em,itemsep=2.2pt,topsep=1pt,parsep=0pt]
\item \begingroup\sloppy{\fontsize{10}{14.6}\selectfont scale-up 用 NVLink 直连，H100 每卡 450 GB/s、B200 900 GB/s，可 load/store 对方显存。\par}\endgroup
\item \begingroup\sloppy{\fontsize{10}{14.6}\selectfont scale-out 用 InfiniBand/RoCE，规模不限但每卡约 50 GB/s，差 9 倍且是消息语义。\par}\endgroup
\item \begingroup\sloppy{\fontsize{10}{14.6}\selectfont 域大小受铜缆一两米距离、交换芯片端口数、几十上百千瓦功耗与散热限制。\par}\endgroup
\item \begingroup\sloppy{\fontsize{10}{14.6}\selectfont GB200 NVL72 是 72 卡、每卡双向 1.8 TB/s；CloudMatrix 384 用 384 张昇腾 910C。\par}\endgroup
\item \begingroup\sloppy{\fontsize{10}{14.6}\selectfont 64 卡 1 GB All-Reduce：8.39 $\rightarrow$ 2.26 ms（3.7 倍）；All-to-All：17.79 $\rightarrow$ 1.12 ms（16 倍）。\par}\endgroup
\item \begingroup\sloppy{\fontsize{10}{14.6}\selectfont 对分带宽 1.6 对 28.8 TB/s（18 倍）；能分层的受益有限，All-to-All 受益最大。\par}\endgroup
\end{itemize}
\meta{公式}{分层 All-Reduce：机内 Reduce-Scatter 缩成 $1/8$ $\to$ 机间对 $1/8$ 做 All-Reduce $\to$ 机内 All-Gather}
\meta{数字}{机内 450 GB/s；机间 50 GB/s；NVL72 72 卡、双向 1.8 TB/s；64 卡 All-Reduce 8.39$\rightarrow$2.26 ms；All-to-All 17.79$\rightarrow$1.12 ms；对分 1.6 vs 28.8 TB/s}
\meta{易错}{超节点不是更快的单卡，只改变卡间连接；「一个机柜」不等于「一个 scale-up 域」，看连接而非物理边界。}
\phantomsection\addcontentsline{toc}{section}{66\quad 性能分析}
\lec{66}{性能分析}
\onel{读时间线四步：GPU 空等、按类别分账、找最耗时 kernel、和屋顶比}
\begin{itemize}[leftmargin=1.35em,itemsep=2.2pt,topsep=1pt,parsep=0pt]
\item \begingroup\sloppy{\fontsize{10}{14.6}\selectfont Profiler 粗看、Nsight Systems 看谁在等谁、Nsight Compute 深挖单个 kernel。\par}\endgroup
\item \begingroup\sloppy{\fontsize{10}{14.6}\selectfont Llama-3-8B、TP=4、decode 一步 3.09 ms，GPU 忙 91\%、空等 0.27 ms。\par}\endgroup
\item \begingroup\sloppy{\fontsize{10}{14.6}\selectfont 空等来自小 kernel 间 CPU 发射跟不上，换 CUDA Graph 后一步降到 2.84 ms。\par}\endgroup
\item \begingroup\sloppy{\fontsize{10}{14.6}\selectfont 分账：矩阵乘 1.69 ms（55\%）、逐元素 0.49、通信 0.35、注意力 0.30；大头 gate\_up 0.80 ms。\par}\endgroup
\item \begingroup\sloppy{\fontsize{10}{14.6}\selectfont gate\_up 一次 24.9 $\mu$s 读 58.7 MB，实际带宽 2.36 TB/s，只到屋顶 3.35 的 70\%。\par}\endgroup
\item \begingroup\sloppy{\fontsize{10}{14.6}\selectfont 测量要 synchronize 或 CUDA 事件、先预热、在真实负载下测，并用端到端计时复核。\par}\endgroup
\end{itemize}
\meta{公式}{还能优化多少 $=\frac{\text{屋顶}}{\text{实测}}$（$2.36/3.35\approx70\%$）；只占 $1\%$ 的 kernel 优化到零也只快 $1\%$}
\meta{数字}{一步 3.09 ms；384 个 kernel；GPU 忙 91\%；空等 0.27 ms；CUDA Graph 2.84 ms；gate\_up 24.9 $\mu$s / 58.7 MB / 2.36 TB/s；屋顶 3.35 TB/s}
\meta{易错}{nvidia-smi 的利用率只说明有 kernel 在跑，不代表用上了带宽和算力；模拟时间线不是测量。}
\phantomsection\addcontentsline{toc}{section}{67\quad 理论 vs 实测}
\lec{67}{理论 vs 实测}
\onel{读内存标定带宽、大矩阵乘标定算力，得到真实屋顶再预测 decode 上限}
\begin{itemize}[leftmargin=1.35em,itemsep=2.2pt,topsep=1pt,parsep=0pt]
\item \begingroup\sloppy{\fontsize{10}{14.6}\selectfont DGX Spark 用 GB10，128 GB 统一内存、带宽 273 GB/s、BF16 约 125 TFLOPS（待核实）。\par}\endgroup
\item \begingroup\sloppy{\fontsize{10}{14.6}\selectfont H100 是 989 TFLOPS、3.35 TB/s；带宽差 12 倍、算力差 8 倍，屋脊点 458 对 295。\par}\endgroup
\item \begingroup\sloppy{\fontsize{10}{14.6}\selectfont 读 4 GB 内存标定带宽，实测通常只有标称的 70\%\textasciitilde{}90\%，之后都要乘这个系数。\par}\endgroup
\item \begingroup\sloppy{\fontsize{10}{14.6}\selectfont 4096$^3$ BF16 矩阵乘标定算力，强度约 1365，cuBLAS 通常达峰值的 60\%\textasciitilde{}80\%。\par}\endgroup
\item \begingroup\sloppy{\fontsize{10}{14.6}\selectfont 线性层 1、8、64 个 token 都约 0.43 ms、512 个约 0.50 ms，强度约等于 token 数。\par}\endgroup
\item \begingroup\sloppy{\fontsize{10}{14.6}\selectfont decode 上限$\approx$带宽÷权重：8B BF16 16.1 GB，DGX Spark 17、H100 209 tok/s。\par}\endgroup
\end{itemize}
\meta{公式}{$t=\max\left(\frac{\text{FLOPs}}{\text{峰值算力}},\ \frac{\text{字节}}{\text{带宽}}\right)$；decode 速度上限 $\approx\frac{\text{带宽}}{\text{权重字节数}}$}
\meta{数字}{273 GB/s；989 TFLOPS / 3.35 TB/s；屋脊点 458 vs 295；读 4 GB 14.65 vs 1.19 ms；4096$^3$ 1.10 vs 0.14 ms；4 bit 8B 60 tok/s；70B 4 bit 39.7 GB 约 7 tok/s}
\meta{易错}{理论值是上限不是应该值，超过 100\% 多半是没同步或数据还在缓存里；统一内存不是 HBM。}
\chapter{推理基础}
{\footnotesize\sffamily\color{grayt}收录第 68--73 讲}\par\vspace{-4pt}
\phantomsection\addcontentsline{toc}{section}{68\quad Prefill 与 Decode}
\lec{68}{Prefill 与 Decode}
\onel{prefill 算力受限、decode 带宽受限，两种瓶颈决定两套优化}
\begin{itemize}[leftmargin=1.35em,itemsep=2.2pt,topsep=1pt,parsep=0pt]
\item \begingroup\sloppy{\fontsize{10}{14.6}\selectfont 耗时 $\approx$ max(计算/算力, 字节/带宽)；H100 的 989 TFLOPS ÷ 3.35 TB/s $\approx$ 295 FLOPs/字节。\par}\endgroup
\item \begingroup\sloppy{\fontsize{10}{14.6}\selectfont 强度大于 295 计算受限，小于 295 带宽受限，阶段瓶颈就看这个比值。\par}\endgroup
\item \begingroup\sloppy{\fontsize{10}{14.6}\selectfont prefill 1000 token：计算 15.44 ms、读写 4.85 ms，强度 927，卡在算力。\par}\endgroup
\item \begingroup\sloppy{\fontsize{10}{14.6}\selectfont decode 1 token：计算 0.016 ms、读写 4.83 ms，强度 0.93，算力只用 0.3\%。\par}\endgroup
\item \begingroup\sloppy{\fontsize{10}{14.6}\selectfont 提示 1000 生成 200：15.4 + 967 $\approx$ 983 ms，decode 占 98\%。\par}\endgroup
\item \begingroup\sloppy{\fontsize{10}{14.6}\selectfont prefill 减计算（FlashAttention、低精度），decode 减读取（批处理、量化、GQA/MLA）。\par}\endgroup
\end{itemize}
\meta{公式}{耗时 $\approx\max\left(\frac{\text{计算量}}{\text{算力}},\ \frac{\text{读写字节}}{\text{带宽}}\right)$；算术强度 $=\frac{\text{FLOPs}}{\text{字节}}$，H100 分界线 $295$}
\meta{数字}{989 TFLOPS；3.35 TB/s；295 FLOPs/字节；prefill 15.44 ms（强度 927）；decode 4.83 ms（强度 0.93，算力 0.3\%）；decode 占 98\%}
\meta{易错}{decode 慢不是因为算得多（计算量只有 prefill 的千分之一），而是每步都要把 16 GB 权重整个读一遍。}
\phantomsection\addcontentsline{toc}{section}{69\quad KV Cache}
\lec{69}{KV Cache}
\onel{只存每层的 k、v，decode 只算新 token，计算量降约 750 倍}
\begin{itemize}[leftmargin=1.35em,itemsep=2.2pt,topsep=1pt,parsep=0pt]
\item \begingroup\sloppy{\fontsize{10}{14.6}\selectfont 因果掩码让历史 token 输出不变，只需缓存每层的 k、v，不缓存 q。\par}\endgroup
\item \begingroup\sloppy{\fontsize{10}{14.6}\selectfont decode 每步只算新 token 的 q、k、v，k、v 入缓存，再用 q 对全部缓存打分加权。\par}\endgroup
\item \begingroup\sloppy{\fontsize{10}{14.6}\selectfont prefill 一次写入整段提示各层所有 token 的 k、v，之后每步只写一组。\par}\endgroup
\item \begingroup\sloppy{\fontsize{10}{14.6}\selectfont 各 1000 token：无缓存 149.95 万位置 22.5 PFLOPs，有缓存 1999 位置 30 TFLOPs。\par}\endgroup
\item \begingroup\sloppy{\fontsize{10}{14.6}\selectfont 代价：每 token 128 KiB，请求结束 250 MiB，随长度与并发线性涨，16$\times$8192 就 16 GiB。\par}\endgroup
\item \begingroup\sloppy{\fontsize{10}{14.6}\selectfont 一层形状 [请求数, KV 头数, 序列长, 每头维度]；decode 用 FlashDecoding 切段并行读缓存。\par}\endgroup
\end{itemize}
\meta{公式}{每 token KV 字节 $=2\times\text{层数}\times\text{KV 头数}\times\text{每头维度}\times\text{字节数}$；Llama-3-8B：$2\times32\times8\times128\times2=128$ KiB}
\meta{数字}{TinyLM 位置 7 对 28；无缓存 1,499,500 位置 / 22.5 PFLOPs；有缓存 1,999 位置 / 30 TFLOPs；约 750 倍；每 token 128 KiB；250 MiB}
\meta{易错}{缓存只在一次请求内有效，请求结束即释放；开头相同的下一个请求默认仍要重新 prefill，除非用前缀缓存。}
\phantomsection\addcontentsline{toc}{section}{70\quad TTFT、TPOT 与吞吐}
\lec{70}{TTFT、TPOT 与吞吐}
\onel{TTFT 看 prefill、TPOT 看 decode；吞吐靠并发，上限看显存}
\begin{itemize}[leftmargin=1.35em,itemsep=2.2pt,topsep=1pt,parsep=0pt]
\item \begingroup\sloppy{\fontsize{10}{14.6}\selectfont 总耗时 $\approx$ TTFT + TPOT $\times$ (生成数 - 1)；TPOT 几十毫秒内用户就觉流畅。\par}\endgroup
\item \begingroup\sloppy{\fontsize{10}{14.6}\selectfont 单请求：TTFT 15.5 ms、TPOT 4.85 ms 即 206 token/s，算力只用 0.3\%。\par}\endgroup
\item \begingroup\sloppy{\fontsize{10}{14.6}\selectfont 批处理共享一次权重读取；KV 不能共享，一步读写 $\approx$ 16 GB + b $\times$ 192 MiB。\par}\endgroup
\item \begingroup\sloppy{\fontsize{10}{14.6}\selectfont 并发 1$\rightarrow$8：TPOT 4.85$\rightarrow$5.27、吞吐 206$\rightarrow$1,517；并发 128：TPOT 12.49、吞吐 10,251。\par}\endgroup
\item \begingroup\sloppy{\fontsize{10}{14.6}\selectfont 显存是硬上限：每请求 2048 token 占 256 MiB KV，80 GB 卡最多约 238 个。\par}\endgroup
\item \begingroup\sloppy{\fontsize{10}{14.6}\selectfont 并发 256 强度也只约 60，长上下文 decode 始终带宽受限。\par}\endgroup
\end{itemize}
\meta{公式}{总耗时 $\approx\text{TTFT}+\text{TPOT}\times(N_{\text{out}}-1)$；总吞吐 $=b/\text{TPOT}$；一步读写 $\approx16\ \text{GB}+b\times\text{每请求 KV}$}
\meta{数字}{TTFT 15.5 ms；TPOT 4.85$\rightarrow$12.49 ms；吞吐 206$\rightarrow$10,251 token/s；上限约 238 个请求（TPOT 19.1 ms、12,463 token/s）；并发 256 强度约 60}
\meta{易错}{吞吐高不等于每个用户都快：并发 128 时总吞吐是 50 倍，单请求却从 206 降到 80 token/s。}
\phantomsection\addcontentsline{toc}{section}{71\quad 显存账本}
\lec{71}{显存账本}
\onel{权重、KV Cache、激活、开销四项算清，剩下的显存全给 KV Cache}
\begin{itemize}[leftmargin=1.35em,itemsep=2.2pt,topsep=1pt,parsep=0pt]
\item \begingroup\sloppy{\fontsize{10}{14.6}\selectfont 80 GB 卡按 90\% 只有 67.1 GiB；8B 的 BF16 权重 16.06 GB 即 15.0 GiB。\par}\endgroup
\item \begingroup\sloppy{\fontsize{10}{14.6}\selectfont 每 token KV = 2$\times$层数$\times$KV 头数$\times$每头维度$\times$字节数，8B 的 BF16 是 128 KiB。\par}\endgroup
\item \begingroup\sloppy{\fontsize{10}{14.6}\selectfont 激活按一次处理的 token 数预留，prefill 8192 token 约 0.91 GiB；其他约 1.9 GiB。\par}\endgroup
\item \begingroup\sloppy{\fontsize{10}{14.6}\selectfont 8B 单卡剩 49.3 GiB 给 KV，约 40.4 万 token，即 99 个 4096 token 请求。\par}\endgroup
\item \begingroup\sloppy{\fontsize{10}{14.6}\selectfont 权重与 KV 都换 FP8：权重 7.5 GiB、每 token 64 KiB，能放 227 个请求。\par}\endgroup
\item \begingroup\sloppy{\fontsize{10}{14.6}\selectfont 70B 的 BF16 131.4 GiB、FP8 65.7 GiB 单卡与两卡都放不下 KV；四卡能放 102 个。\par}\endgroup
\end{itemize}
\meta{公式}{显存 $=\text{权重}+\text{KV Cache}+\text{激活}+\text{其他}$（再留约 $10\%$）；最大并发 $=\frac{\text{剩余 KV 空间}}{\text{每请求 KV}}$}
\meta{数字}{67.1 GiB 可用；权重 15.0 / 激活 0.91 / 其他 1.9 / KV 49.3 GiB；40.4 万 token、99 请求；FP8 227 请求；70B 四卡 102 请求；80 GB $\approx$ 74.5 GiB}
\meta{易错}{能装下权重不等于能部署：70B 的 FP8 权重刚好装进 80 GB，却没有任何地方放 KV Cache。}
\phantomsection\addcontentsline{toc}{section}{72\quad 推理成本模型}
\lec{72}{推理成本模型}
\onel{每百万 token 成本 = 租金 ÷ 吞吐，输出远比输入贵、靠并发降}
\begin{itemize}[leftmargin=1.35em,itemsep=2.2pt,topsep=1pt,parsep=0pt]
\item \begingroup\sloppy{\fontsize{10}{14.6}\selectfont 每百万 token 成本 = 每小时租金 ÷ 每小时 token 数 $\times$ 10⁶，输入按 prefill 另算。\par}\endgroup
\item \begingroup\sloppy{\fontsize{10}{14.6}\selectfont 输入计算受限，8B 在 H100 每秒约 6.5 万 token；输出吞吐 = b / TPOT。\par}\endgroup
\item \begingroup\sloppy{\fontsize{10}{14.6}\selectfont 效率 60\%、\$2.5/小时：并发 1 时输出 \$5.618/百万，并发 128 时 \$0.113。\par}\endgroup
\item \begingroup\sloppy{\fontsize{10}{14.6}\selectfont 输入与并发无关，约 \$0.018/百万；并发 128 时输出是输入的 6 倍多。\par}\endgroup
\item \begingroup\sloppy{\fontsize{10}{14.6}\selectfont 一个请求（1024 进 1024 出、并发 128）\$0.00013；70B 四卡并发 64 约 \$0.94。\par}\endgroup
\item \begingroup\sloppy{\fontsize{10}{14.6}\selectfont 量化和 GQA/MLA 降输出成本，前缀缓存降输入；FP8 输出成本约减半。\par}\endgroup
\end{itemize}
\meta{公式}{每百万 token 成本 $=\frac{\text{每小时租金}}{\text{每小时 token 数}}\times10^{6}$；输出吞吐 $=\frac{b}{\text{TPOT}}\times\text{效率}$}
\meta{数字}{\$2.5/小时；效率 60\%；\$5.618$\rightarrow$\$0.113（并发 1$\rightarrow$128）；输入 \$0.018/百万；\$0.00013/请求；70B \$0.936/百万输出；估算与实测差 2 倍以内}
\meta{易错}{API 价格不等于成本（含利润与闲时浪费）；最便宜是把并发推到显存上限，但那时 TPOT 也最长。}
\phantomsection\addcontentsline{toc}{section}{73\quad 解码策略}
\lec{73}{解码策略}
\onel{采样前用温度调尖平、Top-k/Top-p/Min-p 砍长尾，再归一化抽样。}
\begin{itemize}[leftmargin=1.35em,itemsep=2.2pt,topsep=1pt,parsep=0pt]
\item \begingroup\sloppy{\fontsize{10}{14.6}\selectfont 12 个候选从「好」0.321 到「猫」0.001，长尾词单个很小但加起来不可忽略。\par}\endgroup
\item \begingroup\sloppy{\fontsize{10}{14.6}\selectfont 温度只改尖平：T=0.5「好」0.321$\rightarrow$0.544，T=1.5 降到 0.235，但不删掉任何词。\par}\endgroup
\item \begingroup\sloppy{\fontsize{10}{14.6}\selectfont Top-k 只留概率最高 k 个：k=3 剩好、不错、晴朗，归一化 0.472/0.316/0.212。\par}\endgroup
\item \begingroup\sloppy{\fontsize{10}{14.6}\selectfont Top-p 从高到低累加，到 p 就停：p=0.8 加到第 5 个「冷」0.849，保留 5 个。\par}\endgroup
\item \begingroup\sloppy{\fontsize{10}{14.6}\selectfont Min-p 去掉低于「最大概率 $\times$ min-p」的词：min-p=0.1 门槛 0.0321，保留 7 个。\par}\endgroup
\item \begingroup\sloppy{\fontsize{10}{14.6}\selectfont 顺序为温度$\rightarrow$Top-k$\rightarrow$Top-p$\rightarrow$Min-p$\rightarrow$归一化；推理模型多用 T$\approx$0.6、Top-p$\approx$0.95。\par}\endgroup
\end{itemize}
\meta{公式}{$p_{i}=\operatorname{softmax}(\text{logit}_{i}/T)$；Top-$p$ 取最小 $n$ 使 $\sum_{i\le n}p_{i}\ge p$；Min-$p$ 保留 $p_{i}\ge\text{min\_p}\cdot\max p$}
\meta{数字}{T=0.5「好」0.544；T=1.5「好」0.235、猫 0.004；Top-k=3 剩 3 个；Top-p=0.8 保留 5 个；Min-p=0.1 保留 7 个；T=0.7+Top-p=0.9 时「好」0.458，抽 1000 次 448/274/127/86/65}
\meta{易错}{温度为 0 等于贪心（引擎直接取 argmax，不除以 0）；Top-p 不是「概率大于 p 的词」；k 固定不随模型确定性变，所以 Top-p 更常用；128256 维 logits 的排序很贵，固定种子也难完全复现。}
\chapter{推理引擎与调度}
{\footnotesize\sffamily\color{grayt}收录第 74--82 讲}\par\vspace{-4pt}
\phantomsection\addcontentsline{toc}{section}{74\quad Continuous Batching}
\lec{74}{Continuous Batching}
\onel{调度粒度从「一批请求」缩到「一步迭代」：走一个补一个，批几乎总是满的。}
\begin{itemize}[leftmargin=1.35em,itemsep=2.2pt,topsep=1pt,parsep=0pt]
\item \begingroup\sloppy{\fontsize{10}{14.6}\selectfont 静态批处理整批一起进、一起出，短请求结束后的槽位只能空等，新请求也要排队。\par}\endgroup
\item \begingroup\sloppy{\fontsize{10}{14.6}\selectfont 10 请求 4 槽位：静态 20 步 / 7.9 步 / 56\%，连续 13 步 / 5.0 步 / 87\%。\par}\endgroup
\item \begingroup\sloppy{\fontsize{10}{14.6}\selectfont 理论下限 12 步：一共 45 个 token、每步最多 4 个，连续批处理已很接近。\par}\endgroup
\item \begingroup\sloppy{\fontsize{10}{14.6}\selectfont 来自 2022 年 Orca 论文，vLLM、SGLang、TensorRT-LLM、LMDeploy 都默认启用。\par}\endgroup
\item \begingroup\sloppy{\fontsize{10}{14.6}\selectfont 实现：所有请求 token 拼成一张 [总 token 数 $\times$ d] 矩阵做一次大矩阵乘，只有注意力按请求分开。\par}\endgroup
\item \begingroup\sloppy{\fontsize{10}{14.6}\selectfont 新请求的 prefill 挤进同一步会让该步耗时暴涨，造成 TPOT 抖动（第 77 讲解决）。\par}\endgroup
\end{itemize}
\meta{公式}{槽位利用率 $=\frac{\text{busy}}{\text{slots}\times\text{steps}}$；步数下限 $=\left\lceil\frac{\text{总 token 数}}{\text{槽位数}}\right\rceil$}
\meta{数字}{10 请求、4 槽位、45 个 token；静态 20 步 / 7.9 步 / 56\%；连续 13 步 / 5.0 步 / 87\%；下限 12 步}
\meta{易错}{连续批处理不改变任何请求的计算结果，只决定「谁在同一步一起算」；dynamic batching 是等一小段时间凑批，仍属请求级。}
\phantomsection\addcontentsline{toc}{section}{75\quad PagedAttention}
\lec{75}{PagedAttention}
\onel{把 KV Cache 切成固定块按需分配，块表记录映射，浪费只剩最后一块的空位。}
\begin{itemize}[leftmargin=1.35em,itemsep=2.2pt,topsep=1pt,parsep=0pt]
\item \begingroup\sloppy{\fontsize{10}{14.6}\selectfont 按最大长度 4096 连续预留：64 个请求快照占 262144 个 token，利用率仅 32.9\%。\par}\endgroup
\item \begingroup\sloppy{\fontsize{10}{14.6}\selectfont 理想按最终长度预留 85.5\%；分页每块 16 个 token 到 99.5\%，几乎没有内部碎片。\par}\endgroup
\item \begingroup\sloppy{\fontsize{10}{14.6}\selectfont 第 i 个 token 在逻辑块 $\lfloor$i/16$\rfloor$、块内偏移 i mod 16，查块表得物理块号。\par}\endgroup
\item \begingroup\sloppy{\fontsize{10}{14.6}\selectfont 块内 16 个 token 连续，按 32/128 字节缓存行合并访存；每跨 16 个 token 才查一次块表。\par}\endgroup
\item \begingroup\sloppy{\fontsize{10}{14.6}\selectfont 每 token KV 128 KiB（2$\times$32$\times$8$\times$128$\times$2 字节），一块 16 个 token 即 2 MiB。\par}\endgroup
\item \begingroup\sloppy{\fontsize{10}{14.6}\selectfont 同池子：连续预留最多 98 个请求，分页约 299 个；并行采样共享提示块省 62\%。\par}\endgroup
\end{itemize}
\meta{公式}{逻辑块 $=\left\lfloor i/16\right\rfloor$，块内偏移 $=i\bmod 16$；每 token KV $=2\times\text{层数}\times\text{KV 头数}\times\text{头维}\times2\ \text{字节}=128$ KiB}
\meta{数字}{利用率 32.9\% / 85.5\% / 99.5\%；占用 262144 / 100758 / 86560 个 token；池 403865 个 token；98 $\rightarrow$ 299 个请求；共享 4800 $\rightarrow$ 1840 个 token}
\meta{易错}{分页不减每个 token 的 KV 大小（仍 128 KiB），只减「预留了但没用上」的浪费；共享块只读，写前要 copy-on-write。}
\phantomsection\addcontentsline{toc}{section}{76\quad Prefix Cache}
\lec{76}{Prefix Cache}
\onel{给每个装满的块算链式哈希当前缀指纹，命中的块直接复用，跳过这段 prefill。}
\begin{itemize}[leftmargin=1.35em,itemsep=2.2pt,topsep=1pt,parsep=0pt]
\item \begingroup\sloppy{\fontsize{10}{14.6}\selectfont 因果掩码决定每个 token 的 k、v 只依赖它前面的全部内容，所以只有从头相同的部分能复用。\par}\endgroup
\item \begingroup\sloppy{\fontsize{10}{14.6}\selectfont 哈希必须链式：块的指纹包含前一块的指纹，因为块内 token 的 KV 已融合前面全部上下文。\par}\endgroup
\item \begingroup\sloppy{\fontsize{10}{14.6}\selectfont 查表从头逐块进行，遇到第一个未命中就停，后面的块即使存在也不能用（因果链已断）。\par}\endgroup
\item \begingroup\sloppy{\fontsize{10}{14.6}\selectfont 未装满的末尾块不登记：它还会被写入，共享会冲突，命中粒度恒为 16 token 的整数倍。\par}\endgroup
\item \begingroup\sloppy{\fontsize{10}{14.6}\selectfont 引用计数归 0 的块不释放，进空闲 LRU 队列保留哈希；显存耗尽才从 LRU 头部淘汰叶子块。\par}\endgroup
\item \begingroup\sloppy{\fontsize{10}{14.6}\selectfont 整段命中时最后一个 token 仍要前向：KV Cache 里没有 hidden state 与 logits。\par}\endgroup
\end{itemize}
\meta{公式}{$\text{hash}_{0}=H(\text{None},\ \text{tokens}[0{:}16],\ \text{extra})$；$\text{hash}_{k}=H(\text{hash}_{k-1},\ \text{tokens}[16k{:}16k+16],\ \text{extra})$}
\meta{数字}{块 16 个 token；4 请求 154 个 token 命中 52 个、命中率 34\%；prefill 2100 token 31.9 ms $\rightarrow$ 100 token 4.8 ms；Agent 命中率上限 98\%、实践 90\textasciitilde{}95\%}
\meta{易错}{只能复用从头连续的前缀，中间相同片段不行；复用的是 KV 不是答案；10 台实例随机路由会把命中率打到 10\%。}
\phantomsection\addcontentsline{toc}{section}{77\quad Chunked Prefill}
\lec{77}{Chunked Prefill}
\onel{每步设 token 预算，先保 decode，余下按块做 prefill。}
\begin{itemize}[leftmargin=1.35em,itemsep=2.2pt,topsep=1pt,parsep=0pt]
\item \begingroup\sloppy{\fontsize{10}{14.6}\selectfont decode 一步 6.05 ms；整段 8000 token prefill 塞进来变 138.9 ms。\par}\endgroup
\item \begingroup\sloppy{\fontsize{10}{14.6}\selectfont 预算 512：先给 32 个 decode，剩 480 切给 prefill，分 17 步，每步 9.7 ms。\par}\endgroup
\item \begingroup\sloppy{\fontsize{10}{14.6}\selectfont 长提示 TTFT 几乎不变（146.9 对 138.9 ms）；预算 256 时涨到 223.5 ms。\par}\endgroup
\item \begingroup\sloppy{\fontsize{10}{14.6}\selectfont SARATHI 与 Sarathi-Serve（2023--2024）提出，后者称不停顿批处理。\par}\endgroup
\item \begingroup\sloppy{\fontsize{10}{14.6}\selectfont decode 一步本就带宽受限，还能白塞 363 个 prefill token，耗时不变。\par}\endgroup
\item \begingroup\sloppy{\fontsize{10}{14.6}\selectfont vLLM V1 默认开启，预算参数是 max\_num\_batched\_tokens。\par}\endgroup
\end{itemize}
\meta{公式}{一步耗时 $\approx\max\left(\frac{\text{计算量}}{\text{算力}},\ \frac{\text{读写字节}}{\text{带宽}}\right)$；每步 prefill token $=\text{预算}-\text{decode 请求数}$}
\meta{数字}{平时 6.05 ms；不切块 138.9/138.9 ms；2048$\rightarrow$140.4/37.4；1024$\rightarrow$147.5/19.5；512$\rightarrow$146.9/9.7；256$\rightarrow$223.5/6.4 ms；白捡 363 个 token}
\meta{易错}{切块不改变 KV 与第一个输出 token；它保护的是别人的 TPOT，长提示自己的 TTFT 不会变短，预算太小反而更长。}
\phantomsection\addcontentsline{toc}{section}{78\quad 调度与抢占}
\lec{78}{调度与抢占}
\onel{按需分配块让并发最大，池子满时抢占最后加入的请求，用少量重算换高利用率。}
\begin{itemize}[leftmargin=1.35em,itemsep=2.2pt,topsep=1pt,parsep=0pt]
\item \begingroup\sloppy{\fontsize{10}{14.6}\selectfont 保守做法按「提示长度 + 最大生成长度」预留块，永不缺块但并发低；引擎选按需分配加抢占兜底。\par}\endgroup
\item \begingroup\sloppy{\fontsize{10}{14.6}\selectfont 6 请求 16 块：按需加抢占 13 步、平均延迟 8.7 步、抢占 2 次、重算 16 个 token。\par}\endgroup
\item \begingroup\sloppy{\fontsize{10}{14.6}\selectfont 同样例子按最终长度预留 17 步、平均延迟 10.8 步、零抢占，但 D、E、F 要等到第 7、9 步。\par}\endgroup
\item \begingroup\sloppy{\fontsize{10}{14.6}\selectfont 抢占最后加入者：FCFS 更公平、它算得最少重算最省，且回到等待队列最前面不会被饿死。\par}\endgroup
\item \begingroup\sloppy{\fontsize{10}{14.6}\selectfont 发生抢占的那一步不再放新请求，避免抖动；有的实现留水位线，空闲块不足就提前停止放行。\par}\endgroup
\item \begingroup\sloppy{\fontsize{10}{14.6}\selectfont 2000 token 请求的 KV 262 MB：重算 30.3 ms，换出加换回 10.5 ms，引擎仍偏向重算。\par}\endgroup
\end{itemize}
\meta{公式}{需要的块数 $=\left\lceil\frac{\text{提示长度}+\text{已生成 token 数}}{\text{块大小}}\right\rceil$}
\meta{数字}{16 块 $\times$ 4 token；按需 13 步 / 8.7 步 / 抢占 2 次 / 重算 16 token；预留 17 步 / 10.8 步；262 MB、30.3 ms vs 10.5 ms}
\meta{易错}{抢占不是取消：已生成的 token 保留，回来重算后结果不变；它解决的是 KV 显存不够，算力不够只会让 TPOT 变长。}
\phantomsection\addcontentsline{toc}{section}{79\quad 多 LoRA 服务}
\lec{79}{多 LoRA 服务}
\onel{不合并适配器：基座所有请求合算，小路各用各的 A、B，代价随不同适配器数涨。}
\begin{itemize}[leftmargin=1.35em,itemsep=2.2pt,topsep=1pt,parsep=0pt]
\item \begingroup\sloppy{\fontsize{10}{14.6}\selectfont r=16 时一个适配器 41.9 M 参数、BF16 83.9 MB，只有 8.03 G 基座的 0.52\%。\par}\endgroup
\item \begingroup\sloppy{\fontsize{10}{14.6}\selectfont 计算只多 0.52\%，但 decode 每步要把 batch 里每个不同适配器读一遍，代价正比于个数 D。\par}\endgroup
\item \begingroup\sloppy{\fontsize{10}{14.6}\selectfont 期望 D = $\Sigma$\_i [1 - (1 - p\_i)\textasciicircum{}B]；Zipf 流量 s 越大越集中，D 越小。\par}\endgroup
\item \begingroup\sloppy{\fontsize{10}{14.6}\selectfont N=100、s=1、B=64 时平均 31.4 个适配器，每步多读 2.64 GB。\par}\endgroup
\item \begingroup\sloppy{\fontsize{10}{14.6}\selectfont Punica、S-LoRA（2023）的 SGMV 一次 kernel 算完整个 batch，每层 2 次。\par}\endgroup
\item \begingroup\sloppy{\fontsize{10}{14.6}\selectfont 显存留 max\_loras 个槽位、CPU 缓存 max\_cpu\_loras 个按 LRU 淘汰，超出就跳过请求。\par}\endgroup
\end{itemize}
\meta{公式}{$y=xW+\frac{\alpha}{r}\cdot(xA)B$；$D=\sum_{i=1}^{N}\left[1-(1-p_{i})^{B}\right]$}
\meta{数字}{r=8 21.0M/41.9MB；r=16 41.9M/83.9MB/0.52\%；r=64 167.8M/335.5MB；s=0$\rightarrow$47.4 个适配器、3.98GB；s=1$\rightarrow$31.4 个、2.64GB；一步 4.79$\rightarrow$5.58 ms；拷 84MB$\approx$1.3ms}
\meta{易错}{LoRA 不改变 KV Cache 的维度，但不同适配器算出的 K/V 不同，前缀缓存不能跨适配器共享；SGMV 只减少 kernel 调用，不减读的字节。}
\phantomsection\addcontentsline{toc}{section}{80\quad 约束解码}
\lec{80}{约束解码}
\onel{每步把接上就不合法的 token logit 置 -$\infty$，输出必符合格式。}
\begin{itemize}[leftmargin=1.35em,itemsep=2.2pt,topsep=1pt,parsep=0pt]
\item \begingroup\sloppy{\fontsize{10}{14.6}\selectfont 格式拆成 5 段；状态 = 现在第几段、这段已读几个字符，读入字符时逐段校验。\par}\endgroup
\item \begingroup\sloppy{\fontsize{10}{14.6}\selectfont 判断 token：把它的字符逐个送进状态机，中途出错就不合法，token 常跨段所以要逐字符。\par}\endgroup
\item \begingroup\sloppy{\fontsize{10}{14.6}\selectfont 只有格式完整时才允许输出结束标记，所以最终输出一定能被 json.loads 解析。\par}\endgroup
\item \begingroup\sloppy{\fontsize{10}{14.6}\selectfont 第 1 步屏蔽「好的，」「JSON：」只剩 \{（19\%），第 9 步屏蔽「十二」只剩 12（13\%）。\par}\endgroup
\item \begingroup\sloppy{\fontsize{10}{14.6}\selectfont 词表 128256 个 token，逐字符检查要几十毫秒；为每个状态提前算好掩码（Outlines 的做法）。\par}\endgroup
\item \begingroup\sloppy{\fontsize{10}{14.6}\selectfont 嵌套结构要用上下文无关文法，检查器是带栈的下推自动机；XGrammar 等压到几十微秒。\par}\endgroup
\end{itemize}
\meta{公式}{若 $\text{tokens}+t$ 送入状态机后无合法后继，则 $\text{logit}_{t}=-\infty$；仅当状态到达终止态才允许 $\langle\text{eos}\rangle$}
\meta{数字}{词表 128256；例子 11 步；\{ 原概率 19\%；12 原概率 13\%；逐 token 检查几十 ms $\rightarrow$ 优化后几十 $\mu$s}
\meta{易错}{只保证格式不保证内容；屏蔽会改变分布，把模型逼到低概率 token；生成长度截断仍会破坏格式。}
\phantomsection\addcontentsline{toc}{section}{81\quad 请求生命周期}
\lec{81}{请求生命周期}
\onel{请求过前端、引擎核心、GPU 三进程九个模块，调度$\rightarrow$前向$\rightarrow$采样每步转一圈。}
\begin{itemize}[leftmargin=1.35em,itemsep=2.2pt,topsep=1pt,parsep=0pt]
\item \begingroup\sloppy{\fontsize{10}{14.6}\selectfont 拆成前端、引擎核心、GPU 工作进程三个进程，用 ZeroMQ 通消息，免得 CPU 琐事打断 GPU 循环。\par}\endgroup
\item \begingroup\sloppy{\fontsize{10}{14.6}\selectfont 调度器每步先排 decode 各 1 token，再用剩余预算切块排 prefill，最后从等待队列放新请求。\par}\endgroup
\item \begingroup\sloppy{\fontsize{10}{14.6}\selectfont 放新请求时先查前缀缓存：命中的块直接进块表、引用计数加一，未命中的按需分配新块。\par}\endgroup
\item \begingroup\sloppy{\fontsize{10}{14.6}\selectfont 输入准备把这一步所有请求的 token 拼成长序列，附位置编号、块表与写入位置，供注意力区分请求。\par}\endgroup
\item \begingroup\sloppy{\fontsize{10}{14.6}\selectfont GPU 前向只取各请求最后位置的输出算 logits，采样先加格式掩码、除温度、Top-p 截断再抽 token。\par}\endgroup
\item \begingroup\sloppy{\fontsize{10}{14.6}\selectfont 反分词必须增量：一个汉字可能跨多个 token，要等后续 token 到齐才输出已确定的文字。\par}\endgroup
\end{itemize}
\meta{公式}{$\text{TTFT}=\text{排队时间}+\text{prefill 时间}+\text{首个 token 的采样返回}$；总延迟 $=\text{TTFT}+(\text{输出 token 数}-1)\times\text{TPOT}$}
\meta{数字}{3 个进程、9 个模块；32 层 Transformer；示例提示 12 token、每步预算 10 切成两块；R2 命中 2 块省 8 token}
\meta{易错}{停止词是文字不是 token，只能在反分词时判断，命中时引擎已多算了一两个 token，只是没返回给用户。}
\phantomsection\addcontentsline{toc}{section}{82\quad 推理压测实践}
\lec{82}{推理压测实践}
\onel{可信压测给的是吞吐-延迟曲线，以及在 SLO 下的最大吞吐 goodput。}
\begin{itemize}[leftmargin=1.35em,itemsep=2.2pt,topsep=1pt,parsep=0pt]
\item \begingroup\sloppy{\fontsize{10}{14.6}\selectfont 开环按泊松到达率加压（间隔服从指数分布），闭环固定并发数；闭环掩盖排队，可信压测用开环。\par}\endgroup
\item \begingroup\sloppy{\fontsize{10}{14.6}\selectfont 延迟必须看分布：至少报 P50 和 P99，用户体验由尾部那一次决定。\par}\endgroup
\item \begingroup\sloppy{\fontsize{10}{14.6}\selectfont 吞吐随到达率先线性增长，到容量点饱和，再加压只排队，TTFT P99 跳一两个数量级。\par}\endgroup
\item \begingroup\sloppy{\fontsize{10}{14.6}\selectfont goodput 是满足 SLO 的最大吞吐，说法来自 2024 年 DistServe，比配置、比引擎、比硬件都比它。\par}\endgroup
\item \begingroup\sloppy{\fontsize{10}{14.6}\selectfont 坑：长度分布不符、输出被结束标记截短、前缀缓存作弊、没预热、样本太少、压测机自身成瓶颈。\par}\endgroup
\item \begingroup\sloppy{\fontsize{10}{14.6}\selectfont 过载区的延迟随压测时长无限增长，只说明已经过载，不能拿来比较系统。\par}\endgroup
\end{itemize}
\meta{公式}{$\text{goodput}=\max\{\text{吞吐}\mid\text{TTFT P99}\le\text{SLO}\ \text{且}\ \text{TPOT P99}\le\text{SLO}\}$；过载区延迟随压测时长增长，不可比}
\meta{数字}{Llama-3-8B + H100、KV 池约 40 万 token：40 个/s 时 8316 token/s、TTFT P99 197 ms；45 个/s 时 8570 token/s、P99 1306 ms；60 个/s 时 8780 token/s、P99 4176 ms}
\meta{易错}{TPOT 先按请求求平均再看分布，ITL 把所有间隔放一起；分块 prefill 的偶发长间隔只在 ITL 的 P99 里显形。}
\chapter{算子与Kernel}
{\footnotesize\sffamily\color{grayt}收录第 83--89 讲}\par\vspace{-4pt}
\phantomsection\addcontentsline{toc}{section}{83\quad 算子融合}
\lec{83}{算子融合}
\onel{把几个逐元素 kernel 合成一个，让中间结果留在片上，省掉显存往返。}
\begin{itemize}[leftmargin=1.35em,itemsep=2.2pt,topsep=1pt,parsep=0pt]
\item \begingroup\sloppy{\fontsize{10}{14.6}\selectfont H100 80 GB 显存、3.35 TB/s 带宽，算力/带宽约 295 FLOPs/字节；逐元素运算算术强度不到 1。\par}\endgroup
\item \begingroup\sloppy{\fontsize{10}{14.6}\selectfont 两个 kernel 之间只能经显存交换数据；融合后中间结果留在片上，只读一次输入、写一次输出。\par}\endgroup
\item \begingroup\sloppy{\fontsize{10}{14.6}\selectfont SwiGLU 张量 117 MB：不融合 3 个 kernel 读写 8 遍约 280 $\mu$s，融合后 3 遍约 105 $\mu$s。\par}\endgroup
\item \begingroup\sloppy{\fontsize{10}{14.6}\selectfont 残差相加 + RMSNorm 用 7 个 kernel 读写 10 遍约 100 $\mu$s，融合后 4 遍约 40 $\mu$s。\par}\endgroup
\item \begingroup\sloppy{\fontsize{10}{14.6}\selectfont 常见融合：silu\_and\_mul、fused\_add\_rms\_norm、矩阵乘尾部融合、合并 QKV 与 gate/up 投影。\par}\endgroup
\item \begingroup\sloppy{\fontsize{10}{14.6}\selectfont 边界在要整行或整张量的运算（softmax 分母、量化 scale）；读写量少了，峰值显存不一定少。\par}\endgroup
\end{itemize}
\meta{公式}{$\text{耗时}\approx\text{读写遍数}\times\frac{\text{张量字节}}{\text{带宽}}$；整体加速 $=\frac{1}{(1-f)+f/s}$（Amdahl）}
\meta{数字}{SwiGLU 117 MB 读写 8$\rightarrow$3 遍、280$\rightarrow$105 $\mu$s（940 MB）；RMSNorm 34 MB 10$\rightarrow$4 遍、100$\rightarrow$40 $\mu$s；一层 481$\rightarrow$185 $\mu$s，同层矩阵乘约 1812 $\mu$s；persistent kernel 126$\rightarrow$114 $\mu$s；Qwen3-8B 请求 817 ms，激活 13.6$\rightarrow$12.2 ms}
\meta{易错}{kernel 加速比不等于请求加速比：热点占 20\% 时上限 1.25 倍，SwiGLU 快 2.6 倍，请求只快约 0.16\%。}
\phantomsection\addcontentsline{toc}{section}{84\quad FlashAttention}
\lec{84}{FlashAttention}
\onel{分块算注意力配在线 softmax，S 不落显存，s$^2$ 级读写降到四分之一。}
\begin{itemize}[leftmargin=1.35em,itemsep=2.2pt,topsep=1pt,parsep=0pt]
\item \begingroup\sloppy{\fontsize{10}{14.6}\selectfont S = QKᵀ/$\surd$d 是 s$\times$s：s = 8192 时一个头 128 MiB，32 个头共 4 GiB。\par}\endgroup
\item \begingroup\sloppy{\fontsize{10}{14.6}\selectfont 标准实现 5 个算子把 S 读写 8 遍：s = 8192 时约 1082 MB、323 $\mu$s，真正计算只要 35 $\mu$s。\par}\endgroup
\item \begingroup\sloppy{\fontsize{10}{14.6}\selectfont 在线 softmax 维护 m、l、acc，最大值变大时旧值乘 e\textasciicircum{}(m\_旧 - m\_新) 修正，结果精确等价。\par}\endgroup
\item \begingroup\sloppy{\fontsize{10}{14.6}\selectfont FlashAttention（2022）把 S 留在片上：读写约 277 MB、83 $\mu$s，显存从 s$^2$ 级降到 s 级。\par}\endgroup
\item \begingroup\sloppy{\fontsize{10}{14.6}\selectfont 99 KB 共享内存下 b = 64 放 82 行 Q、读写 404 MiB；b = 1 只读 260 MiB，更新次数却多 41 倍。\par}\endgroup
\item \begingroup\sloppy{\fontsize{10}{14.6}\selectfont FlashAttention-2（2023）改循环顺序与并行，-3（2024）用 TMA 异步搬 K、V 并支持 FP8。\par}\endgroup
\end{itemize}
\meta{公式}{$l\leftarrow l\cdot e^{m-m_{\text{new}}}+\sum e^{x-m_{\text{new}}}$，$\text{acc}\leftarrow\text{acc}\cdot e^{m-m_{\text{new}}}+\sum e^{x-m_{\text{new}}}\cdot v$，输出 $=\text{acc}/l$}
\meta{数字}{s = 8192、d = 128：1082 MB/323 $\mu$s $\rightarrow$ 277 MB/83 $\mu$s；实测 RTX PRO 6000 快约 33 倍，新增显存 848$\rightarrow$22 MiB；A100 上达峰值 50\%\textasciitilde{}73\%；b = 64 时在线更新 12,800 次、b = 1 时 524,288 次}
\meta{易错}{既不是近似也没减少计算量：乘加次数不变、复杂度仍是 s$^2$，省的只是显存读写和占用。}
\phantomsection\addcontentsline{toc}{section}{85\quad Decode 注意力 Kernel}
\lec{85}{Decode 注意力 Kernel}
\onel{decode 只有 1 个查询，按块表读散落 KV，再沿上下文切段用满 SM。}
\begin{itemize}[leftmargin=1.35em,itemsep=2.2pt,topsep=1pt,parsep=0pt]
\item \begingroup\sloppy{\fontsize{10}{14.6}\selectfont 每读一个 k 只与 q 做一次点积，算术强度个位数，远低于 295 FLOPs/字节，完全受带宽限制。\par}\endgroup
\item \begingroup\sloppy{\fontsize{10}{14.6}\selectfont 第 i 个 token 在物理块 table[i ÷ B] 的第 i mod B 个位置，vLLM 默认 B = 16。\par}\endgroup
\item \begingroup\sloppy{\fontsize{10}{14.6}\selectfont 按请求$\times$KV 头分工：批 1、32K 上下文只有 8 个工作单元，一层要读 128 MiB KV，约 661 $\mu$s。\par}\endgroup
\item \begingroup\sloppy{\fontsize{10}{14.6}\selectfont Flash-Decoding（2023，split-K）每 512 token 切一段，512 个工作单元，约 40 $\mu$s。\par}\endgroup
\item \begingroup\sloppy{\fontsize{10}{14.6}\selectfont 切段多写一组 (m, l, acc) 再加一次归约；批 64、4K 上下文时切不切都约 321 $\mu$s。\par}\endgroup
\item \begingroup\sloppy{\fontsize{10}{14.6}\selectfont 切段改变浮点结果：段数与合并次序让同一请求在不同 batch 下输出不同，要做 batch invariance。\par}\endgroup
\end{itemize}
\meta{公式}{$\text{物理块}=\text{table}[\,i\div B\,]$，块内位置 $=i\bmod B$；合并 $l=\sum_s l_s\cdot e^{m_s-M}$，$\text{acc}=\sum_s \text{acc}_s\cdot e^{m_s-M}$}
\meta{数字}{B = 16；每 token 每层 KV 4096 字节；批1/32K 661$\rightarrow$40 $\mu$s；批1/4K 83$\rightarrow$10 $\mu$s；批64/4K 321/321 $\mu$s；RMSNorm 切 8 段 24$\rightarrow$32.1 MiB、1$\rightarrow$3 个 kernel、误差 111$\rightarrow$9 ULP}
\meta{易错}{Flash-Decoding 不为 prefill 服务；合并是精确的，但段数随 batch 变，同一前缀在不同 batch 里概率会不同。}
\phantomsection\addcontentsline{toc}{section}{86\quad CUDA 编程模型}
\lec{86}{CUDA 编程模型}
\onel{CUDA 用线程、块、网格三级组织，一个块住一个 SM，块内靠共享内存合作。}
\begin{itemize}[leftmargin=1.35em,itemsep=2.2pt,topsep=1pt,parsep=0pt]
\item \begingroup\sloppy{\fontsize{10}{14.6}\selectfont i = blockIdx.x $\times$ blockDim.x + threadIdx.x 决定线程负责的元素，越界线程跳过。\par}\endgroup
\item \begingroup\sloppy{\fontsize{10}{14.6}\selectfont 块最多 1024 个线程，一个块整个跑在一个 SM 上；块之间不能直接合作、顺序无保证。\par}\endgroup
\item \begingroup\sloppy{\fontsize{10}{14.6}\selectfont 线程每 32 个编成一个 warp、同执行一条指令：避免分支发散，相邻线程读相邻地址。\par}\endgroup
\item \begingroup\sloppy{\fontsize{10}{14.6}\selectfont 存储三级：寄存器私有最快、共享内存块内共享（几十到两百多 KB）、全局内存即显存。\par}\endgroup
\item \begingroup\sloppy{\fontsize{10}{14.6}\selectfont 按 T$\times$T 分块在共享内存复用，显存读量降到 1/T：K = 4096 从 274.9 GB 降到 4.3 GB。\par}\endgroup
\item \begingroup\sloppy{\fontsize{10}{14.6}\selectfont 共享内存 32 个 bank：行跨度 32 时读一列撞同一 bank 要 32 轮，补到 33 或 swizzle 后 1 轮。\par}\endgroup
\end{itemize}
\meta{公式}{$i=\text{blockIdx}.x\times\text{blockDim}.x+\text{threadIdx}.x$；显存读取量 $\approx2K^{3}/T$ 个元素；$T=t_{\text{搬}}+(n-1)\cdot\max(t_{\text{搬}},\ t_{\text{算}})+t_{\text{算}}$}
\meta{数字}{H100 132 SM；SM 资源 2048 线程槽 / 65,536 寄存器 / 228 KB 共享内存；128$\times$128 块驻留 2 块、占用率 25\%；在途需 15.2 KB、16 warp 给 32 KiB；K = 4096 分块 82$\rightarrow$1.3 ms；双缓冲 6.28$\rightarrow$5.44 $\mu$s}
\meta{易错}{CUDA 的块是线程组，PagedAttention 的块是 KV 存储段；占用率高不等于快，在途字节够了再加没有收益。}
\phantomsection\addcontentsline{toc}{section}{87\quad 矩阵乘法 Kernel 与 Tensor Core}
\lec{87}{矩阵乘法 Kernel 与 Tensor Core}
\onel{三层分块喂 Tensor Core，块按波分给 SM，波次悬崖与块太少都浪费。}
\begin{itemize}[leftmargin=1.35em,itemsep=2.2pt,topsep=1pt,parsep=0pt]
\item \begingroup\sloppy{\fontsize{10}{14.6}\selectfont 三层分块：线程块级走共享内存、warp 级累加在寄存器、指令级交给 Tensor Core。\par}\endgroup
\item \begingroup\sloppy{\fontsize{10}{14.6}\selectfont A100 的 mma 一个 warp 算 16$\times$8；H100 的 wgmma 四个 warp 算 64$\times$N，A、B 从共享内存异步读。\par}\endgroup
\item \begingroup\sloppy{\fontsize{10}{14.6}\selectfont 波次悬崖：1408$^2$ 一波 121 块、利用率 91.7\%；1536$^2$ 两波 144 块、54.5\%，时间翻倍。\par}\endgroup
\item \begingroup\sloppy{\fontsize{10}{14.6}\selectfont decode 16$\times$1024$\times$4096 只有 16 块，13 TFLOPS 即峰值 1.3\%，比理想 2.6 $\mu$s 慢 4 倍。\par}\endgroup
\item \begingroup\sloppy{\fontsize{10}{14.6}\selectfont split-K 沿 K 切 8 段得 128 块，2.9 $\mu$s 接近只读一遍权重；Flash-Decoding 同一思路。\par}\endgroup
\item \begingroup\sloppy{\fontsize{10}{14.6}\selectfont 库的对策：cuBLAS 按形状选配置、CUTLASS 模板实例化、Stream-K 均分、decode 档位特化。\par}\endgroup
\end{itemize}
\meta{公式}{$\text{波次利用率}=\frac{\text{块数}}{\text{波数}\times\text{SM}}$；算术强度 $\approx\frac{2MNK}{2KN}\approx M$；split-K 的块数 $=\lceil M/M_t\rceil\cdot\lceil N/N_t\rceil\cdot s$}
\meta{数字}{H100 132 SM；BF16 峰值 989 TFLOPS、拐点 295 FLOP/字节；1408$^2$ 与 1536$^2$ 为 17.9/35.8 $\mu$s、907/539 TFLOPS；16$\times$1024 为 10.5 $\mu$s、13 TFLOPS；split-K=8 为 2.9 $\mu$s、47 TFLOPS}
\meta{易错}{用上 Tensor Core 不等于跑到峰值；性能曲线是阶梯形：1536 比 1408 慢一倍，再长到 2048 又和 1536 一样。}
\phantomsection\addcontentsline{toc}{section}{88\quad Triton 入门}
\lec{88}{Triton 入门}
\onel{以一块数据为单位写 kernel，线程分工与共享内存交给编译器。}
\begin{itemize}[leftmargin=1.35em,itemsep=2.2pt,topsep=1pt,parsep=0pt]
\item \begingroup\sloppy{\fontsize{10}{14.6}\selectfont Triton（OpenAI 2021）嵌在 Python 里：program\_id 定位块，arange 出整块下标，mask 管越界。\par}\endgroup
\item \begingroup\sloppy{\fontsize{10}{14.6}\selectfont 融合 SwiGLU 十几行：load g、u，片上算 silu(g)$\cdot$u，只 store 结果，读写 8 遍降到 3 遍。\par}\endgroup
\item \begingroup\sloppy{\fontsize{10}{14.6}\selectfont RMSNorm 一个 program 一行，tl.sum 做块内归约，先转 FP32 保精度，other=0 不影响平方和。\par}\endgroup
\item \begingroup\sloppy{\fontsize{10}{14.6}\selectfont 矩阵乘用二维网格沿 K 循环 load A、B 小块，tl.dot 累加到 FP32 累加器，编译器出 Tensor Core 指令。\par}\endgroup
\item \begingroup\sloppy{\fontsize{10}{14.6}\selectfont autotune 试块大小、num\_warps、num\_stages；候选先按资源排除、再估算、最后按真实调用比例实测。\par}\endgroup
\item \begingroup\sloppy{\fontsize{10}{14.6}\selectfont 动态形状分通用、分桶、特化三种：前 65 组通用、第 66 组起分桶、第 91 组起特化。\par}\endgroup
\end{itemize}
\meta{公式}{$\text{平均时间}=5p+20(1-p)\ \mu\text{s}$，优于 $10\ \mu\text{s}$ 需 $p>2/3$；总时间 $=\text{准备}+\text{组数}\times\text{每组时间}$}
\meta{数字}{n = 1000、BLOCK = 256 $\rightarrow$ 4 个程序、末块 232 个有效；转置 16$\times$16 7.9 $\mu$s vs 32$\times$32 2.9 $\mu$s；分桶多算 45\%（8192 vs 5632 行）；每组 16.9/12.3/6.8 ms；搜索 1.2 亿次调用才回本}
\meta{易错}{Triton 的程序不是线程，也和 NVIDIA 的 Triton 推理服务器无关；块形状是编译时常数，一般取 2 的幂。}
\phantomsection\addcontentsline{toc}{section}{89\quad CUDA Graph}
\lec{89}{CUDA Graph}
\onel{CPU 发射跟不上 GPU 时，把整步 kernel 序列录成图回放即可。}
\begin{itemize}[leftmargin=1.35em,itemsep=2.2pt,topsep=1pt,parsep=0pt]
\item \begingroup\sloppy{\fontsize{10}{14.6}\selectfont eager 模式每个 kernel 要过 Python、框架和驱动，CPU 花几微秒。\par}\endgroup
\item \begingroup\sloppy{\fontsize{10}{14.6}\selectfont kernel 比发射还短时，GPU 算完一个只能等下一个，一步由 CPU 决定。\par}\endgroup
\item \begingroup\sloppy{\fontsize{10}{14.6}\selectfont 捕获只记 kernel、顺序、参数和地址，回放时 CPU 一次调用、几十微秒。\par}\endgroup
\item \begingroup\sloppy{\fontsize{10}{14.6}\selectfont 8 卡并行 8B：545 kernel，GPU 1.93 ms，图 1.95 ms。\par}\endgroup
\item \begingroup\sloppy{\fontsize{10}{14.6}\selectfont 图和地址要固定，按批大小各捕一张，实际批向上补，prefill 一般不用。\par}\endgroup
\item \begingroup\sloppy{\fontsize{10}{14.6}\selectfont vLLM v0.6.0 拆 API 进程、多步调度、异步输出，吞吐提到 2.7 倍。\par}\endgroup
\end{itemize}
\meta{公式}{第 $i$ 个 kernel 开始 $=\max(\text{它被发射完},\ \text{第 }i-1\text{ 个做完})$；图重放的临界输入 $2X\div B=15\ \mu\text{s}\to X\approx12.8\ \text{MiB}$}
\meta{数字}{481 个 kernel/5.18 ms（单卡 8B）；545 个/1.93 ms（8 卡，图 1.95 ms）；241 个/1.09 ms（1B，图 1.11 ms）；回放省 15 $\mu$s；流水 2020 $\mu$s，GPU 提速后 2005 $\mu$s；主机也提速 505 $\mu$s}
\meta{易错}{CUDA Graph 不让 kernel 变快，只填 kernel 之间的空隙；输入超过约 12.8 MiB 时复制开销会吃掉全部收益。}
\chapter{模型压缩与加速}
{\footnotesize\sffamily\color{grayt}收录第 90--98 讲}\par\vspace{-4pt}
\phantomsection\addcontentsline{toc}{section}{90\quad 量化基础}
\lec{90}{量化基础}
\onel{整数量化用 scale 把实数映射到格点，误差不超过 scale÷2。}
\begin{itemize}[leftmargin=1.35em,itemsep=2.2pt,topsep=1pt,parsep=0pt]
\item \begingroup\sloppy{\fontsize{10}{14.6}\selectfont 对称量化：scale = max|x| ÷ (2\textasciicircum{}(b-1)-1)，用时乘回 scale。\par}\endgroup
\item \begingroup\sloppy{\fontsize{10}{14.6}\selectfont INT8：scale 0.01819，最大误差 0.0077，不超过 scale÷2。\par}\endgroup
\item \begingroup\sloppy{\fontsize{10}{14.6}\selectfont 位宽越低误差越陡：INT4 0.0725，INT3 0.1362，INT2 已不可用。\par}\endgroup
\item \begingroup\sloppy{\fontsize{10}{14.6}\selectfont 非对称量化多存零点 z，把 [min,max] 映射到 [0,2\textasciicircum{}b-1]，用满格点。\par}\endgroup
\item \begingroup\sloppy{\fontsize{10}{14.6}\selectfont 一个离群值把 2.31 换成 12，其余 7 个数的平均误差从 0.07 涨到 0.49。\par}\endgroup
\item \begingroup\sloppy{\fontsize{10}{14.6}\selectfont W4A16 只省读权重的带宽，decode 收益大、prefill 几乎没有。\par}\endgroup
\end{itemize}
\meta{公式}{$\text{scale}=\frac{\max|x|}{2^{b-1}-1}$，$q=\operatorname{round}(x\div\text{scale})$，$\hat{x}=q\times\text{scale}$，误差 $\le\text{scale}\div2$；非对称 $q=\operatorname{round}(x\div\text{scale})+z$，$\hat{x}=(q-z)\times\text{scale}$}
\meta{数字}{Llama-3-8B：BF16 16.06 GB/4.79 ms、INT8 8.03 GB/2.40 ms、INT4 4.14 GB/1.24 ms；70B 由 140 GB 降到约 35 GB；每 128 个配 16 位 scale 合 4.125 比特；每 4 个一组把 0.49 降到 0.17}
\meta{易错}{存成 INT4 不等于用 INT4 计算：W4A16 读进片上先反量化成 BF16 再乘，省的是字节不是计算。}
\phantomsection\addcontentsline{toc}{section}{91\quad 权重量化：GPTQ 与 AWQ}
\lec{91}{权重量化：GPTQ 与 AWQ}
\onel{权重量化该最小化输出误差而非权重误差，激活大的通道那行最关键。}
\begin{itemize}[leftmargin=1.35em,itemsep=2.2pt,topsep=1pt,parsep=0pt]
\item \begingroup\sloppy{\fontsize{10}{14.6}\selectfont RTN 就是除以 scale 再四舍五入，8 比特几乎无损，4 比特明显掉点。\par}\endgroup
\item \begingroup\sloppy{\fontsize{10}{14.6}\selectfont 输出误差 = $\Sigma$ⱼ xⱼ(Wⱼ - Ŵⱼ)，激活大的通道把该行误差放大到输出。\par}\endgroup
\item \begingroup\sloppy{\fontsize{10}{14.6}\selectfont AWQ 用 x$\cdot$W=(x/s)$\cdot$(s$\cdot$W) 把重要的行放大再量化，s 与激活大小挂钩。\par}\endgroup
\item \begingroup\sloppy{\fontsize{10}{14.6}\selectfont AWQ 的 s 可离线合并进前一层 RMSNorm，4 比特输出误差降到 6.34\%。\par}\endgroup
\item \begingroup\sloppy{\fontsize{10}{14.6}\selectfont GPTQ 逐行量化，按 H=XᵀX 的逆把取整误差摊给未量化的行补偿。\par}\endgroup
\item \begingroup\sloppy{\fontsize{10}{14.6}\selectfont 先量化 H 对角线最大的行：3 比特时 RTN 23.99\%，GPTQ 13.07\%。\par}\endgroup
\end{itemize}
\meta{公式}{$x\cdot W=(x/s)\cdot(s\cdot W)$；$s_j=(\text{平均}|x_j|)^{\alpha}$；补偿矩阵来自 $H=X^{\top}X$ 的逆；输出误差 $=\frac{\lVert X\cdot W-X\cdot\hat{W}\rVert}{\lVert X\cdot W\rVert}$}
\meta{数字}{校准 256 token、4 个离群通道大 20 倍；4 比特输出误差 10.56\%/6.34\%/5.30\%；GPTQ 权重误差 0.016 反而最大；3 比特 23.99\%/14.71\%/13.07\%；175B 量化几个 GPU 小时}
\meta{易错}{GPTQ 的权重平均误差反而最大，因为它有意把权重调离原值去抵消输出误差------权重误差不是优化目标。}
\phantomsection\addcontentsline{toc}{section}{92\quad 激活量化与 SmoothQuant}
\lec{92}{激活量化与 SmoothQuant}
\onel{整数矩阵乘只允许激活每 token 一个 scale，得把难度匀给权重。}
\begin{itemize}[leftmargin=1.35em,itemsep=2.2pt,topsep=1pt,parsep=0pt]
\item \begingroup\sloppy{\fontsize{10}{14.6}\selectfont 整数矩阵乘要求 scale 在求和维不变，所以激活只能每 token 一个。\par}\endgroup
\item \begingroup\sloppy{\fontsize{10}{14.6}\selectfont LLM.int8() 2022：6\textasciitilde{}7B 以上每层激活有几个固定离群通道，大几十倍。\par}\endgroup
\item \begingroup\sloppy{\fontsize{10}{14.6}\selectfont 直接 W8A8：整体激活误差 1.52\%，普通通道 17.40\%，输出 1.70\%。\par}\endgroup
\item \begingroup\sloppy{\fontsize{10}{14.6}\selectfont SmoothQuant 用 (X/s)$\cdot$(s$\cdot$W) 等价变换，把激活的难度匀给权重。\par}\endgroup
\item \begingroup\sloppy{\fontsize{10}{14.6}\selectfont $\alpha$ = 0.5 时两边最大值÷中位数都是 7.7，输出误差 1.70\% 降到 0.50\%。\par}\endgroup
\item \begingroup\sloppy{\fontsize{10}{14.6}\selectfont s 离线算好可合并进前一层，推理零开销；W4A4 还得靠正交旋转。\par}\endgroup
\end{itemize}
\meta{公式}{$X\cdot W=(X/s)\cdot(s\cdot W)$；$s_j=\frac{\max|X_j|^{\alpha}}{\max|W_j|^{1-\alpha}}$；$\alpha=0.5$ 时两边都是 $\sqrt{\max|X_j|\cdot\max|W_j|}$}
\meta{数字}{普通通道 max $\approx$ 3.0、离群通道 $\approx$ 178；直接 W8A8 输出 1.70\%；$\alpha$ = 0 / 0.5 / 1 时输出误差 1.75\% / 0.50\% / 1.88\%，两边难度 59.8:1.0 $\rightarrow$ 7.7:7.7 $\rightarrow$ 1.0:59.8}
\meta{易错}{离群通道不是噪声，去掉它们模型效果大降，只能安置不能截断；逐通道量化激活硬件做不了。}
\phantomsection\addcontentsline{toc}{section}{93\quad FP8 与 FP4}
\lec{93}{FP8 与 FP4}
\onel{浮点格点在对数坐标上均匀，FP4 只有 8 个值，必须配细粒度分组缩放。}
\begin{itemize}[leftmargin=1.35em,itemsep=2.2pt,topsep=1pt,parsep=0pt]
\item \begingroup\sloppy{\fontsize{10}{14.6}\selectfont 值 = $\pm$(1 + M/2\textasciicircum{}m) $\times$ 2\textasciicircum{}(E-bias)，格点在对数坐标上均匀。\par}\endgroup
\item \begingroup\sloppy{\fontsize{10}{14.6}\selectfont E4M3 最大 448、精度好，推理多用它；E5M2 最大 57344、范围大。\par}\endgroup
\item \begingroup\sloppy{\fontsize{10}{14.6}\selectfont FP4 E2M1 只有 0、0.5、1、1.5、2、3、4、6 八个非负值。\par}\endgroup
\item \begingroup\sloppy{\fontsize{10}{14.6}\selectfont 重尾 4096 个数整体 scale：INT4 51.25\%，FP4 34.02\%。\par}\endgroup
\item \begingroup\sloppy{\fontsize{10}{14.6}\selectfont MXFP4 每 32 个数一组，scale 只能取 2 的幂，合 4.25 位。\par}\endgroup
\item \begingroup\sloppy{\fontsize{10}{14.6}\selectfont NVFP4 每 16 个数一组、scale 用 E4M3，合 4.5 位。\par}\endgroup
\end{itemize}
\meta{公式}{$(-1)^{S}\times(1+M/2^{m})\times2^{E-\text{bias}}$；次正规数 $(-1)^{S}\times(M/2^{m})\times2^{1-\text{bias}}$；平均位宽 $=4+8\div\text{组大小}$}
\meta{数字}{E4M3：127 个非负值、最小 2$^-$⁹ $\approx$ 0.00195、每段 8 个格点；E5M2：124 个、最小 1.5$\times$10$^-$⁵；整体 scale 误差 INT8 3.22\% 对 E4M3 2.44\%；每 32 个一组 INT4 15.09\%、MXFP4 15.91\%；NVFP4 9.25\%；H100 FP8 1979 对 BF16 989 TFLOPS；8B 用 FP8 约 8 GB、NVFP4 约 4.5 GB}
\meta{易错}{分组 scale 由 Blackwell 的 Tensor Core 直接乘进矩阵乘，组再细也不拖慢计算；FP8 在 A100 上只能转回 BF16 算，省显存不省计算。}
\phantomsection\addcontentsline{toc}{section}{94\quad KV Cache 量化}
\lec{94}{KV Cache 量化}
\onel{KV 字节减半则同显存并发翻倍，满载 TPOT 几乎不变、吞吐翻倍。}
\begin{itemize}[leftmargin=1.35em,itemsep=2.2pt,topsep=1pt,parsep=0pt]
\item \begingroup\sloppy{\fontsize{10}{14.6}\selectfont 80 GB 卡跑 8B：权重 16 GB，KV 池 49.3 GiB。\par}\endgroup
\item \begingroup\sloppy{\fontsize{10}{14.6}\selectfont 每 token 的 KV 128 KiB，4096 token 的请求只能放 98 个。\par}\endgroup
\item \begingroup\sloppy{\fontsize{10}{14.6}\selectfont FP8 把每 token 压到 64 KiB，池子能放 197 个请求。\par}\endgroup
\item \begingroup\sloppy{\fontsize{10}{14.6}\selectfont 满载时 decode 主要读 KV，字节减半 TPOT 几乎不变、吞吐翻倍。\par}\endgroup
\item \begingroup\sloppy{\fontsize{10}{14.6}\selectfont K 有固定离群通道（与 RoPE 慢转维度有关），V 均匀，故 K 按通道分。\par}\endgroup
\item \begingroup\sloppy{\fontsize{10}{14.6}\selectfont INT4：K、V 都按 token 量化误差 64.71\%，K 按通道 18.84\%。\par}\endgroup
\end{itemize}
\meta{公式}{$\text{可放 token 数}=\text{KV 池}\div\text{每 token 字节}$；满载一步时间 $\approx\text{KV 总字节}\div\text{带宽}$，故字节减半时 TPOT 不变、吞吐翻倍}
\meta{数字}{每 token 128 / 64 / 33 / 17 KiB；池子 403,865 / 807,731 / 1,566,508 个 token；98 / 197 / 382 / 742 个请求；TPOT 20.5 / 20.6 / 20.6 / 20.6 ms；吞吐 4,781 / 9,573 / 18,564 / 36,042 token/s；注意力输出误差 FP8 3.87\%、INT8 5.00\%}
\meta{易错}{KV 量化不让单个请求变快：并发低时 decode 在读权重，TPOT 几乎不动；收益全在同样显存服务更多请求和更长上下文。}
\phantomsection\addcontentsline{toc}{section}{95\quad KV Cache 压缩与驱逐}
\lec{95}{KV Cache 压缩与驱逐}
\onel{注意力集中在开头几个汇聚点，驱逐必须保留它们，丢了损失暴涨。}
\begin{itemize}[leftmargin=1.35em,itemsep=2.2pt,topsep=1pt,parsep=0pt]
\item \begingroup\sloppy{\fontsize{10}{14.6}\selectfont 128K 请求 KV 16 GiB，只留 4096 个 token 降到 0.5 GiB。\par}\endgroup
\item \begingroup\sloppy{\fontsize{10}{14.6}\selectfont 尾巴 63.5\% 的注意力落在开头 4 个 token 上，它们只占 0.9\%。\par}\endgroup
\item \begingroup\sloppy{\fontsize{10}{14.6}\selectfont 汇聚点源于 softmax 要凑满权重和 1，内容不重要、位置才重要。\par}\endgroup
\item \begingroup\sloppy{\fontsize{10}{14.6}\selectfont 只留最近 B 个是灾难：B = 256 也多 2.81；保留开头 4 个再加最近。\par}\endgroup
\item \begingroup\sloppy{\fontsize{10}{14.6}\selectfont H2O 按累计注意力挑，SnapKV 用最后 32 个查询打分。\par}\endgroup
\item \begingroup\sloppy{\fontsize{10}{14.6}\selectfont 驱逐和 PagedAttention 的 16 token 块冲突，主流引擎默认不做。\par}\endgroup
\end{itemize}
\meta{公式}{损失增量换算困惑度：$e^{\Delta}$ 倍，$\Delta=0.539\to1.71\times$、$0.134\to1.14\times$、$0.013\to1.01\times$}
\meta{数字}{每 token 128 KiB；128K 一条请求 16 GiB、只留 4096 个是 3.1\%；不驱逐损失 3.316；B = 32 时四种策略 +4.49 / +0.54 / +0.57 / +0.50；B = 256 时 +2.81 / +0.10 / +0.01 / +0.03}
\meta{易错}{驱逐扔掉的 token 再也找不回来，语言建模损失好看不代表长文问答也好；随机丢 128 个恰好丢了第 0 个，直接多 4.72。}
\phantomsection\addcontentsline{toc}{section}{96\quad 投机解码}
\lec{96}{投机解码}
\onel{小模型猜 k 个、大模型一次前向验证，输出与大模型自己生成完全一样。}
\begin{itemize}[leftmargin=1.35em,itemsep=2.2pt,topsep=1pt,parsep=0pt]
\item \begingroup\sloppy{\fontsize{10}{14.6}\selectfont 小批量 decode 一步都花在读权重上，处理 1 个和 32 个几乎同价。\par}\endgroup
\item \begingroup\sloppy{\fontsize{10}{14.6}\selectfont 小模型照常猜 k 个 token，大模型一次前向同时给出 k+1 个位置的预测。\par}\endgroup
\item \begingroup\sloppy{\fontsize{10}{14.6}\selectfont 从前往后比对，第一个错的位置换成大模型的预测，后面草稿全作废。\par}\endgroup
\item \begingroup\sloppy{\fontsize{10}{14.6}\selectfont 每轮至少前进 1 个、最多 k+1 个，被作废的草稿要回退 KV 位置。\par}\endgroup
\item \begingroup\sloppy{\fontsize{10}{14.6}\selectfont 期望前进 E=(1-$\alpha$\textasciicircum{}(k+1))÷(1-$\alpha$)，加速比 E÷(1+k$\cdot$c)。\par}\endgroup
\item \begingroup\sloppy{\fontsize{10}{14.6}\selectfont 采样用拒绝采样：以 min(1,p/q) 接受，拒绝时从 max(0,p-q) 重抽。\par}\endgroup
\end{itemize}
\meta{公式}{$E=\frac{1-\alpha^{k+1}}{1-\alpha}$；加速比 $=E\div(1+k\cdot c)$；接受率 $\alpha=\sum\min(p,\ q)$}
\meta{数字}{1 / 5 / 32 个 token 前向都约 4.83 ms；$\alpha$ = 0.8、k = 4 $\rightarrow$ E = 3.362、2.80$\times$；$\alpha$ = 0.8 最佳 k = 8 $\rightarrow$ 3.09$\times$；$\alpha$ = 0.6 最佳 k = 4 $\rightarrow$ 1.92$\times$；$\alpha$ = 0.9、k = 10 $\rightarrow$ 4.57$\times$；接受率 0.699 对理论 0.700}
\meta{易错}{加速比不等于每轮前进的 token 数，还要除以 1 + k$\cdot$c；k 不是越大越好，$\alpha$ 低时多猜纯属浪费。}
\phantomsection\addcontentsline{toc}{section}{97\quad EAGLE 与 Medusa}
\lec{97}{EAGLE 与 Medusa}
\onel{Medusa 多头跳猜、EAGLE 特征层逐步猜，树形注意力验证整棵树。}
\begin{itemize}[leftmargin=1.35em,itemsep=2.2pt,topsep=1pt,parsep=0pt]
\item \begingroup\sloppy{\fontsize{10}{14.6}\selectfont Medusa 在最后一层隐藏状态接小头，一次前向并行猜后面 4 个位置。\par}\endgroup
\item \begingroup\sloppy{\fontsize{10}{14.6}\selectfont 越靠后的头越不准：不知道中间几个 token，一轮平均只前进 1.99 个。\par}\endgroup
\item \begingroup\sloppy{\fontsize{10}{14.6}\selectfont 每头取前 2 名组成候选树，深度 3 共 14 个候选，一次前向全验证。\par}\endgroup
\item \begingroup\sloppy{\fontsize{10}{14.6}\selectfont 树形注意力只换掩码：每个候选只看得见前文和自己的祖先，结果不变。\par}\endgroup
\item \begingroup\sloppy{\fontsize{10}{14.6}\selectfont 沿大模型认可的路径从根往下走，能走多深就接受多少，链只接受第 1 名。\par}\endgroup
\item \begingroup\sloppy{\fontsize{10}{14.6}\selectfont EAGLE 草稿网约一层 Transformer，输入当前特征加一步 token 嵌入。\par}\endgroup
\end{itemize}
\meta{公式}{一轮期望前进 $=1+\sum_{i=1}^{k}\prod_{j=1}^{i}p_j$（链取 top-1，树取 top-2）；深度 $d$ 的候选数 $=2^{1}+\cdots+2^{d}$}
\meta{数字}{top-1 命中率 0.60/0.45/0.35/0.28；top-2 0.78/0.62/0.52/0.44；链 1.99；深度 2/3/4 树 2.26/2.52/2.63；候选 6/14/30；MTP 第二个 token 接受率 85\%\textasciitilde{}90\%、约 1.8 倍}
\meta{易错}{树形注意力不是新的注意力机制，只是换掩码让多条路径挤进一次前向；Medusa 的头也不是 MoE 专家。}
\phantomsection\addcontentsline{toc}{section}{98\quad 蒸馏与剪枝}
\lec{98}{蒸馏与剪枝}
\onel{蒸馏拟合老师分布、剪枝去权重；只有规则稀疏才在硬件上真变快。}
\begin{itemize}[leftmargin=1.35em,itemsep=2.2pt,topsep=1pt,parsep=0pt]
\item \begingroup\sloppy{\fontsize{10}{14.6}\selectfont 老师分布含暗知识（Hinton，2015）：猫 0.744、狗 0.166，狗也说得通。\par}\endgroup
\item \begingroup\sloppy{\fontsize{10}{14.6}\selectfont 师生 logits 同除温度 T 再 softmax，损失是 KL(p老师‖p学生)。\par}\endgroup
\item \begingroup\sloppy{\fontsize{10}{14.6}\selectfont 只用标准答案时梯度 p-y 把狗和汽车一起压低；蒸馏梯度反而提高狗。\par}\endgroup
\item \begingroup\sloppy{\fontsize{10}{14.6}\selectfont 大模型多用数据蒸馏：老师生成回答，学生在上面做 SFT。\par}\endgroup
\item \begingroup\sloppy{\fontsize{10}{14.6}\selectfont Wanda 用 |w|$\cdot$‖x‖ 判重要性，只需少量校准数据、不用再训练。\par}\endgroup
\item \begingroup\sloppy{\fontsize{10}{14.6}\selectfont 只有 2:4 与整块结构化能真加速；实践先剪再蒸馏，如 8B 剪成 4B。\par}\endgroup
\end{itemize}
\meta{公式}{$p=\operatorname{softmax}(z/T)$，$L=\text{KL}(p_{\text{老师}}\,\Vert\,p_{\text{学生}})$；重要性 $=|w|\cdot\|x_j\|$}
\meta{数字}{T=1 猫 0.744/狗 0.166；T=4 猫 0.336/狗 0.231；剪一半误差 27.12\%(|w|)、10.72\%(Wanda)、13.92\%(2:4)、23.43\%(结构化)；2:4 理论快 2 倍、权重近半；R1 的 80 万条数据}
\meta{易错}{蒸馏不是压缩已有模型，而是训练另一个小模型；零散的 0 在普通矩阵乘法里照样要算。}
\chapter{分布式与集合通信}
{\footnotesize\sffamily\color{grayt}收录第 99--109 讲}\par\vspace{-4pt}
\phantomsection\addcontentsline{toc}{section}{99\quad 为什么需要多卡}
\lec{99}{为什么需要多卡}
\onel{一卡不够有三因：放不下、decode 受带宽限、流量扛不住；前两个切分。}
\begin{itemize}[leftmargin=1.35em,itemsep=2.2pt,topsep=1pt,parsep=0pt]
\item \begingroup\sloppy{\fontsize{10}{14.6}\selectfont 权重 = 参数量 $\times$ 字节数：70B 的 BF16 一张 H100 装不下。\par}\endgroup
\item \begingroup\sloppy{\fontsize{10}{14.6}\selectfont 每 token KV（KiB）：70B 320、405B 504、V3 仅 68.6。\par}\endgroup
\item \begingroup\sloppy{\fontsize{10}{14.6}\selectfont 按每卡 90\% 可用算：70B 要 4 张 H100，405B 14 张，V3 19 张。\par}\endgroup
\item \begingroup\sloppy{\fontsize{10}{14.6}\selectfont 单卡读一遍 70B 权重 42 ms、405B 242 ms，即每秒 24 与 4 个。\par}\endgroup
\item \begingroup\sloppy{\fontsize{10}{14.6}\selectfont 权重分到 8 卡同时读，70B 降到 5.3 ms、405B 30.3 ms。\par}\endgroup
\item \begingroup\sloppy{\fontsize{10}{14.6}\selectfont 训练每参数 16 字节：8B 要 128 GB、70B 要 1130 GB。\par}\endgroup
\end{itemize}
\meta{公式}{$\text{权重}=\text{参数量}\times\text{字节数}$；decode 单请求 token/s $\approx\beta\div\text{权重字节}$；训练状态 $\approx16$ 字节/参数；卡数 $=\lceil(\text{权重}+\text{KV})\div(80\ \text{GB}\times0.9)\rceil$}
\meta{数字}{H100 80 GB；每 token KV 70B 320 KiB、405B 504 KiB、V3 68.6 KiB；32 个 8K 请求的 KV 85.9/135.3/18.4 GB；卡数 4/14/19；NVLink 450 GB/s、跨机 50 GB/s}
\meta{易错}{刚装得下只是最低要求，此时没给 KV Cache 留空间；多卡也不是免费加速，通信可能吃掉收益。}
\phantomsection\addcontentsline{toc}{section}{100\quad 集合通信原语}
\lec{100}{集合通信原语}
\onel{六种原语归为求和、拼接、交换；All-Reduce 是后两者之和。}
\begin{itemize}[leftmargin=1.35em,itemsep=2.2pt,topsep=1pt,parsep=0pt]
\item \begingroup\sloppy{\fontsize{10}{14.6}\selectfont 4 卡按位置相加得 [60, 64, 68, 72]：Reduce 只放卡 0。\par}\endgroup
\item \begingroup\sloppy{\fontsize{10}{14.6}\selectfont All-Reduce 让每张卡都拿到完整的和，是最常用的原语。\par}\endgroup
\item \begingroup\sloppy{\fontsize{10}{14.6}\selectfont Reduce-Scatter 求和后每卡只留一段，All-Gather 拼成完整。\par}\endgroup
\item \begingroup\sloppy{\fontsize{10}{14.6}\selectfont 求和后切段再拼起就是 All-Reduce，ZeRO 只用前半省一半通信。\par}\endgroup
\item \begingroup\sloppy{\fontsize{10}{14.6}\selectfont All-to-All 把第 j 段发给卡 j，等效转置数据矩阵，MoE 靠它。\par}\endgroup
\item \begingroup\sloppy{\fontsize{10}{14.6}\selectfont 每卡至少收发：All-Reduce 2(p-1)/p$\cdot$n，其余 (p-1)/p$\cdot$n。\par}\endgroup
\end{itemize}
\meta{公式}{$\text{All-Reduce}=\text{Reduce-Scatter}+\text{All-Gather}$；All-Reduce $2(p-1)/p\cdot n$，其余三种 $(p-1)/p\cdot n$}
\meta{数字}{p=4：[60, 64, 68, 72]；All-Reduce 1.5n，其余 0.75n；张量并行每层 2 次、MoE 每层 2 次}
\meta{易错}{All-Gather 后各卡相同、All-to-All 后各不相同；集合通信要所有卡到场，最慢的卡决定速度。}
\phantomsection\addcontentsline{toc}{section}{101\quad Ring 与 Tree All-Reduce}
\lec{101}{Ring 与 Tree All-Reduce}
\onel{Ring 带宽达下限但步数多，Tree 只需约 2log$_2$p 步赢小消息。}
\begin{itemize}[leftmargin=1.35em,itemsep=2.2pt,topsep=1pt,parsep=0pt]
\item \begingroup\sloppy{\fontsize{10}{14.6}\selectfont 朴素做法卡 0 收发 2(p-1)$\cdot$n，时间随卡数线性增长。\par}\endgroup
\item \begingroup\sloppy{\fontsize{10}{14.6}\selectfont Ring 把数据切 p 块，先逐卡累加再逐块传一圈，共 2(p-1) 步。\par}\endgroup
\item \begingroup\sloppy{\fontsize{10}{14.6}\selectfont Ring 每卡共发 2(p-1)/p$\cdot$n，正好达到下限，所有链路同时传。\par}\endgroup
\item \begingroup\sloppy{\fontsize{10}{14.6}\selectfont Tree 叶子向根相加再向叶广播，步数约 2$\cdot$log$_2$p，512 卡时 18 步。\par}\endgroup
\item \begingroup\sloppy{\fontsize{10}{14.6}\selectfont NCCL 用双二叉树均衡负担，NVSwitch 的 SHARP 在交换机里做加法。\par}\endgroup
\item \begingroup\sloppy{\fontsize{10}{14.6}\selectfont 分层通信让慢速网络只传 1/8：机内 RS、跨机 AR、机内 AG。\par}\endgroup
\end{itemize}
\meta{公式}{$T_{\text{ring}}=2(p-1)\cdot(\alpha+n/(p\cdot\beta))$；$T_{\text{tree}}\approx2\lceil\log_2 p\rceil\cdot\alpha+2n/\beta$}
\meta{数字}{1 GB/512 卡：朴素 20 s、Ring 45 ms、树 40 ms；64 KB/512 卡：Ring 5.1 ms、树 0.09 ms；4 卡每卡 1.5n、朴素 6n}
\meta{易错}{带宽最优不等于最快：小消息延迟项主导时 Tree 反超；环只是逻辑上的，不要求物理成环。}
\phantomsection\addcontentsline{toc}{section}{102\quad 通信代价模型}
\lec{102}{通信代价模型}
\onel{T = $\alpha$ + m/$\beta$，分界点 m* = $\alpha$$\cdot$$\beta$：小消息卡延迟，大消息卡带宽。}
\begin{itemize}[leftmargin=1.35em,itemsep=2.2pt,topsep=1pt,parsep=0pt]
\item \begingroup\sloppy{\fontsize{10}{14.6}\selectfont $\alpha$ 是每条消息的固定开销（几微秒），$\beta$ 是单向带宽，规格要先换算。\par}\endgroup
\item \begingroup\sloppy{\fontsize{10}{14.6}\selectfont m* = $\alpha$$\cdot$$\beta$：NVLink 约 1.35 MB，IB 约 300 KB，以下卡延迟。\par}\endgroup
\item \begingroup\sloppy{\fontsize{10}{14.6}\selectfont NVLink 发 8 KB 3.0 $\mu$s，99.4\% 是延迟，带宽仅 2.7 GB/s。\par}\endgroup
\item \begingroup\sloppy{\fontsize{10}{14.6}\selectfont 集合通信把每步相加：Ring 是 2(p-1)$\cdot$$\alpha$ + 2(p-1)/p$\cdot$m/$\beta$。\par}\endgroup
\item \begingroup\sloppy{\fontsize{10}{14.6}\selectfont 一步到位只付一次 $\alpha$，但每卡要收 (p-1)$\cdot$m，小消息时最快。\par}\endgroup
\item \begingroup\sloppy{\fontsize{10}{14.6}\selectfont 70B 张量并行 decode：Ring 7.05 ms，一步到位 1.78 ms。\par}\endgroup
\end{itemize}
\meta{公式}{$T=\alpha+m/\beta$；$m^{*}=\alpha\cdot\beta$；$T_{\text{ring}}=2(p-1)\cdot(\alpha+m/(p\beta))$}
\meta{数字}{$\alpha$=3 $\mu$s/$\beta$=450 GB/s（NVLink）、$\alpha$=6 $\mu$s/$\beta$=50 GB/s（IB）；512 KB 延迟占 72.5\%；每步 160 次 524 KB；16 GB 梯度 16 节点 0.60 s；算法带宽 26.7、总线带宽 50.0 GB/s}
\meta{易错}{标称常是双向总和（NVLink 900 GB/s 即单向 450）；看 nccl-tests 要用总线带宽才可与链路带宽比。}
\phantomsection\addcontentsline{toc}{section}{103\quad 数据并行与 ZeRO}
\lec{103}{数据并行与 ZeRO}
\onel{数据并行复制完整模型并 All-Reduce 梯度，ZeRO 依次切三块状态。}
\begin{itemize}[leftmargin=1.35em,itemsep=2.2pt,topsep=1pt,parsep=0pt]
\item \begingroup\sloppy{\fontsize{10}{14.6}\selectfont N 卡各存完整模型、各算 1/N 数据，梯度取平均等于单卡大批次。\par}\endgroup
\item \begingroup\sloppy{\fontsize{10}{14.6}\selectfont 每参数 16 字节：BF16 权重与梯度各 2，FP32 主权重与 Adam 共 12。\par}\endgroup
\item \begingroup\sloppy{\fontsize{10}{14.6}\selectfont ZeRO-1 只切优化器状态，每卡 44.2 GB，参数更新后再 All-Gather。\par}\endgroup
\item \begingroup\sloppy{\fontsize{10}{14.6}\selectfont ZeRO-2 改用 Reduce-Scatter 同步梯度，每卡只存 1/N。\par}\endgroup
\item \begingroup\sloppy{\fontsize{10}{14.6}\selectfont ZeRO-3 连参数也切，用时 All-Gather、算完丢弃，每卡 16.1 GB。\par}\endgroup
\item \begingroup\sloppy{\fontsize{10}{14.6}\selectfont ZeRO 不管激活，长序列时要靠激活重算或张量、序列并行。\par}\endgroup
\end{itemize}
\meta{公式}{每卡状态：数据并行 $16\Psi$；ZeRO-1 $4\Psi+12\Psi/N$；ZeRO-2 $2\Psi+14\Psi/N$；ZeRO-3 $16\Psi/N$}
\meta{数字}{8B、8 卡：128.5 $\rightarrow$ 44.2 $\rightarrow$ 30.1 $\rightarrow$ 16.1 GB；每步每卡通信 32/32/32/48 GB；ZeRO-3 是数据并行的 1.5 倍}
\meta{易错}{ZeRO 不改变计算结果；卡数增加时没切的部分不缩小，ZeRO-1 的权重与梯度始终 4Ψ。}
\phantomsection\addcontentsline{toc}{section}{104\quad 混合精度训练与激活重计算}
\lec{104}{混合精度训练与激活重计算}
\onel{BF16 算矩阵乘、FP32 记主参数，16 字节/参数；重计算多算 33\%。}
\begin{itemize}[leftmargin=1.35em,itemsep=2.2pt,topsep=1pt,parsep=0pt]
\item \begingroup\sloppy{\fontsize{10}{14.6}\selectfont FP16 指数只 5 位，2$\times$10$^-$⁸ 的梯度下溢成 0，要靠损失缩放。\par}\endgroup
\item \begingroup\sloppy{\fontsize{10}{14.6}\selectfont BF16 尾数 7 位，1 附近间隔 2$^-$⁷，权重加 10$^-$$^4$ 走千步仍是 1。\par}\endgroup
\item \begingroup\sloppy{\fontsize{10}{14.6}\selectfont 所以要留 FP32 主参数：2+2+4+8 = 16 字节/参数由此而来。\par}\endgroup
\item \begingroup\sloppy{\fontsize{10}{14.6}\selectfont 每层激活 $\approx$ s$\cdot$b$\cdot$h$\cdot$(34 + 5as/h) 字节，后一项随 s$^2$ 涨。\par}\endgroup
\item \begingroup\sloppy{\fontsize{10}{14.6}\selectfont 全部重计算只留每层输入，训练从 6N 变 8N，多算 33\%。\par}\endgroup
\item \begingroup\sloppy{\fontsize{10}{14.6}\selectfont FlashAttention 天然去掉平方项，与选择性重计算等效。\par}\endgroup
\end{itemize}
\meta{公式}{每层激活 $\approx s\cdot b\cdot h\cdot(34+5as/h)$ 字节；训练 FLOPs $6N\to8N$（多 33\%）}
\meta{数字}{H100 BF16 989 vs FP32 67 TFLOPS；8B s=8192：380.1 / 36.5 / 3.29 GB，多算 8.3\% 与 33\%；s=32768：146.0 / 13.15 GB}
\meta{易错}{重计算省的是激活而不是参数状态（后者靠 ZeRO）；多算是一次前向，33\% 而非翻倍。}
\phantomsection\addcontentsline{toc}{section}{105\quad 张量并行}
\lec{105}{张量并行}
\onel{TP 把矩阵乘法拆到 t 张卡，每卡只读 1/t 权重。}
\begin{itemize}[leftmargin=1.35em,itemsep=2.2pt,topsep=1pt,parsep=0pt]
\item \begingroup\sloppy{\fontsize{10}{14.6}\selectfont 按列切 A 得到完整的部分列；按行切只得部分和，要先相加。\par}\endgroup
\item \begingroup\sloppy{\fontsize{10}{14.6}\selectfont FFN 里 A 按列、B 按行，激活逐元素各卡独立算，只在末尾通信一次。\par}\endgroup
\item \begingroup\sloppy{\fontsize{10}{14.6}\selectfont 反过来 A 按行切，ReLU 之前就得先通信，FFN 要通信两次。\par}\endgroup
\item \begingroup\sloppy{\fontsize{10}{14.6}\selectfont 注意力 QKV 按头切、输出投影按行切，每层前向共 2 次通信。\par}\endgroup
\item \begingroup\sloppy{\fontsize{10}{14.6}\selectfont KV Cache 随头切开；GQA 只有 8 个 KV 头，TP=8 后不能再切。\par}\endgroup
\item \begingroup\sloppy{\fontsize{10}{14.6}\selectfont 70B 用 TP=8：每卡 17.6 GB，一步 5.3 ms 加通信 1.78 ms。\par}\endgroup
\end{itemize}
\meta{公式}{按列切 $Y=[X\cdot A_0,\ X\cdot A_1]$；按行切 $Y=X_0\cdot A_0+X_1\cdot A_1$，需 All-Reduce}
\meta{数字}{TP=1 每卡 141 GB、42.1 ms；TP=2 70.5 GB、21.7 ms；TP=4 35.2 GB、11.6 ms；TP=8 17.6 GB、7.0 ms；通信占比 3\%$\rightarrow$25\%；2019 年 Megatron-LM}
\meta{易错}{TP 要求每张卡都参与每一层，任何一张慢全组都等；它不减少每卡处理的 token 数，这点和数据并行相反。}
\phantomsection\addcontentsline{toc}{section}{106\quad 流水并行}
\lec{106}{流水并行}
\onel{PP 按层切模型，段间只传激活，气泡占 (p-1)/(m+p-1)。}
\begin{itemize}[leftmargin=1.35em,itemsep=2.2pt,topsep=1pt,parsep=0pt]
\item \begingroup\sloppy{\fontsize{10}{14.6}\selectfont 按层切成 p 段，段间只点对点传一次激活，通信量小，适合跨机。\par}\endgroup
\item \begingroup\sloppy{\fontsize{10}{14.6}\selectfont 只有一批数据时同时只有一张卡干活，4 段利用率仅 25\%，空等叫气泡。\par}\endgroup
\item \begingroup\sloppy{\fontsize{10}{14.6}\selectfont GPipe 先做完所有前向再做反向、最后一次更新，气泡只在灌满排空。\par}\endgroup
\item \begingroup\sloppy{\fontsize{10}{14.6}\selectfont 1F1B 先热身再一前一后交替，每卡同时存的激活从 m 份降到最多 p 份。\par}\endgroup
\item \begingroup\sloppy{\fontsize{10}{14.6}\selectfont 推理时 PP 不缩短单请求延迟：70B 四段要 42.1 ms。\par}\endgroup
\item \begingroup\sloppy{\fontsize{10}{14.6}\selectfont 实际是机内 TP、机间 PP：405B 用两台 8 卡机器，TP=8 加 PP=2。\par}\endgroup
\end{itemize}
\meta{公式}{$\text{气泡占比}=\frac{p-1}{m+p-1}$}
\meta{数字}{p=4：m=1 总时长 12、气泡 75\%；m=4 21、42.9\%；m=8 33、27.3\%；m=16 57、15.8\%；同一请求一步：PP 四段要读完 141 GB 约 42.1 ms，TP=4 约 10.5 ms；段间 32 个 token 的隐藏状态 524 KB、跨机约 16 $\mu$s}
\meta{易错}{气泡公式假设每段耗时相等，段不均衡时实际气泡更大；微批次只把梯度累加，不改变训练结果。}
\phantomsection\addcontentsline{toc}{section}{107\quad 序列与上下文并行}
\lec{107}{序列与上下文并行}
\onel{SP 切逐 token 运算，CP 把注意力沿序列切开。}
\begin{itemize}[leftmargin=1.35em,itemsep=2.2pt,topsep=1pt,parsep=0pt]
\item \begingroup\sloppy{\fontsize{10}{14.6}\selectfont Megatron 2022 的 SP 把 RMSNorm、残差沿序列切开。\par}\endgroup
\item \begingroup\sloppy{\fontsize{10}{14.6}\selectfont All-Gather 拼回序列，再用 Reduce-Scatter 切开，通信量不变。\par}\endgroup
\item \begingroup\sloppy{\fontsize{10}{14.6}\selectfont Ring Attention 让 K、V 沿环转 P 步，每步在线 softmax 累加。\par}\endgroup
\item \begingroup\sloppy{\fontsize{10}{14.6}\selectfont 因果掩码下按顺序切各卡算量 36、100、164、228，之字形后都是 132。\par}\endgroup
\item \begingroup\sloppy{\fontsize{10}{14.6}\selectfont Ulysses 用两次 All-to-All 在按序列切和按头切之间切换。\par}\endgroup
\item \begingroup\sloppy{\fontsize{10}{14.6}\selectfont decode 时每卡算自己那段 KV 的部分 softmax，再跨卡合并。\par}\endgroup
\end{itemize}
\meta{数字}{1M token 的 Llama-3-8B KV 122 GiB，8 卡每卡 15 GiB，每层每步传 488 MiB；Ulysses 并行度受头数限制}
\meta{易错}{「序列并行」有多种用法，读文档要看它切了什么；之字形切分只在因果掩码下需要。}
\phantomsection\addcontentsline{toc}{section}{108\quad 专家并行与 All-to-All}
\lec{108}{专家并行与 All-to-All}
\onel{EP 把专家分到各卡，每个 MoE 层两次 All-to-All。}
\begin{itemize}[leftmargin=1.35em,itemsep=2.2pt,topsep=1pt,parsep=0pt]
\item \begingroup\sloppy{\fontsize{10}{14.6}\selectfont 每层 256 个路由专家、每 token 选 8 个，总参数 6710 亿。\par}\endgroup
\item \begingroup\sloppy{\fontsize{10}{14.6}\selectfont 一个专家约 4400 万参数、一层 113 亿；切小后 GPU 跑不满。\par}\endgroup
\item \begingroup\sloppy{\fontsize{10}{14.6}\selectfont 流程是路由、dispatch 发副本、专家计算、combine 加权送回。\par}\endgroup
\item \begingroup\sloppy{\fontsize{10}{14.6}\selectfont 路由不保证均匀：最忙的卡收 35 个副本，平均只有 16 个。\par}\endgroup
\item \begingroup\sloppy{\fontsize{10}{14.6}\selectfont decode 用 320 张卡 EP320，几乎每卡一个专家，并复制热门专家。\par}\endgroup
\item \begingroup\sloppy{\fontsize{10}{14.6}\selectfont 对策是限制跨机范围、先机内汇集、算通重叠和 DeepEP。\par}\endgroup
\end{itemize}
\meta{数字}{每层分发 7.3 MB（FP8）、合并 14.7 MB（BF16），50 GB/s 下 0.44 ms、58 层约 26 ms；prefill 32 卡 EP32、decode 320 卡 EP320；每 token 最多发往 4 个节点}
\meta{易错}{EP 只切 MoE 层，注意力还要别的并行方式；All-to-All 的数据量取决于路由，不是固定值。}
\phantomsection\addcontentsline{toc}{section}{109\quad 并行策略与拓扑映射}
\lec{109}{并行策略与拓扑映射}
\onel{通信越频繁越靠里：TP 在 NVLink，PP、DP 放在跨机网络。}
\begin{itemize}[leftmargin=1.35em,itemsep=2.2pt,topsep=1pt,parsep=0pt]
\item \begingroup\sloppy{\fontsize{10}{14.6}\selectfont 机内 NVLink 每向 450 GB/s，机间 50 GB/s，差 9 倍。\par}\endgroup
\item \begingroup\sloppy{\fontsize{10}{14.6}\selectfont TP 每层两次 All-Reduce 最难重叠，PP、DP 量小容易藏住。\par}\endgroup
\item \begingroup\sloppy{\fontsize{10}{14.6}\selectfont 同样 TP=8，组放机内 7.0 ms；故意让组跨机变 21.6 ms。\par}\endgroup
\item \begingroup\sloppy{\fontsize{10}{14.6}\selectfont TP=16 跨机，通信 31.95 ms 吃掉一切，TPOT 34.6 ms。\par}\endgroup
\item \begingroup\sloppy{\fontsize{10}{14.6}\selectfont TP=4 机内 4 份，单请求 11.6 ms，集群吞吐却更高。\par}\endgroup
\item \begingroup\sloppy{\fontsize{10}{14.6}\selectfont Llama 3 405B 训练按 TP、CP、PP、DP 从里到外排。\par}\endgroup
\end{itemize}
\meta{数字}{TP=8 机内$\times$2：17.6 GB、7.0 ms、9083 tok/s；TP=8 组跨机：21.6 ms、2958 tok/s；TP=8$\times$PP=2：7.1 ms、9062 tok/s；TP=4$\times$4：11.6 ms、11,071 tok/s；GB200 NVL72 让 72 卡共享一个 NVLink 域}
\meta{易错}{并行度大小和放在哪些卡上是两件事；rank 编号没和物理拓扑对齐不会报错，只会悄悄变慢。}
\chapter{分布式推理与集群}
{\footnotesize\sffamily\color{grayt}收录第 110--116 讲}\par\vspace{-4pt}
\phantomsection\addcontentsline{toc}{section}{110\quad 推理中的张量并行}
\lec{110}{推理中的张量并行}
\onel{TP 用更多卡换低延迟，但每卡吞吐随卡数下降。}
\begin{itemize}[leftmargin=1.35em,itemsep=2.2pt,topsep=1pt,parsep=0pt]
\item \begingroup\sloppy{\fontsize{10}{14.6}\selectfont 每卡只读 1/t 权重和 KV、只算 1/t，代价是每层两次 All-Reduce。\par}\endgroup
\item \begingroup\sloppy{\fontsize{10}{14.6}\selectfont 70B 批 64：TP 从 4 到 8，TPOT 15.3 降到 10.0 ms。\par}\endgroup
\item \begingroup\sloppy{\fontsize{10}{14.6}\selectfont TP=8 的 10.0 ms 里约三成是通信，批越小占比越高。\par}\endgroup
\item \begingroup\sloppy{\fontsize{10}{14.6}\selectfont 选型顺序：先放得下，再用 TP 满足 TPOT，剩下的卡做副本。\par}\endgroup
\item \begingroup\sloppy{\fontsize{10}{14.6}\selectfont prefill 计算受限：70B 的 TTFT 从 669 ms 降到 245 ms。\par}\endgroup
\item \begingroup\sloppy{\fontsize{10}{14.6}\selectfont MLA 的 KV 压成共享潜向量没法按头切，改用 DP 注意力。\par}\endgroup
\end{itemize}
\meta{公式}{$\text{TPOT}(t)\approx\frac{\text{单卡耗时}}{t}+2\times\text{层数}\times\text{All-Reduce}(\text{批大小}\times d\times2\ \text{字节})$}
\meta{数字}{8B 批 64：TP=1 9.9 ms/6450 tok/s，TP=8 2.0 ms/4094 tok/s；70B TP=8 每卡显存 23.0 GB；TP 4$\rightarrow$8 扩展效率约 77\%；8 份 TP=1 共约 51,600 tok/s}
\meta{易错}{说「TP 让推理变快」要指明是延迟变快，每卡效率在下降；批大时 KV 和计算占主导，TP 效率反而好一些。}
\phantomsection\addcontentsline{toc}{section}{111\quad MoE 推理与大规模专家并行}
\lec{111}{MoE 推理与大规模专家并行}
\onel{大规模 EP 下专家算得快，瓶颈在网络：每 token 每步约 10 MB。}
\begin{itemize}[leftmargin=1.35em,itemsep=2.2pt,topsep=1pt,parsep=0pt]
\item \begingroup\sloppy{\fontsize{10}{14.6}\selectfont EP=64 每卡每层 4 个专家，FP8 权重 176 MB，读一遍 0.053 ms。\par}\endgroup
\item \begingroup\sloppy{\fontsize{10}{14.6}\selectfont 每 token 每层要送收 8 份隐藏状态，约 172 KB。\par}\endgroup
\item \begingroup\sloppy{\fontsize{10}{14.6}\selectfont 50 GB/s 下每层 All-to-All 0.45 ms，是专家计算的八倍多。\par}\endgroup
\item \begingroup\sloppy{\fontsize{10}{14.6}\selectfont 带宽提到 450 GB/s，每卡吞吐涨近 5 倍。\par}\endgroup
\item \begingroup\sloppy{\fontsize{10}{14.6}\selectfont 每卡吞吐上限 = 带宽 ÷ 每 token 流量，50 GB/s 时约 5011。\par}\endgroup
\item \begingroup\sloppy{\fontsize{10}{14.6}\selectfont EP 越大越不均：4 倍热点在 EP8 只多 8\%，EP256 就是 4 倍。\par}\endgroup
\end{itemize}
\meta{公式}{每卡吞吐 $\le\frac{\text{网络带宽}}{\text{每 token 每步流量}}$（$58$ 层 $\times172\ \text{KB}\approx10.0\ \text{MB}$）}
\meta{数字}{每层专家 0.053 ms vs All-to-All 0.45 ms；一步 31.1 ms、通信约 26 ms；25$\rightarrow$450 GB/s 时 2405$\rightarrow$22,371 tok/s；每传 1 字节只有 4096 次运算，H100 FP8÷50 GB/s=39,580；冗余专家机内第 3 批回本、跨机第 23 批；CPU/GPU 交点 AVX-512 71、AMX 688 个 token}
\meta{易错}{MoE 省的是计算不是通信；加大批次流量同比例涨、上限不变，只能摊薄读专家权重。}
\phantomsection\addcontentsline{toc}{section}{112\quad PD 分离}
\lec{112}{PD 分离}
\onel{PD 分离把 prefill 与 decode 分到不同实例。}
\begin{itemize}[leftmargin=1.35em,itemsep=2.2pt,topsep=1pt,parsep=0pt]
\item \begingroup\sloppy{\fontsize{10}{14.6}\selectfont 混在一起时，带 prefill 块的一步从 5 ms 变成 30 ms。\par}\endgroup
\item \begingroup\sloppy{\fontsize{10}{14.6}\selectfont prefill 偏好大算力和大 TP，decode 偏好大带宽大显存。\par}\endgroup
\item \begingroup\sloppy{\fontsize{10}{14.6}\selectfont 12 请求/s 时，混合部署 token 间隔 P99 23.7 ms。\par}\endgroup
\item \begingroup\sloppy{\fontsize{10}{14.6}\selectfont PD 分离后同一场景 P99 降到 6 ms 以内，吞吐不变。\par}\endgroup
\item \begingroup\sloppy{\fontsize{10}{14.6}\selectfont 配比错更差：1P+3D 的 prefill 排队，TTFT P99 涨到 120 ms。\par}\endgroup
\item \begingroup\sloppy{\fontsize{10}{14.6}\selectfont KV 搬一次：2000 token 的 8B 模型 262 MB、约 5.2 ms。\par}\endgroup
\end{itemize}
\meta{公式}{$\mu=\min(\mu_P,\ \mu_D,\ \text{链路带宽}\div\text{每请求 KV})$}
\meta{数字}{Qwen3-8B 每 token KV 144 KiB、8K 输入 1.125 GiB、25 GB/s 传 48.3 ms；A100 prefill 0.856 s/decode 1.764 GPU 秒，H20 1.805 s/0.878；异构 4A100+4H20：共置 3.02、PD 4.55、对调 2.22 请求/s；同构 8$\times$A100 共置 3.05、3P5D 2.83、分块捎带上限 4.53}
\meta{易错}{同构集群上 PD 分离不提高吞吐（每请求 GPU 秒不变、卡只能整数分），提吞吐要靠分块 prefill，PD 买的是隔离。}
\phantomsection\addcontentsline{toc}{section}{113\quad KV Cache 传输与分层存储}
\lec{113}{KV Cache 传输与分层存储}
\onel{读回已算KV还是重算，看存储带宽是否高于分界带宽，据此分层缓存。}
\begin{itemize}[leftmargin=1.35em,itemsep=2.2pt,topsep=1pt,parsep=0pt]
\item \begingroup\sloppy{\fontsize{10}{14.6}\selectfont 分界带宽=每token KV÷每token重算时间，快于它读回划算，慢于它不如重算。\par}\endgroup
\item \begingroup\sloppy{\fontsize{10}{14.6}\selectfont 模型越大、KV越小（GQA、MLA），分界带宽越低，越慢的存储也值得用。\par}\endgroup
\item \begingroup\sloppy{\fontsize{10}{14.6}\selectfont 容量下限是r$\times$T：KV产生速率r就是分界带宽，HBM只留得住几十秒。\par}\endgroup
\item \begingroup\sloppy{\fontsize{10}{14.6}\selectfont 读回可与计算重叠：逐层预取，再提前读入读比算多出的几层，读取被完全盖住。\par}\endgroup
\item \begingroup\sloppy{\fontsize{10}{14.6}\selectfont SSD瓶颈不是带宽而是寿命，1 DWPD每秒只准写几十MB，写入前要准入判断。\par}\endgroup
\item \begingroup\sloppy{\fontsize{10}{14.6}\selectfont 传输要合并成大块、走GPUDirect RDMA不绕CPU、用块哈希目录记谁有哪段前缀。\par}\endgroup
\end{itemize}
\meta{公式}{$\text{分界带宽}=\frac{\text{每 token KV}}{\text{每 token 重算时间}}$；$C\ge r\times T$；$T_{\text{写}}+k\cdot T_{\text{读}}<k\cdot T_{\text{重算}}$；$\text{缓冲}=\text{链路带宽}\times(T_{\text{读}}-T_{\text{算}})$}
\meta{数字}{三个模型分界带宽 8.6 / 2.3 / 1.9 GB/s（128 / 320 / 68.6 KiB 每 token）；8K 前缀 1.05 GB、重算 121 ms、CPU 读回 21 ms、SSD 105 ms、远端 350 ms；r=1.41 GB/s；保留 45 / 195 / 2721 s；10 s 要 14.1 GB、600 s 要 847 GB；链路 25 GB/s 复用 1 次回本、2 GB/s 3 次、1.25 GB/s 不回本；读 48.3 ms、算 30.9 ms、先读后算 79.2 ms、逐层 49.2 ms、提前读入 436 MB 即 14 层；SSD 写入 44.4 MB/s 是 r 的 3.1\%；队列预取命中 86\% 对 LRU 58\%}
\meta{易错}{SSD 带宽够但寿命不够：按 1 DWPD 把所有新 KV 都写进去，5 年寿命不到两个月就写完了。}
\phantomsection\addcontentsline{toc}{section}{114\quad 路由与负载均衡}
\lec{114}{路由与负载均衡}
\onel{路由要兼顾前缀命中与负载均衡，只在负载不超上限的实例里选命中最多的。}
\begin{itemize}[leftmargin=1.35em,itemsep=2.2pt,topsep=1pt,parsep=0pt]
\item \begingroup\sloppy{\fontsize{10}{14.6}\selectfont 负载：12个应用各2000token提示、300会话聊1--8轮，共1288请求、4实例。\par}\endgroup
\item \begingroup\sloppy{\fontsize{10}{14.6}\selectfont 轮询与最空闲负载最均匀，但同一会话的下一轮散落各处，历史前缀命中不了。\par}\endgroup
\item \begingroup\sloppy{\fontsize{10}{14.6}\selectfont 按应用固定命中率最高，但热点集中，最忙实例是平均的1.86倍，拖住整个集群。\par}\endgroup
\item \begingroup\sloppy{\fontsize{10}{14.6}\selectfont 缓存感知只在负载不超平均1.2倍的实例里选命中最多的，命中96.8\%、负载1.12倍。\par}\endgroup
\item \begingroup\sloppy{\fontsize{10}{14.6}\selectfont 路由要比首token=max(排队,状态就绪)+剩余计算，只看缓存或空闲都会错。\par}\endgroup
\item \begingroup\sloppy{\fontsize{10}{14.6}\selectfont 路由器可自建近似前缀树，也可让实例上报块哈希；负载要按命中后还要算多少token衡量。\par}\endgroup
\end{itemize}
\meta{公式}{$\text{首 token}=\max(Q,\ R)+C$；$R=\text{查找}+\text{KV}\div\text{链路带宽}+\text{PCIe 搬运}$；优化目标是 goodput，命中率只是手段}
\meta{数字}{五种策略命中率/最忙倍数：71.6/1.01、71.4/1.00、84.8/1.12、92.1/1.86、96.8/1.12；281 / 908 / 282 / 138 ms；追平需 6.30 GB/s；串行 150 ms 对并行 90 ms；p=0.9 均值 196.5 ms 但 10\% 请求要 968 ms}
\meta{易错}{平均命中率高不代表尾延迟好：缓存以概率 p 还在，那 1-p 的请求整段重算，p99 由它决定。}
\phantomsection\addcontentsline{toc}{section}{115\quad 超节点上的推理}
\lec{115}{超节点上的推理}
\onel{超节点只改连法：放不进8卡的模型和MoE受益，放得进8卡的更慢。}
\begin{itemize}[leftmargin=1.35em,itemsep=2.2pt,topsep=1pt,parsep=0pt]
\item \begingroup\sloppy{\fontsize{10}{14.6}\selectfont 超节点只是把NVLink域放大，单卡带宽与算力不变，跨机仍走50GB/s网卡。\par}\endgroup
\item \begingroup\sloppy{\fontsize{10}{14.6}\selectfont Ring All-Reduce走2(p-1)轮，耗时=轮数$\times$每轮延迟再加传输时间。\par}\endgroup
\item \begingroup\sloppy{\fontsize{10}{14.6}\selectfont 小消息的All-Reduce耗时几乎全是轮数$\times$延迟，带宽项仅40ns，加带宽没用。\par}\endgroup
\item \begingroup\sloppy{\fontsize{10}{14.6}\selectfont 32B的TP8一步5.15ms，域内TP16仍5.64ms：KV头只有8个，切不动。\par}\endgroup
\item \begingroup\sloppy{\fontsize{10}{14.6}\selectfont 405B换域内TP16，一步从39.13降到27.43ms，TP才缩短延迟。\par}\endgroup
\item \begingroup\sloppy{\fontsize{10}{14.6}\selectfont MoE靠容量取胜：EP组越大每卡专家越少、会话越多，一步时间几乎固定。\par}\endgroup
\end{itemize}
\meta{公式}{$T_{\text{AR}}=2(p-1)\times(\alpha+m/(p\cdot\beta))$；$\text{可用带宽}=\min(\text{链路带宽},\ u\cdot q/L)$}
\meta{数字}{TP8/跨机TP16/域内TP16：5.15 / 6.67 / 5.64 ms；TP16 追平需每次 AR $\leq$20.9 $\mu$s（域内 24.7、跨机 32.74）；405B 39.13$\rightarrow$27.43 ms；DeepSeek-V3 的 S=16/64/256：34/87/101 会话、1436/3504/3996 tok/s；64 卡分 8 台 2248、重叠 2754，域内重叠反降到 2815；单会话 5.96 ms、168 tok/s；环面 336 对交换 112 $\mu$s；用满 50 GB/s 需 1469 个在途请求}
\meta{易错}{超节点不是更快的 GPU：单卡显存带宽不变，KV 头不够分时 TP16 在域内也比 TP8 慢。}
\phantomsection\addcontentsline{toc}{section}{116\quad SLO 与容量规划}
\lec{116}{SLO 与容量规划}
\onel{先写SLO，再由TPOT上限与Little定律算单实例请求率，按余量算卡数。}
\begin{itemize}[leftmargin=1.35em,itemsep=2.2pt,topsep=1pt,parsep=0pt]
\item \begingroup\sloppy{\fontsize{10}{14.6}\selectfont SLO写成TTFT与TPOT的P99；TPOT上限决定同时能decode多少请求。\par}\endgroup
\item \begingroup\sloppy{\fontsize{10}{14.6}\selectfont decode只能用剩余1-f的时间，故TPOT=t(b)÷(1-f)。\par}\endgroup
\item \begingroup\sloppy{\fontsize{10}{14.6}\selectfont 8B单卡每实例15.3请求/s，同时136个在decode，prefill占39\%。\par}\endgroup
\item \begingroup\sloppy{\fontsize{10}{14.6}\selectfont 实例数=$\lceil$总请求率÷(单实例请求率$\times$使用率)$\rceil$，留30\%余量守P99。\par}\endgroup
\item \begingroup\sloppy{\fontsize{10}{14.6}\selectfont 启动积压$\lambda$$\cdot$T\_s，排空=$\lambda$$\cdot$T\_s÷($\mu$-$\lambda$)，靠余量不靠服务率。\par}\endgroup
\item \begingroup\sloppy{\fontsize{10}{14.6}\selectfont 成本分母是合格请求：四分之一合格时每千请求从0.49涨到1.95元。\par}\endgroup
\end{itemize}
\meta{公式}{$b=\lambda\times\text{输出长度}\times\text{TPOT}$；$\text{TPOT}=t(b)\div(1-f)$；实例数 $=\lceil\text{QPS}\div(\lambda\times\text{使用率})\rceil$；排空 $=\lambda\cdot T_s\div(\mu-\lambda)$}
\meta{数字}{15.3 请求/s、136 在 decode、TPOT 29.7 ms；TPOT 15 ms$\rightarrow$9.7/s、8 张卡；输入 4000$\rightarrow$4.0/s、18 张；200 请求/s$\rightarrow$19 张；5 张 H100 约 69 美元/百万请求；排空 43.3 / 62.2 / 68.8 s；$\mu$ 至少 4.20；捕获执行图 50,000 步持平；迁移 1.66 / 6.76 s；重配置 11.6 / 9.7 / 9.0 s}
\meta{易错}{排空靠的是余量 $\mu$-$\lambda$ 而不是服务率：PD 只比到达率高 8\%，效率掉一成就到第 68.8 s 才排空。}
\chapter{手写迷你推理引擎}
{\footnotesize\sffamily\color{grayt}收录第 117--125 讲}\par\vspace{-4pt}
\phantomsection\addcontentsline{toc}{section}{117\quad 一个推理引擎的组成}
\lec{117}{一个推理引擎的组成}
\onel{引擎骨架是一个循环，按模型执行、状态、调度、内存管理、加速五部分拆开。}
\begin{itemize}[leftmargin=1.35em,itemsep=2.2pt,topsep=1pt,parsep=0pt]
\item \begingroup\sloppy{\fontsize{10}{14.6}\selectfont 请求走八步：分词、进队列、调度、准备KV位置、前向、采样、更新状态，循环到结束。\par}\endgroup
\item \begingroup\sloppy{\fontsize{10}{14.6}\selectfont 五个部分：模型执行、状态、调度、内存管理、加速手段，其余都只是组织工作。\par}\endgroup
\item \begingroup\sloppy{\fontsize{10}{14.6}\selectfont 优化各落在循环某一步：批处理改调度、分页改KV分配、量化改权重。\par}\endgroup
\item \begingroup\sloppy{\fontsize{10}{14.6}\selectfont 六版按依赖加零件，每版都能跑，且先过逐token对齐再测速度。\par}\endgroup
\item \begingroup\sloppy{\fontsize{10}{14.6}\selectfont 选Qwen3-0.6B：28层、1024维、GQA16对8，权重1.5GB。\par}\endgroup
\item \begingroup\sloppy{\fontsize{10}{14.6}\selectfont 硬标准是4条提示各贪心生成32个token全一致，只有v4量化允许不同。\par}\endgroup
\end{itemize}
\meta{数字}{v0--v5 共 6 版，最后一版约 400 行 Python；28 层、1024 维、16 查询头/8 KV 头、权重 1.5 GB；4 条提示$\times$32 token 逐一对齐；CPU FP32、Apple M2 16 GB}
\meta{易错}{用 FP32 不是为快而是为了对齐：掩码差一位或位置从 1 开始数，输出照样通顺却全错。}
\phantomsection\addcontentsline{toc}{section}{118\quad 从权重文件到第一个token}
\lec{118}{从权重文件到第一个token}
\onel{v0用四个文件把1.5GB权重变成前向，59行算出Paris。}
\begin{itemize}[leftmargin=1.35em,itemsep=2.2pt,topsep=1pt,parsep=0pt]
\item \begingroup\sloppy{\fontsize{10}{14.6}\selectfont 权重文件共311个张量、752M参数，BF16存储，头部是一段JSON。\par}\endgroup
\item \begingroup\sloppy{\fontsize{10}{14.6}\selectfont q输出2048维（16头$\times$128），k、v只1024维（8头$\times$128），即GQA。\par}\endgroup
\item \begingroup\sloppy{\fontsize{10}{14.6}\selectfont 头维$\times$头数=2048大于隐藏维1024，头维度不等于d÷头数。\par}\endgroup
\item \begingroup\sloppy{\fontsize{10}{14.6}\selectfont q\_norm、k\_norm各128个数，在RoPE之前对每个头做RMSNorm。\par}\endgroup
\item \begingroup\sloppy{\fontsize{10}{14.6}\selectfont lm\_head与嵌入表同权重，文件存了两份，752M里1.56亿重复。\par}\endgroup
\item \begingroup\sloppy{\fontsize{10}{14.6}\selectfont 四个易错：RoPE前后两半配对、repeat\_interleave、位置从0起。\par}\endgroup
\end{itemize}
\meta{数字}{311 张量、752 M 参数；词表 151936、d=1024、FFN 3072；config 33 行、weights 13 行、model 59 行、engine 32 行；32 token 6.95 s、每 token 217 ms；199 token 提示第 1 步 419 ms、第 64 步 724 ms、共 34 s}
\meta{易错}{v0 慢在重算：总计算约 64$\times$230$\approx$15,000 个 token 的前向，真正新算的只有 263 个。}
\phantomsection\addcontentsline{toc}{section}{119\quad KV Cache}
\lec{119}{KV Cache}
\onel{只缓存每层K、V，新token只算自己的qkv，生成从34s降到4.55s。}
\begin{itemize}[leftmargin=1.35em,itemsep=2.2pt,topsep=1pt,parsep=0pt]
\item \begingroup\sloppy{\fontsize{10}{14.6}\selectfont 只有注意力跨位置，而k、v一旦算出就不再变，所以只需缓存每层的K、V。\par}\endgroup
\item \begingroup\sloppy{\fontsize{10}{14.6}\selectfont 缓存形状 [层数, 最大长度, KV头数, 头维度]，一个请求一份连续内存。\par}\endgroup
\item \begingroup\sloppy{\fontsize{10}{14.6}\selectfont FP32下每token 224KiB，2048token请求448MiB，BF16减半。\par}\endgroup
\item \begingroup\sloppy{\fontsize{10}{14.6}\selectfont 注意力多两行：先写新token的K、V，再读回全部；存RoPE之后的k。\par}\endgroup
\item \begingroup\sloppy{\fontsize{10}{14.6}\selectfont 掩码从T$\times$T下三角改成「查询位置p只看键位置$\leq$p」，输出只算最后位置。\par}\endgroup
\item \begingroup\sloppy{\fontsize{10}{14.6}\selectfont decode每步要读一遍全部3GB权重，所以63ms远慢于prefill。\par}\endgroup
\end{itemize}
\meta{数字}{v0 419 / 417 / 724 ms、共 34.12 s；v1 357 / 63 / 66 ms、共 4.55 s（1/7.5）；每 token KV 224 KiB、2048 token 448 MiB；kv\_cache.py 18 行；model.py +10/-7 行}
\meta{易错}{缓存按最大长度一次分配，请求提前结束时后面没用到的空间照样占着，这正是下一讲的毛病。}
\phantomsection\addcontentsline{toc}{section}{120\quad Continuous Batching}
\lec{120}{Continuous Batching}
\onel{每步重定batch，新token拼一维、只有注意力分段，吞吐到1.9倍。}
\begin{itemize}[leftmargin=1.35em,itemsep=2.2pt,topsep=1pt,parsep=0pt]
\item \begingroup\sloppy{\fontsize{10}{14.6}\selectfont 引擎拆成请求状态、调度器与step()三件，状态里关键是已算位置数。\par}\endgroup
\item \begingroup\sloppy{\fontsize{10}{14.6}\selectfont 每步先让运行中的请求各decode一个token，有空位再放新请求做prefill。\par}\endgroup
\item \begingroup\sloppy{\fontsize{10}{14.6}\selectfont 模型侧把所有新token拼成一维并记下每段，线性层一次算完。\par}\endgroup
\item \begingroup\sloppy{\fontsize{10}{14.6}\selectfont 只有注意力必须按请求分段，每个请求只和自己的上下文做注意力。\par}\endgroup
\item \begingroup\sloppy{\fontsize{10}{14.6}\selectfont 上限是每步请求数与每步token数，static开关用来对比静态批。\par}\endgroup
\item \begingroup\sloppy{\fontsize{10}{14.6}\selectfont 同样792个token：逐个792步、静态批192步、continuous只160步。\par}\endgroup
\end{itemize}
\meta{数字}{792 个 token（8--96 不等）；792 / 192 / 160 步；14.5 / 22.4 / 27.5 token/s；每步 69 ms $\rightarrow$ 约 180 ms（请求 5 倍、耗时 2.6 倍）；槽空置 48\% 对 38\%}
\meta{易错}{瓶颈接着变成按最大长度预留整份 KV：声明 512 实际只生成 50，九成空间空着却挡住后面的请求。}
\phantomsection\addcontentsline{toc}{section}{121\quad 分页KV 管理}
\lec{121}{分页KV 管理}
\onel{共享块池加块表取代按最大长度预留，块不够就抢占最晚请求重算。}
\begin{itemize}[leftmargin=1.35em,itemsep=2.2pt,topsep=1pt,parsep=0pt]
\item \begingroup\sloppy{\fontsize{10}{14.6}\selectfont v2 按提示加最大输出整块预留，能跑几个请求由最大长度而非实际用量决定。\par}\endgroup
\item \begingroup\sloppy{\fontsize{10}{14.6}\selectfont 块池改成 [层, 块, 16, KV 头, 头维] 一整块，K、V 各一份。\par}\endgroup
\item \begingroup\sloppy{\fontsize{10}{14.6}\selectfont 位置 p 在 table[p // 16] 块的第 p \% 16 格。\par}\endgroup
\item \begingroup\sloppy{\fontsize{10}{14.6}\selectfont SeqKV 接口同 v1 的 KVCache，按块表换算，model.py 不用改。\par}\endgroup
\item \begingroup\sloppy{\fontsize{10}{14.6}\selectfont 调度器按空闲块数放行，拿不到块就抢占最晚进来的请求并还块。\par}\endgroup
\item \begingroup\sloppy{\fontsize{10}{14.6}\selectfont 最先进来的请求永不被抢占，KV 只是缓存，重算结果完全相同。\par}\endgroup
\end{itemize}
\meta{公式}{位置 $p\to$ 块号 $\text{table}[\lfloor p/16\rfloor]$，块内偏移 $p\bmod 16$；块池 $\text{shape}=(28,\text{块数},16,8,128)$}
\meta{数字}{FP32 下单块 16 $\times$ 224 KiB = 3.5 MiB；压测 20 块（320 位置）共 70 MiB；KV 只有 320 位置时 v2 17.0 token/s、240 步、每步 3.3 个请求，v3 22.8 token/s、167 步、每步 4.7 个、抢占 9 次（+34\%）；块池缩到 6 块发生 2 次抢占；4/4 条与 HuggingFace 逐 token 一致}
\meta{易错}{分页的价值只在 KV 空间成为瓶颈时显现；空间足够大时 v2 与 v3 没有区别。}
\phantomsection\addcontentsline{toc}{section}{122\quad 前缀缓存}
\lec{122}{前缀缓存}
\onel{装满且算完的块按链式哈希登记，新请求从头逐块命中即共用。}
\begin{itemize}[leftmargin=1.35em,itemsep=2.2pt,topsep=1pt,parsep=0pt]
\item \begingroup\sloppy{\fontsize{10}{14.6}\selectfont 块 KV 由之前的全部 token 加本块 16 个 token 唯一决定。\par}\endgroup
\item \begingroup\sloppy{\fontsize{10}{14.6}\selectfont 每块哈希都含前一块哈希，链式连到开头：相同即整个前缀相同。\par}\endgroup
\item \begingroup\sloppy{\fontsize{10}{14.6}\selectfont 只登记装满且已算完的块，维护哈希$\rightarrow$块号、块号$\rightarrow$哈希两张表。\par}\endgroup
\item \begingroup\sloppy{\fontsize{10}{14.6}\selectfont 从头逐块查，连续命中才继续、第一次没中就停。\par}\endgroup
\item \begingroup\sloppy{\fontsize{10}{14.6}\selectfont 引用计数让多请求共用块，归零后内容与登记保留，被覆盖才作废。\par}\endgroup
\item \begingroup\sloppy{\fontsize{10}{14.6}\selectfont 含最后一个 token 的块总自己算，否则得不到该位置的打分。\par}\endgroup
\end{itemize}
\meta{公式}{$h_0=\operatorname{hash}(\varnothing,\ \text{第 }0\text{ 块 16 个 token})$；$h_i=\operatorname{hash}(h_{i-1},\ \text{第 }i\text{ 块 16 个 token})$}
\meta{数字}{后 7 个请求提示共 2356 token：关闭时命中 0、prefill 2356、5.5 s；打开时命中 2240（7 $\times$ 20 块 $\times$ 16）、只 prefill 116、2.4 s；系统提示约 330 token；两种输出完全一致 True}
\meta{易错}{时间只少一半多而非 95\%：prefill 省了，但每请求仍要 decode 8 个 token，decode 仍读完整 330 多个位置的 K、V。}
\phantomsection\addcontentsline{toc}{section}{123\quad 权重量化}
\lec{123}{权重量化}
\onel{每通道一个 scale 的 INT8 量化，显存降到 1/4 但更慢。}
\begin{itemize}[leftmargin=1.35em,itemsep=2.2pt,topsep=1pt,parsep=0pt]
\item \begingroup\sloppy{\fontsize{10}{14.6}\selectfont 只量化每层 7 个线性层，嵌入表、输出层、RMSNorm 都保持原样。\par}\endgroup
\item \begingroup\sloppy{\fontsize{10}{14.6}\selectfont 对称量化取 -127 到 127，每行 scale = 最大绝对值 ÷ 127。\par}\endgroup
\item \begingroup\sloppy{\fontsize{10}{14.6}\selectfont 按行取 scale 使缩放可后置：y = (x$\cdot$qᵀ) $\odot$ s，不必还原矩阵。\par}\endgroup
\item \begingroup\sloppy{\fontsize{10}{14.6}\selectfont 线性层 1680 $\rightarrow$ 421 MiB，整模型 2.8 GiB $\rightarrow$ 约 1.6 GiB。\par}\endgroup
\item \begingroup\sloppy{\fontsize{10}{14.6}\selectfont 没有 int8 kernel 时每次乘前要转回 FP32，CPU 上反而慢 20\%。\par}\endgroup
\item \begingroup\sloppy{\fontsize{10}{14.6}\selectfont 提速要靠融合 kernel 直接读 int8，显存只读 FP32 的四分之一。\par}\endgroup
\end{itemize}
\meta{公式}{$s=\operatorname{amax}(\text{行})/127$；$q=\operatorname{clamp}(\operatorname{round}(W/s),-127,127)$；$y=(x\cdot q^{\top})\odot s$}
\meta{数字}{线性层 1680 MiB $\rightarrow$ 421 MiB（3.99 倍），多出的 0.01 是 scale 开销；128 步最高分 token 全同、平均 KL 散度 0.0011；单请求 64 token：FP32 4.27 s、INT8 5.11 s；4 条检查 3 条逐 token 一致；省下约 1.2 GiB 可多放约 360 个块}
\meta{易错}{「逐 token 一致」对量化太苛刻：应改成读同一串 token、比每步打分与 KL 散度。}
\phantomsection\addcontentsline{toc}{section}{124\quad 投机解码}
\lec{124}{投机解码}
\onel{n-gram 找草稿、一次前向验证，贪心采样下输出与不投机相同。}
\begin{itemize}[leftmargin=1.35em,itemsep=2.2pt,topsep=1pt,parsep=0pt]
\item \begingroup\sloppy{\fontsize{10}{14.6}\selectfont 起草取结尾 n 个 token（3 退到 1），从近往远找，返回其后最多 k 个。\par}\endgroup
\item \begingroup\sloppy{\fontsize{10}{14.6}\selectfont 验证把最后一个真实 token 与草稿连续送入，逐个比对贪心预测。\par}\endgroup
\item \begingroup\sloppy{\fontsize{10}{14.6}\selectfont 接受到第一个不一致处，再加该位置的模型预测；全错也有 p$_0$。\par}\endgroup
\item \begingroup\sloppy{\fontsize{10}{14.6}\selectfont 回滚只推进 num\_computed，被拒位置的 KV 下一步被覆盖。\par}\endgroup
\item \begingroup\sloppy{\fontsize{10}{14.6}\selectfont 起草是纯 Python 列表比较，比一次前向便宜几个数量级。\par}\endgroup
\item \begingroup\sloppy{\fontsize{10}{14.6}\selectfont forward 多 out\_idx，调度器按 1 + 草稿数要块与算预算。\par}\endgroup
\end{itemize}
\meta{公式}{送入 $t,d_1,\dots,d_k$ 得 $p_0,\dots,p_k$；$p_{i-1}=d_i$ 则接受 $d_i$，接受到第一个不一致处再加上该位置的 $p_j$}
\meta{数字}{改写代码：128 token 从 127 步降到 34 步，接受 93 个草稿，平均每次前向 3.66 个 token，9.4 s $\rightarrow$ 6.6 s；自由生成 53 步、2.37 个，9.2 s $\rightarrow$ 8.0 s；两种情况输出都一致}
\meta{易错}{步数降到约四分之一、时间只降到七成：CPU 上多算 4 个位置不免费，batch 大时草稿更贵。}
\phantomsection\addcontentsline{toc}{section}{125\quad 时间花在哪里}
\lec{125}{时间花在哪里}
\onel{CPU 上 76\% 时间在线性层、24\% 在注意力，Python 调度可忽略。}
\begin{itemize}[leftmargin=1.35em,itemsep=2.2pt,topsep=1pt,parsep=0pt]
\item \begingroup\sloppy{\fontsize{10}{14.6}\selectfont 六个版本约 400 行，v0 119 $\rightarrow$ v5 407 行，最大增长是 v3 块管理器。\par}\endgroup
\item \begingroup\sloppy{\fontsize{10}{14.6}\selectfont v1 的 KV Cache 把同任务从 34.1 s 降到 4.55 s。\par}\endgroup
\item \begingroup\sloppy{\fontsize{10}{14.6}\selectfont 每版收益都有前提：负载不匹配时，优化甚至会成为负担。\par}\endgroup
\item \begingroup\sloppy{\fontsize{10}{14.6}\selectfont 8 请求各 64 token 共 13.1 s：注意力 24\%、线性层 76\%。\par}\endgroup
\item \begingroup\sloppy{\fontsize{10}{14.6}\selectfont 注意力偏高因逐请求 Python 循环，且 read 每层每步复制块。\par}\endgroup
\item \begingroup\sloppy{\fontsize{10}{14.6}\selectfont 单请求 64 token 只比 HuggingFace 慢 0.6 s。\par}\endgroup
\end{itemize}
\meta{数字}{13.1 s / 64 步；注意力 3.1 s、线性层等 10.0 s、Python 26 ms；v2 27.5 token/s 为逐个跑的 1.9 倍；v3 再 +34\%；前缀缓存省 95\% 重复 prefill；v5 每次前向 3.66 个 token；v1 快 7.5 倍、decode 每步约 65 ms}
\meta{易错}{GPU 上一步 decode 只需几毫秒，同样的几十微秒 Python 开销就变得可观：vLLM v0.6.0 剖析里曾占一步的六成。}
\chapter{解读nano-vllm}
{\footnotesize\sffamily\color{grayt}收录第 126--128 讲}\par\vspace{-4pt}
\phantomsection\addcontentsline{toc}{section}{126\quad nano-vllm 的整体结构}
\lec{126}{nano-vllm 的整体结构}
\onel{拆成调度器、执行器、引擎三段，prefill 优先、一步只做一种。}
\begin{itemize}[leftmargin=1.35em,itemsep=2.2pt,topsep=1pt,parsep=0pt]
\item \begingroup\sloppy{\fontsize{10}{14.6}\selectfont 引擎在主进程 step()：问调度器跑谁、执行器跑模型采样、回填状态。\par}\endgroup
\item \begingroup\sloppy{\fontsize{10}{14.6}\selectfont 拆开是因为张量并行：每卡一份执行器，rank 0 广播名单。\par}\endgroup
\item \begingroup\sloppy{\fontsize{10}{14.6}\selectfont 一步要么全 prefill 要么全 decode，两种不混在一起。\par}\endgroup
\item \begingroup\sloppy{\fontsize{10}{14.6}\selectfont 这样 decode 步形状规整，才能按 batch 分档录 CUDA Graph。\par}\endgroup
\item \begingroup\sloppy{\fontsize{10}{14.6}\selectfont 比 miniinfer 多了分块 prefill（只切等待队首）和块大小 256。\par}\endgroup
\item \begingroup\sloppy{\fontsize{10}{14.6}\selectfont 块管理几乎一样：链式哈希、引用计数、抢最晚者、倒着还块。\par}\endgroup
\end{itemize}
\meta{数字}{nano-vllm 1195 行 vs miniinfer v5 407 行；调度一侧 247 行 vs 180 行，多出的约 800 行在执行侧（执行器 256、模型 479、注意力 66）；默认每步 16384 token、512 请求、上下文 4096、90\% 显存、块 256；README 测 RTX 4070 上 vLLM 1361.84、nano-vllm 1434.13 token/s}
\meta{易错}{块大小 256 让前缀缓存粒度变粗：约 330 token 的共用系统提示只能命中 1 块，miniinfer 能命中 20 块。}
\phantomsection\addcontentsline{toc}{section}{127\quad nano-vllm 多做了什么}
\lec{127}{nano-vllm 多做了什么}
\onel{执行侧补 GPU 写法：预热量显存分 KV、变长注意力、分档录图。}
\begin{itemize}[leftmargin=1.35em,itemsep=2.2pt,topsep=1pt,parsep=0pt]
\item \begingroup\sloppy{\fontsize{10}{14.6}\selectfont 启动按最坏 prefill 预热一遍记下激活峰值，剩余显存全切成 KV 块。\par}\endgroup
\item \begingroup\sloppy{\fontsize{10}{14.6}\selectfont 一步的请求压成 slot\_mapping、块表等一维数组再进 kernel。\par}\endgroup
\item \begingroup\sloppy{\fontsize{10}{14.6}\selectfont 写 KV 用 Triton kernel 跳过 -1，否则覆盖最后一块的 KV。\par}\endgroup
\item \begingroup\sloppy{\fontsize{10}{14.6}\selectfont prefill 走变长 FlashAttention，decode 按块表读 KV。\par}\endgroup
\item \begingroup\sloppy{\fontsize{10}{14.6}\selectfont decode 按 batch 分档录图，补齐行靠槽位 -1 与长度 0。\par}\endgroup
\item \begingroup\sloppy{\fontsize{10}{14.6}\selectfont 张量并行按列切与行切配对，每层两次 All-Reduce。\par}\endgroup
\end{itemize}
\meta{公式}{槽位 $\text{slot}=$ 块号 $\times$ 块大小 $+$ 块内偏移；$\operatorname{argmax}_i\,p_i/E_i$（$E_i\sim\operatorname{Exp}(1)$ 独立）$\iff$ 按概率 $p_i$ 抽出 $i$}
\meta{数字}{36 张图（1、2、4、8，再每 16 一档到 512）；37 个请求用 48 档、补 11 行、多算 23\% 的行；16 查询头 + 2 $\times$ 8 KV 头 $\rightarrow$ 合并矩阵 4096 行，双卡各 2048 行（8 查询头、4 K 头、4 V 头）；decode 一步 58 次集合通信（28 层 $\times$ 2 + 2）}
\meta{易错}{块大小 256 是 FlashAttention 分页接口的硬要求（assert 块大小 \% 256 == 0），代价是末块平均浪费约 128 个位置。}
\phantomsection\addcontentsline{toc}{section}{128\quad 动手改 nano-vllm}
\lec{128}{动手改 nano-vllm}
\onel{投机移植的关键是草稿单列、验证复用带前缀缓存的 prefill。}
\begin{itemize}[leftmargin=1.35em,itemsep=2.2pt,topsep=1pt,parsep=0pt]
\item \begingroup\sloppy{\fontsize{10}{14.6}\selectfont 草稿不放进 token\_ids：它牵动 num\_tokens、生成数与前缀哈希。\par}\endgroup
\item \begingroup\sloppy{\fontsize{10}{14.6}\selectfont 验证复用带前缀缓存的 prefill，代价是不能用 CUDA Graph。\par}\endgroup
\item \begingroup\sloppy{\fontsize{10}{14.6}\selectfont Context 加 all\_logits，输出层对每个送入位置都算打分。\par}\endgroup
\item \begingroup\sloppy{\fontsize{10}{14.6}\selectfont 草稿确定，接受规则化为：照常采样 x，x 等于草稿就接受。\par}\endgroup
\item \begingroup\sloppy{\fontsize{10}{14.6}\selectfont 块不够先放弃草稿再抢占；验证后用 trim 退还多占的块。\par}\endgroup
\item \begingroup\sloppy{\fontsize{10}{14.6}\selectfont 回滚只推进 num\_cached\_tokens，且必须先追加 token 再登记哈希。\par}\endgroup
\end{itemize}
\meta{数字}{nano-vllm 新增 132 行、改 8 个文件（含新增 spec.py 27 行）；调度器 +42/-11、执行器 +38/-3、块管理器 +17/-4；miniinfer 只新增 1 个文件、改 4 个；温度取 10$^-$⁵，三请求各 256 token}
\meta{易错}{两处不改就出错却最易漏：序列化要带上 draft（否则张量并行各卡形状不一致），run\_model 判断用不用 CUDA Graph 要看 Context。}
\chapter{主流推理引擎}
{\footnotesize\sffamily\color{grayt}收录第 129--133 讲}\par\vspace{-4pt}
\phantomsection\addcontentsline{toc}{section}{129\quad vLLM 的架构}
\lec{129}{vLLM 的架构}
\onel{V1 拆三类进程，调度器只维护「已算」「应有」两个数，每步补差额。}
\begin{itemize}[leftmargin=1.35em,itemsep=2.2pt,topsep=1pt,parsep=0pt]
\item \begingroup\sloppy{\fontsize{10}{14.6}\selectfont API 服务器分词反分词、EngineCore 调度、Worker 每卡一个。\par}\endgroup
\item \begingroup\sloppy{\fontsize{10}{14.6}\selectfont 调度器没有 prefill 和 decode 阶段，每请求只有「已算」「应有」两个数。\par}\endgroup
\item \begingroup\sloppy{\fontsize{10}{14.6}\selectfont 同一个差额统一 prefill、分块 prefill、前缀缓存、decode 和投机。\par}\endgroup
\item \begingroup\sloppy{\fontsize{10}{14.6}\selectfont 块哈希是 SHA-256 链式哈希，多带 LoRA 编号、图片哈希、缓存盐。\par}\endgroup
\item \begingroup\sloppy{\fontsize{10}{14.6}\selectfont 异步调度先占位再算：调度器不知道上一步采出什么就先排下一步。\par}\endgroup
\item \begingroup\sloppy{\fontsize{10}{14.6}\selectfont SchedulerOutput 只发增量：新请求全量、老请求只发新增块号。\par}\endgroup
\end{itemize}
\meta{公式}{$\text{num\_tokens\_with\_spec}=$ 提示 $+$ 已生成 $+$ 草稿；每步差额 $=\text{num\_tokens\_with\_spec}-\text{num\_computed\_tokens}$}
\meta{数字}{vllm/v1 目录 170876 行；调度器 3569 行；nano-vllm 全仓 1197 行；一步预算 1024 token；一步 17 ms 变 max(12, 3+2)=12 ms，吞吐 1.42 倍}
\meta{易错}{long\_prefill\_token\_threshold 默认是 0，表示不限；不分块时一个长提示会把整步的 token 预算吃光。}
\phantomsection\addcontentsline{toc}{section}{130\quad vLLM 的分布式服务}
\lec{130}{vLLM 的分布式服务}
\onel{服务五层：卡间广播、副本间负载均衡、MoE 锁步、connector 接 PD。}
\begin{itemize}[leftmargin=1.35em,itemsep=2.2pt,topsep=1pt,parsep=0pt]
\item \begingroup\sloppy{\fontsize{10}{14.6}\selectfont 副本内靠共享内存广播队列驱动各卡，只让唯一指定的 worker 回结果。\par}\endgroup
\item \begingroup\sloppy{\fontsize{10}{14.6}\selectfont 稠密模型副本彼此独立；MoE 副本每步有 All-to-All，必须锁步、空批次陪跑。\par}\endgroup
\item \begingroup\sloppy{\fontsize{10}{14.6}\selectfont 挑副本看快照、本地计数、在途数和 KV 压力罚分，快照最小间隔 100 ms。\par}\endgroup
\item \begingroup\sloppy{\fontsize{10}{14.6}\selectfont 全体同步用 2 元素张量做一次 All-Reduce，一位报忙、一位报想暂停。\par}\endgroup
\item \begingroup\sloppy{\fontsize{10}{14.6}\selectfont 每 16 步才同步一次是否全体空闲，代价是空闲后多陪跑十几个空批次。\par}\endgroup
\item \begingroup\sloppy{\fontsize{10}{14.6}\selectfont connector 分两侧：调度器问外部还有多少 KV，worker 做加载和保存。\par}\endgroup
\end{itemize}
\meta{公式}{$\text{sync\_dp\_state}$：2 元素张量做一次 SUM All-Reduce，$[0]=1$ 表示有未完成请求（$\text{SUM}>0$ 即或），$[1]=1$ 表示想暂停（$\text{SUM}=\text{dp\_size}$ 全体同意）}
\meta{数字}{两节点 8 卡、4 个副本 $\times$ TP 4；快照间隔 100 ms；dp\_sync\_interval 默认 16；平均排队只看快照 497 ms，加本地计数降到 123 ms}
\meta{易错}{博客和文档会过时：现在只有 MoE 副本才锁步，博客里的 SharedStorageConnector 已被 example\_connector.py 取代。}
\phantomsection\addcontentsline{toc}{section}{131\quad SGLang 的架构}
\lec{131}{SGLang 的架构}
\onel{每卡一个调度器进程，前缀缓存用基数树，调度与执行重叠即零开销调度。}
\begin{itemize}[leftmargin=1.35em,itemsep=2.2pt,topsep=1pt,parsep=0pt]
\item \begingroup\sloppy{\fontsize{10}{14.6}\selectfont 每卡一个调度器进程，调度和执行同进程；rank 0 收请求后广播。\par}\endgroup
\item \begingroup\sloppy{\fontsize{10}{14.6}\selectfont 省掉调度结果每步跨进程发送，代价是每卡重复算一遍调度且结果必须一致。\par}\endgroup
\item \begingroup\sloppy{\fontsize{10}{14.6}\selectfont 节点存 token 段和 KV 槽位，孩子按首 token 索引，停在中间就劈开。\par}\endgroup
\item \begingroup\sloppy{\fontsize{10}{14.6}\selectfont 淘汰从可淘汰叶子按 LRU 弹堆；父节点变叶子后再入堆，公共前缀最后被淘汰。\par}\endgroup
\item \begingroup\sloppy{\fontsize{10}{14.6}\selectfont 两级池：req\_to\_token 逐 token 记槽位，第 0 行留给补出来的空请求。\par}\endgroup
\item \begingroup\sloppy{\fontsize{10}{14.6}\selectfont 零开销调度先发射这一步，再处理上一步结果，i+1 步输入直接引用显存里的输出。\par}\endgroup
\end{itemize}
\meta{公式}{命中 $=$ 沿树逐节点比 $\text{key}$ 的公共前缀；节点 $\text{value}$ 是该段 token 的 KV 槽位，$\text{lock\_ref}>0$ 不可淘汰}
\meta{数字}{GPU 一步 12 ms、CPU 共 5 ms：串行 17 ms（GPU 忙 71\%），重叠 12 ms、吞吐 1.42 倍；块哈希粒度 16--256 token，基数树是 1 token；lpm 在等待队列超 128 时退化}
\meta{易错}{这个提交里默认策略是 fcfs，缓存感知的排序要显式打开；它能提高命中率，但会让前缀罕见的请求等更久。}
\phantomsection\addcontentsline{toc}{section}{132\quad 其他引擎与规模化部署}
\lec{132}{其他引擎与规模化部署}
\onel{选引擎看它押注在哪一层：kernel、量化、硬件覆盖面还是编译器。}
\begin{itemize}[leftmargin=1.35em,itemsep=2.2pt,topsep=1pt,parsep=0pt]
\item \begingroup\sloppy{\fontsize{10}{14.6}\selectfont TensorRT-LLM 押 kernel：强在 NVIDIA 新卡的 FP8、FP4。\par}\endgroup
\item \begingroup\sloppy{\fontsize{10}{14.6}\selectfont LMDeploy：TurboMind 押 4 位权重量化和 KV 量化的 kernel。\par}\endgroup
\item \begingroup\sloppy{\fontsize{10}{14.6}\selectfont llama.cpp 用 GGUF 和 k-quant，目标是本地单用户低并发。\par}\endgroup
\item \begingroup\sloppy{\fontsize{10}{14.6}\selectfont MLC-LLM 押编译器：TVM 为每种硬件生成 kernel，能到手机和 WebGPU。\par}\endgroup
\item \begingroup\sloppy{\fontsize{10}{14.6}\selectfont 大规模 EP 四样拼：DP 注意力、All-to-All 通信库、EPLB、微批重叠。\par}\endgroup
\item \begingroup\sloppy{\fontsize{10}{14.6}\selectfont 编排层是 Dynamo、llm-d 和 SGLang 路由器，缓存感知路由也在这一层。\par}\endgroup
\end{itemize}
\meta{公式}{大规模 $\text{EP}=\text{DP}$ 注意力 $+$ DeepEP 类 All-to-All $+$ EPLB $+$ 微批重叠}
\meta{数字}{TensorRT-LLM 只支持 NVIDIA；LMDeploy 押 4 位权重量化；llama.cpp 十几种 k-quant；MLC-LLM 用 TVM}
\meta{易错}{这些定位是定性的，不是压测排名；四个主流引擎都具备 continuous batching、分页 KV、前缀缓存和多卡并行，差别在押注方向。}
\phantomsection\addcontentsline{toc}{section}{133\quad 怎么读一个大型引擎的源码}
\lec{133}{怎么读一个大型引擎的源码}
\onel{带着小引擎的地图进去，只追一个请求的路径；定位功能从配置项搜起，用运行验证。}
\begin{itemize}[leftmargin=1.35em,itemsep=2.2pt,topsep=1pt,parsep=0pt]
\item \begingroup\sloppy{\fontsize{10}{14.6}\selectfont 带小引擎的地图进代码：3569 行调度器里只认跑谁、算几个 token、块不够、何时结束。\par}\endgroup
\item \begingroup\sloppy{\fontsize{10}{14.6}\selectfont 顺着一个请求走而不是顺着目录走：只追分词、入队、调度、送模型、接 token。\par}\endgroup
\item \begingroup\sloppy{\fontsize{10}{14.6}\selectfont 先读设计注释和数据结构，再读函数；用单测和 git log -L 弄懂为什么。\par}\endgroup
\item \begingroup\sloppy{\fontsize{10}{14.6}\selectfont 定位功能从配置项搜起：先读定义处的文档串和默认值，再跳到使用处。\par}\endgroup
\item \begingroup\sloppy{\fontsize{10}{14.6}\selectfont 验证靠最小配置、DEBUG 日志、Profiler 和只改一个参数的对照。\par}\endgroup
\item \begingroup\sloppy{\fontsize{10}{14.6}\selectfont 压测坑：离线吞吐和在线延迟是两回事，随机提示测不出前缀缓存，长度要固定。\par}\endgroup
\end{itemize}
\meta{公式}{定位一个功能：配置项定义（文档字符串与默认值）$\to$ 传递 $\to$ 使用处，读两头、略中间}
\meta{数字}{vLLM 17 万行；long\_prefill\_token\_threshold 出现在 3 个文件共 21 处；第五章从 119 行 v0 到 400 行 miniinfer、1200 行 nano-vllm}
\meta{易错}{不控制输出长度时不同引擎在不同位置遇到结束符，生成的 token 数不同，吞吐没有可比性；压测要 ignore\_eos 固定长度。}
\chapter{适配原理：把模型跑通}
{\footnotesize\sffamily\color{grayt}收录第 134--141 讲}\par\vspace{-4pt}
\phantomsection\addcontentsline{toc}{section}{134\quad 换一块卡要重做什么}
\lec{134}{换一块卡要重做什么}
\onel{换卡只重做九层里接口下方的实现，nano-vllm 里只有 4.4\% 要改。}
\begin{itemize}[leftmargin=1.35em,itemsep=2.2pt,topsep=1pt,parsep=0pt]
\item \begingroup\sloppy{\fontsize{10}{14.6}\selectfont 九层：模型、引擎、框架、图编译、算子库、kernel 语言、通信库、运行时、工具。\par}\endgroup
\item \begingroup\sloppy{\fontsize{10}{14.6}\selectfont 模型、引擎、框架接口、图表示在接口上方，换卡不用改，工作量在下方实现。\par}\endgroup
\item \begingroup\sloppy{\fontsize{10}{14.6}\selectfont 1195 行里 52 行与 NVIDIA 绑定，占 4.4\%，其中设备运行时 23 行。\par}\endgroup
\item \begingroup\sloppy{\fontsize{10}{14.6}\selectfont 最重的是那 3 行 flash\_attn：要么用厂商移植版，要么自己手写一遍。\par}\endgroup
\item \begingroup\sloppy{\fontsize{10}{14.6}\selectfont 兼容路线沿用 1 层、翻译 7 层、重写 1 层；NPU 路线 1/2/6。\par}\endgroup
\item \begingroup\sloppy{\fontsize{10}{14.6}\selectfont 顺序是跑通、跑对、跑快，不能颠倒：没有基线就分不清错误是优化引入的还是本来就有。\par}\endgroup
\end{itemize}
\meta{公式}{拐点（$\text{FLOP}/\text{字节}$）$=\frac{\text{峰值算力}}{\text{带宽}}$；decode 一步线性层的算术强度 $\approx M$}
\meta{数字}{1195 行中 52 行（4.4\%）；兼容路线 1/7/1；NPU 路线 1/2/6；MPS 上少一次逐层同步 decode 每步 55 ms 降到 49 ms}
\meta{易错}{52 行改动量很小、背后的工作量很大：适配周期由这些行背后的库在新硬件上成熟到什么程度决定。}
\phantomsection\addcontentsline{toc}{section}{135\quad 国产与NVIDIA 硬件差异}
\lec{135}{国产与NVIDIA 硬件差异}
\onel{换卡要重算线程组宽度、指令形状、片上容量三组常数，Roofline 拐点不变。}
\begin{itemize}[leftmargin=1.35em,itemsep=2.2pt,topsep=1pt,parsep=0pt]
\item \begingroup\sloppy{\fontsize{10}{14.6}\selectfont GPGPU 用 SIMT，线程自己搬数据，靠很多小块占满全部计算单元。\par}\endgroup
\item \begingroup\sloppy{\fontsize{10}{14.6}\selectfont NPU 放弃 SIMT：标量、向量、矩阵单元分工，搬运单元显式搬数据、排流水。\par}\endgroup
\item \begingroup\sloppy{\fontsize{10}{14.6}\selectfont 宽度 32 变 64：跨通道归约步数从 5 变 6，掩码要从 32 位换成 64 位。\par}\endgroup
\item \begingroup\sloppy{\fontsize{10}{14.6}\selectfont 16$\times$16 输出 256 个元素：32 通道每通道 8 个累加器，64 通道只有 4 个。\par}\endgroup
\item \begingroup\sloppy{\fontsize{10}{14.6}\selectfont 片上 228 KB 降到 64 KB，注意力的 K/V 块从 304 行掉到 85 行。\par}\endgroup
\item \begingroup\sloppy{\fontsize{10}{14.6}\selectfont Roofline 拐点 = 算力 ÷ 带宽，三块卡 295、245、250，差别很小。\par}\endgroup
\end{itemize}
\meta{公式}{拐点 $=\frac{\text{峰值算力}}{\text{带宽}}$；算术强度 $100$ 时可达 $\min(989,100\times 3.35)=335\ \text{TFLOPS}$（峰值的 $33.9\%$）}
\meta{数字}{H100 989 TFLOPS、3.35 TB/s、132 CU、32 宽、228 KB；GPGPU 示例 1300/5.3/304/64/64；NPU 示例 400/1.6/24/0/512；M=16 只跑到峰值 5.4\%--6.5\%}
\meta{易错}{参数表峰值有两个坑：稀疏算力是稠密的两倍，不同数据类型峰值依次翻倍，比两块卡要确认同类型、同稠密值。}
\phantomsection\addcontentsline{toc}{section}{136\quad 编程模型与CUDA 迁移}
\lec{136}{编程模型与CUDA 迁移}
\onel{CUDA 锁定不在语言，而在算子库、kernel 里的硬件常数、框架外的工具链。}
\begin{itemize}[leftmargin=1.35em,itemsep=2.2pt,topsep=1pt,parsep=0pt]
\item \begingroup\sloppy{\fontsize{10}{14.6}\selectfont 锁定在三处：调好参数的算子库、写进 kernel 的硬件常数、框架外的工具链。\par}\endgroup
\item \begingroup\sloppy{\fontsize{10}{14.6}\selectfont 硬件常数常是裸字面量：32、228 KB、m16n8k16，散在分块逻辑里。\par}\endgroup
\item \begingroup\sloppy{\fontsize{10}{14.6}\selectfont HIPIFY 有两版：perl 版按正则替换，clang 版用 C++ 前端解析、更准。\par}\endgroup
\item \begingroup\sloppy{\fontsize{10}{14.6}\selectfont 写死 \% 32 的代码工具不动，64 通道的卡上要改用 warpSize。\par}\endgroup
\item \begingroup\sloppy{\fontsize{10}{14.6}\selectfont 归约 32 通道 5 步、64 通道 6 步；表决掩码必须 64 位。\par}\endgroup
\item \begingroup\sloppy{\fontsize{10}{14.6}\selectfont 重写路线没有 threadIdx：声明缓冲、搬运、计算、搬回，再排双缓冲流水。\par}\endgroup
\end{itemize}
\meta{公式}{归约步数 $=\log_2 w$：$w=32$ 时 5 步、$w=64$ 时 6 步；表决掩码需要 $w$ 位，$\text{uint32\_t}$ 只够 32 条通道}
\meta{数字}{16 行样例：自动 6、人工确认 7、必须重写 3；海光插件 4250 行里架构内建 11 处、架构宏 4 处、写死 64 共 3 处}
\meta{易错}{和架构绑定的只有十几处，却全落在最热的小 M 矩阵乘 kernel 里；而且这类错误是静默的：编译通过、不报错、输出有值但是错的。}
\phantomsection\addcontentsline{toc}{section}{137\quad 框架怎么接入新硬件}
\lec{137}{框架怎么接入新硬件}
\onel{接新卡＝在算子分发表里填一列，要么沿用 CUDA 键，要么注册新设备。}
\begin{itemize}[leftmargin=1.35em,itemsep=2.2pt,topsep=1pt,parsep=0pt]
\item \begingroup\sloppy{\fontsize{10}{14.6}\selectfont 分发键三类：设备键（CPU、CUDA、PrivateUse1）、功能键、复合键。\par}\endgroup
\item \begingroup\sloppy{\fontsize{10}{14.6}\selectfont 沿用 CUDA 键：用 HIP 重编 PyTorch，torch.cuda 接口照旧。\par}\endgroup
\item \begingroup\sloppy{\fontsize{10}{14.6}\selectfont PrivateUse1 路线：算子、分配器、流、事件、拷贝都要自己实现。\par}\endgroup
\item \begingroup\sloppy{\fontsize{10}{14.6}\selectfont rms\_norm 设备键全空却能跑：注册在复合键，拆成 pow 等基础算子。\par}\endgroup
\item \begingroup\sloppy{\fontsize{10}{14.6}\selectfont 缺算子四条岔路：报错、回退 CPU、用现有算子组合、自己写 kernel。\par}\endgroup
\item \begingroup\sloppy{\fontsize{10}{14.6}\selectfont 自定义算子要配 fake/meta 实现，并正确声明 mutates\_args。\par}\endgroup
\end{itemize}
\meta{公式}{$x\,\text{@}\,w\to\text{aten::mm}\to$ 分发器查表（设备键 $+$ 功能键）$\to$ 该键上的实现 $\to$ 厂商算子库 $\to$ 运行时发射 kernel}
\meta{数字}{aten::mm 在 CPU、MPS、Meta、SparseCPU 各一行；Qwen3 前向 31 种 aten 算子、设备上真执行 23 种；FP8 矩阵乘回退 CPU 22 ms、组合 1.5 ms}
\meta{易错}{原地写 KV Cache 的算子若声明成不改输入，图编译会认为它无副作用，把写入优化掉或重排，KV 写丢。}
\phantomsection\addcontentsline{toc}{section}{138\quad 图模式与编译}
\lec{138}{图模式与编译}
\onel{图中断决定有没有图模式：去掉每层的 int 取值，27 张图降到 2 张。}
\begin{itemize}[leftmargin=1.35em,itemsep=2.2pt,topsep=1pt,parsep=0pt]
\item \begingroup\sloppy{\fontsize{10}{14.6}\selectfont 一步 decode 下发 2212 次，CPU 17.7 ms 已超过 GPU。\par}\endgroup
\item \begingroup\sloppy{\fontsize{10}{14.6}\selectfont 整图下发只付一次入口 20 $\mu$s，一步 13.27 ms，快 1.33 倍。\par}\endgroup
\item \begingroup\sloppy{\fontsize{10}{14.6}\selectfont torch.compile 分捕获、变换、生成三段，厂商只接最后一段。\par}\endgroup
\item \begingroup\sloppy{\fontsize{10}{14.6}\selectfont 后端就是一个函数：收 (gm, example\_inputs)，返回可调用对象。\par}\endgroup
\item \begingroup\sloppy{\fontsize{10}{14.6}\selectfont v1 切成 27 张图、24 处中断；port 版只剩 2 张。\par}\endgroup
\item \begingroup\sloppy{\fontsize{10}{14.6}\selectfont 形状先特化、变了再变符号量，四个长度加 decode 只编 3 张图。\par}\endgroup
\end{itemize}
\meta{数字}{2212 次 $\times$ 8 $\mu$s $\approx$ 17.7 ms；MPS 一步 49 ms；下发小 add 5.4 $\mu$s 对 int(tensor) 141.2 $\mu$s（26 倍）}
\meta{易错}{被切成 27 段的前向等于回到逐算子执行：段间不能融合、不能跨段规划片上缓冲、不能一次下发。}
\phantomsection\addcontentsline{toc}{section}{139\quad 精度对齐}
\lec{139}{精度对齐}
\onel{精度目标不是逐位一致，而是每处误差都能归因；逐层按位置取最差才能定位。}
\begin{itemize}[leftmargin=1.35em,itemsep=2.2pt,topsep=1pt,parsep=0pt]
\item \begingroup\sloppy{\fontsize{10}{14.6}\selectfont 三级检查：逐层相对误差、贪心 token 比对、任务级评测，各有盲区。\par}\endgroup
\item \begingroup\sloppy{\fontsize{10}{14.6}\selectfont 逐位一致不成立：累加顺序、中间精度、超越函数实现都不同。\par}\endgroup
\item \begingroup\sloppy{\fontsize{10}{14.6}\selectfont bug 在第 4 层：6544$^2$ 在 FP16 溢出成 inf，整层输出变 0。\par}\endgroup
\item \begingroup\sloppy{\fontsize{10}{14.6}\selectfont BF16 指数 8 位能到 1e38 不溢出；修法是平方和先升 FP32。\par}\endgroup
\item \begingroup\sloppy{\fontsize{10}{14.6}\selectfont 按位置分别算再取最差：整张量 4e-3 时最差位置已是 1.0。\par}\endgroup
\item \begingroup\sloppy{\fontsize{10}{14.6}\selectfont 阈值按精度定：FP32 1e-5、BF16 5e-2、FP16 1e-2。\par}\endgroup
\end{itemize}
\meta{公式}{相对误差 $=\frac{\lVert a-b\rVert}{\lVert b\rVert}$；$\text{RMSNorm}(x)=\frac{x}{\sqrt{\operatorname{mean}(x^2)+\varepsilon}}\odot w$}
\meta{数字}{MPS FP32 1e-6 量级、BF16 2e-2 且不随层增；FP16 问题版第 4 层 1.1e0，修好 3e-3；BF16 尾数 7 位、FP16 10 位}
\meta{易错}{溢出型错误只在换精度后出现，残差网络里不崩只是某层失效，整张量平均误差看不出来。}
\phantomsection\addcontentsline{toc}{section}{140\quad 算子缺口：盘点与自写}
\lec{140}{算子缺口：盘点与自写}
\onel{在分发器记账：3932 次调用里只有 23 种、2221 次真要在设备上算。}
\begin{itemize}[leftmargin=1.35em,itemsep=2.2pt,topsep=1pt,parsep=0pt]
\item \begingroup\sloppy{\fontsize{10}{14.6}\selectfont 用 TorchDispatchMode 截住每个 aten 调用记次数，并捕获回退警告。\par}\endgroup
\item \begingroup\sloppy{\fontsize{10}{14.6}\selectfont 视图算子只改形状步长、不启 kernel：permute 477、view 448 次。\par}\endgroup
\item \begingroup\sloppy{\fontsize{10}{14.6}\selectfont 只有原子算子构成缺口，复合算子拆开就自动可用。\par}\endgroup
\item \begingroup\sloppy{\fontsize{10}{14.6}\selectfont 缺口集中在三处：量化矩阵乘、融合注意力、MoE 分组算子。\par}\endgroup
\item \begingroup\sloppy{\fontsize{10}{14.6}\selectfont 补法代价：回退 CPU 22273 $\mu$s、组合 1459 $\mu$s、不量化 526 $\mu$s。\par}\endgroup
\item \begingroup\sloppy{\fontsize{10}{14.6}\selectfont 组合与原生误差都 3.78e-2，误差来自 FP8 量化本身。\par}\endgroup
\end{itemize}
\meta{公式}{整体加速 $=\frac{1}{(1-p)+p/s}$，$p$ 是算子时间占比、$s$ 是它自身的加速倍数}
\meta{数字}{3932 次调用、视图 1711 次、要算 2221 次；mm 占 62\%、bmm 9\%、逐元素 8\%；回退慢 15 倍}
\meta{易错}{组合能保住正确性但保不住收益：没有原生 FP8 矩阵乘，FP8 只省显存，不省带宽也不省时间。}
\phantomsection\addcontentsline{toc}{section}{141\quad 实操$\cdot$把 miniinfer 移植到非 CUDA}
\lec{141}{实操$\cdot$把 miniinfer 移植到非 CUDA}
\onel{五步：设备收口、逐 token 比对、算子盘点、逐精度对齐、先治同步点。}
\begin{itemize}[leftmargin=1.35em,itemsep=2.2pt,topsep=1pt,parsep=0pt]
\item \begingroup\sloppy{\fontsize{10}{14.6}\selectfont device.py 收设备 31 行，engine 改 4 处、model 改 3 处。\par}\endgroup
\item \begingroup\sloppy{\fontsize{10}{14.6}\selectfont 设备顺序 CUDA $\rightarrow$ MPS $\rightarrow$ CPU；海光 DCU 也自称 cuda。\par}\endgroup
\item \begingroup\sloppy{\fontsize{10}{14.6}\selectfont FP32 下四条提示全部与 HuggingFace 逐 token 一致。\par}\endgroup
\item \begingroup\sloppy{\fontsize{10}{14.6}\selectfont FP16 在 RMSNorm 平方和溢出，升 FP32 后 3e-3，比 BF16 还小。\par}\endgroup
\item \begingroup\sloppy{\fontsize{10}{14.6}\selectfont int(pos[-1]) 值 13\%：同步 140 $\mu$s 对下发 5.2 $\mu$s。\par}\endgroup
\item \begingroup\sloppy{\fontsize{10}{14.6}\selectfont FP32$\rightarrow$FP16 decode 只快 1.5 倍，带宽 48 对 36 GB/s。\par}\endgroup
\end{itemize}
\meta{数字}{prefill 355$\rightarrow$170$\rightarrow$141 ms；decode 64.9$\rightarrow$55.5$\rightarrow$49.3$\rightarrow$32.9 ms；权重 596 M、FP32 2.38 GB、FP16 1.19 GB}
\meta{易错}{CPU 上没有原生 BF16 矩阵乘，prefill 反比 FP32 慢约 10 倍：低精度在缺原生支持的硬件上更慢。}
\chapter{调优原理：把模型跑快}
{\footnotesize\sffamily\color{grayt}收录第 142--147 讲}\par\vspace{-4pt}
\phantomsection\addcontentsline{toc}{section}{142\quad 在新硬件上做性能分析}
\lec{142}{在新硬件上做性能分析}
\onel{时间线两行有空隙就是下发瓶颈；decode 先治同步，再看访存，最后算力。}
\begin{itemize}[leftmargin=1.35em,itemsep=2.2pt,topsep=1pt,parsep=0pt]
\item \begingroup\sloppy{\fontsize{10}{14.6}\selectfont 看三样：两行关系、最长几个 kernel、空隙的位置与规律。\par}\endgroup
\item \begingroup\sloppy{\fontsize{10}{14.6}\selectfont decode 顺序：下发同步$\rightarrow$访存$\rightarrow$通信$\rightarrow$计算；prefill 反过来。\par}\endgroup
\item \begingroup\sloppy{\fontsize{10}{14.6}\selectfont 第一名是 aten::copy\_：63.3\%、844 次，来自类型转换。\par}\endgroup
\item \begingroup\sloppy{\fontsize{10}{14.6}\selectfont 按段看注意力 47.5\%、FFN 21.6\%；按算子查哪类操作有问题。\par}\endgroup
\item \begingroup\sloppy{\fontsize{10}{14.6}\selectfont 异步后端分段计时必须显式同步，否则记到被阻塞的那次调用上。\par}\endgroup
\item \begingroup\sloppy{\fontsize{10}{14.6}\selectfont 判据两条：实际带宽＝字节÷耗时，算力＝2MNK÷耗时。\par}\endgroup
\end{itemize}
\meta{公式}{实际带宽 $=\frac{\text{读写字节}}{\text{耗时}}$；实际算力 $=\frac{2MNK}{\text{耗时}}$}
\meta{数字}{一步 25.4 ms（FP16、199 token）；每段同步后 61.0 ms；下发小 add 3.6 $\mu$s 对 int(张量) 158.1 $\mu$s（44 倍）}
\meta{易错}{占用率高不等于快；带宽到标称 60\% 以下说明访存不连续、跨步或没对齐。}
\phantomsection\addcontentsline{toc}{section}{143\quad 矩阵乘法与注意力的适配}
\lec{143}{矩阵乘法与注意力的适配}
\onel{decode M 只有几十，大分块补零太多，专用 kernel 快 1.9 倍。}
\begin{itemize}[leftmargin=1.35em,itemsep=2.2pt,topsep=1pt,parsep=0pt]
\item \begingroup\sloppy{\fontsize{10}{14.6}\selectfont 三个硬约束：指令形状与对齐、数据排布、每条通道的累加器数。\par}\endgroup
\item \begingroup\sloppy{\fontsize{10}{14.6}\selectfont 海光小 M kernel：M 在 [1,64]、N 只支持几种、K 必须 128 的倍数。\par}\endgroup
\item \begingroup\sloppy{\fontsize{10}{14.6}\selectfont lm\_head 264 MB：库 0.93 TB/s，只沿 N 分块 1.73 TB/s。\par}\endgroup
\item \begingroup\sloppy{\fontsize{10}{14.6}\selectfont 256$\times$256 在 M=16 补 94\% 的零，8.3 M 格子只有 0.5 M 有用。\par}\endgroup
\item \begingroup\sloppy{\fontsize{10}{14.6}\selectfont 注意力按负载分两种：prefill 与 decode 的并行度来源不同。\par}\endgroup
\item \begingroup\sloppy{\fontsize{10}{14.6}\selectfont 厂商插件给一组后端加选择顺序，每个候选先过配置校验。\par}\endgroup
\end{itemize}
\meta{公式}{实际带宽 $=\frac{\text{权重字节}}{\text{耗时}}$；实际算力 $=\frac{2MNK}{\text{耗时}}$；算术强度 $\approx M$，$M=16$ 时远在拐点左侧}
\meta{数字}{M=16 时 285 对 153 $\mu$s、算力 14.8 对 27.6 TFLOPS；M=64 时 190 $\mu$s；容量 101,376 字节最多 82 行 Q，读写 404 MiB，在线更新 12,800 次}
\meta{易错}{注意力块大小与 KV 分页块大小耦合，不匹配时要么报错要么性能很差；布局还分头优先与序列优先。}
\phantomsection\addcontentsline{toc}{section}{144\quad 用 Triton 与厂商语言写 kernel}
\lec{144}{用 Triton 与厂商语言写 kernel}
\onel{同一个 RMSNorm，融合 kernel 快 4.6 到 10.6 倍。}
\begin{itemize}[leftmargin=1.35em,itemsep=2.2pt,topsep=1pt,parsep=0pt]
\item \begingroup\sloppy{\fontsize{10}{14.6}\selectfont 两条分水岭：要不要写线程编号、搬运是不是显式的。\par}\endgroup
\item \begingroup\sloppy{\fontsize{10}{14.6}\selectfont 框架组合把 RMSNorm 摊成 5 个 kernel，张量读写五遍。\par}\endgroup
\item \begingroup\sloppy{\fontsize{10}{14.6}\selectfont decode 一行 25.9$\rightarrow$5.6 $\mu$s，差的是下发五次变一次。\par}\endgroup
\item \begingroup\sloppy{\fontsize{10}{14.6}\selectfont prefill 199 行 122.2$\rightarrow$11.5 $\mu$s，差的是读写五遍变一遍。\par}\endgroup
\item \begingroup\sloppy{\fontsize{10}{14.6}\selectfont 量带宽要让数据超缓存：800 KB 141 GB/s，32 MB 只有 85 GB/s。\par}\endgroup
\item \begingroup\sloppy{\fontsize{10}{14.6}\selectfont Triton 的折扣：目标架构要有后端，块大小与 num\_warps 要重调。\par}\endgroup
\end{itemize}
\meta{数字}{误差均 5e-8；decode 25.9/22.3/5.6 $\mu$s；prefill 122.2/125.4/11.5 $\mu$s；4096 行 3181.6/394.3 $\mu$s}
\meta{易错}{在组合写法里做微调只值 14\%，写融合 kernel 值 4.6 倍，要改就改到位；手写 kernel 还要跟架构代号走。}
\phantomsection\addcontentsline{toc}{section}{145\quad 低精度与量化的取舍}
\lec{145}{低精度与量化的取舍}
\onel{量化只在反量化融进矩阵乘 kernel 时才省时间，否则只省显存容量。}
\begin{itemize}[leftmargin=1.35em,itemsep=2.2pt,topsep=1pt,parsep=0pt]
\item \begingroup\sloppy{\fontsize{10}{14.6}\selectfont 精度分三档：原生支持、能存不能算、完全不支持；中间档显存减半却不提速。\par}\endgroup
\item \begingroup\sloppy{\fontsize{10}{14.6}\selectfont 有原生 kernel 时 FP16 比 FP32 快 2.51 倍，权重与时间都减半。\par}\endgroup
\item \begingroup\sloppy{\fontsize{10}{14.6}\selectfont W8A16 要把整个权重反量化成新张量，写 34 MB 再读 34 MB，比读 67 MB 更贵。\par}\endgroup
\item \begingroup\sloppy{\fontsize{10}{14.6}\selectfont decode 每步时间下限 = 要读的权重字节 ÷ 带宽。\par}\endgroup
\item \begingroup\sloppy{\fontsize{10}{14.6}\selectfont 只能存不能算时三种精度都挤在 20\textasciitilde{}24 ms，位宽越低这一列越平。\par}\endgroup
\item \begingroup\sloppy{\fontsize{10}{14.6}\selectfont KV Cache 量化要写 KV 和注意力两个 kernel 都支持，常有专用打包布局。\par}\endgroup
\end{itemize}
\meta{公式}{$t=\frac{w}{bw}$（$w=P\times$ 每参数字节）；无原生 kernel：$t_{\text{bad}}=\frac{w+2\cdot P\cdot 2\cdot 2}{bw}$}
\meta{数字}{FP32 1186.6 $\mu$s；FP16 491.4 $\mu$s；BF16 544.7 $\mu$s；INT8 W8A16 4971.1 $\mu$s；W8A8 3922.3 $\mu$s；FP8 2055.4 $\mu$s；3.35 TB/s 下 BF16/FP8/INT4 = 4.79/2.40/1.20 ms，209/417/834 token/s}
\meta{易错}{能存不能算最难识别：模型能加载、显存真降了一半、输出也正确，只有压测速度才发现更慢。}
\phantomsection\addcontentsline{toc}{section}{146\quad 集合通信库与多卡互联}
\lec{146}{集合通信库与多卡互联}
\onel{通信库这层接口几乎不改，要重新量的是 $\alpha$、$\beta$、P2P 和拓扑。}
\begin{itemize}[leftmargin=1.35em,itemsep=2.2pt,topsep=1pt,parsep=0pt]
\item \begingroup\sloppy{\fontsize{10}{14.6}\selectfont 换卡只需重量 $\alpha$ 与 $\beta$：几 KB 小消息量 $\alpha$、几百 MB 大消息量 $\beta$，分界是 $\alpha$$\cdot$$\beta$。\par}\endgroup
\item \begingroup\sloppy{\fontsize{10}{14.6}\selectfont 兼容路线 backend 仍写 nccl、底下换 RCCL；重写路线注册新后端，调用方照旧。\par}\endgroup
\item \begingroup\sloppy{\fontsize{10}{14.6}\selectfont All-Reduce 消息仅 256 KB，Ring 14 轮的固定开销占 97\%。\par}\endgroup
\item \begingroup\sloppy{\fontsize{10}{14.6}\selectfont 自定义 All-Reduce 区间：2/4/6/8 卡、$\leq$8 MiB、全连接、16 字节对齐。\par}\endgroup
\item \begingroup\sloppy{\fontsize{10}{14.6}\selectfont 非全连接时一步到位用不了、跳数乘进 $\alpha$，跨机改用两阶段。\par}\endgroup
\item \begingroup\sloppy{\fontsize{10}{14.6}\selectfont 重叠三招：独立流、按微批次切、把通信融进算子，要在时间线上确认真并行。\par}\endgroup
\end{itemize}
\meta{公式}{$t=\alpha+m/\beta$；小消息分界线 $m=\alpha\beta$；Ring 轮数 $2(p-1)$}
\meta{数字}{NVLink $\alpha$=3 $\mu$s、$\beta$=450 GB/s：8 卡 512 KB 的 Ring 44.0、Tree 20.3、一步到位 11.0 $\mu$s；16 MB：Ring 227.4、一步到位 77.7 $\mu$s；NCCL All-Gather 每次约 55 $\mu$s}
\meta{易错}{自定义 All-Reduce 不是更快的 All-Reduce：卡数、消息、链路、对齐四条之外要么不可用、要么更慢。}
\phantomsection\addcontentsline{toc}{section}{147\quad 推理引擎的硬件插件}
\lec{147}{推理引擎的硬件插件}
\onel{接新硬件靠 entry point 插件：实现平台类与 worker，不改引擎源码。}
\begin{itemize}[leftmargin=1.35em,itemsep=2.2pt,topsep=1pt,parsep=0pt]
\item \begingroup\sloppy{\fontsize{10}{14.6}\selectfont 引擎用 entry point 留入口：平台插件返回平台类全限定名，卡不可用返回 None。\par}\endgroup
\item \begingroup\sloppy{\fontsize{10}{14.6}\selectfont 平台类最少实现五项：设备类型与名字、改配置、注意力后端类、通信器类。\par}\endgroup
\item \begingroup\sloppy{\fontsize{10}{14.6}\selectfont 调度、KV 块管理、前缀缓存、请求生命周期全在上游，插件只回答这块卡怎么做。\par}\endgroup
\item \begingroup\sloppy{\fontsize{10}{14.6}\selectfont 沿用 CUDA 分发键能白拿引擎所有 CUDA 分支；注册新设备则几乎处处要重写。\par}\endgroup
\item \begingroup\sloppy{\fontsize{10}{14.6}\selectfont 改上游行为按代价排：改配置、运行时替换（精确、幂等、失败锁存、可报告）、fork。\par}\endgroup
\item \begingroup\sloppy{\fontsize{10}{14.6}\selectfont 运行时替换必须配版本关卡：只读元数据、主次补丁号一致，rc/dev/本地后缀放宽。\par}\endgroup
\end{itemize}
\meta{数字}{vLLM 17 万行；插件 255134 行，测试 115629 行，补丁 195 个，整模块替换 16 对，平台类 1684 行；nano-vllm 1195 行}
\meta{易错}{入口函数必须可重入且发现阶段不能 import torch/vLLM：一旦提前导入，整模块替换就来不及安排。}
\chapter{案例：海光DCU}
{\footnotesize\sffamily\color{grayt}收录第 148--153 讲}\par\vspace{-4pt}
\phantomsection\addcontentsline{toc}{section}{148\quad 海光 DCU 的软硬件栈}
\lec{148}{海光 DCU 的软硬件栈}
\onel{DCU 走兼容 GPGPU 路线，九层栈逐层有对应物，卡内是波前 64 加 MMAC。}
\begin{itemize}[leftmargin=1.35em,itemsep=2.2pt,topsep=1pt,parsep=0pt]
\item \begingroup\sloppy{\fontsize{10}{14.6}\selectfont DCU 走 GPGPU 架构，编程模型与 CUDA 同构，对外仍自称 cuda 设备。\par}\endgroup
\item \begingroup\sloppy{\fontsize{10}{14.6}\selectfont 九层栈沿用 1 层、翻译 6 层、重写 2 层；重写的是算子库和运行时驱动。\par}\endgroup
\item \begingroup\sloppy{\fontsize{10}{14.6}\selectfont DTK 相当于 CUDA Toolkit，默认装在 /opt/dtk。\par}\endgroup
\item \begingroup\sloppy{\fontsize{10}{14.6}\selectfont 一个波前是 64 条通道（kWaveSize = 64），NVIDIA 的 warp 是 32。\par}\endgroup
\item \begingroup\sloppy{\fontsize{10}{14.6}\selectfont MMAC 形状 16$\times$16$\times$16、BF16 输入 FP32 累加，每通道 4 个累加器。\par}\endgroup
\item \begingroup\sloppy{\fontsize{10}{14.6}\selectfont 工作量一大半在 DAS 层：triton、flash\_attn、deepgemm 等十几个组件。\par}\endgroup
\end{itemize}
\meta{数字}{gfx928 2 处、gfx936 10 处、gfx938 87 处、gfx950 6 处；DTK 26.04/26.10；目标引擎 vLLM 0.28.1；插件 25 万行（提交 81cbe8f）}
\meta{易错}{波前宽度 64 被当作同步语义的依据，不只是性能参数：一个程序就是一个波前，写入才有序。}
\phantomsection\addcontentsline{toc}{section}{149\quad 从 CUDA 到 DCU：迁移与版本矩阵}
\lec{149}{从 CUDA 到 DCU：迁移与版本矩阵}
\onel{CUDA 代码按三类处理迁到 DCU，六层版本必须成套才能起来。}
\begin{itemize}[leftmargin=1.35em,itemsep=2.2pt,topsep=1pt,parsep=0pt]
\item \begingroup\sloppy{\fontsize{10}{14.6}\selectfont 迁移三类：API 与库名机械替换、波前 64 写法人工确认、内联汇编与矩阵指令重写。\par}\endgroup
\item \begingroup\sloppy{\fontsize{10}{14.6}\selectfont 16 行样例里自动 6 行、人工确认 7 行、重写 3 行；hipify 类工具能过九成行。\par}\endgroup
\item \begingroup\sloppy{\fontsize{10}{14.6}\selectfont 写死 32 不报错、只让结果偏小：0xffffffff 掩码和折半归约循环都要按 64 改。\par}\endgroup
\item \begingroup\sloppy{\fontsize{10}{14.6}\selectfont MMAC 内建按代分支：gfx928/936 用 amdgcn 版，gfx938 用 hcu 版。\par}\endgroup
\item \begingroup\sloppy{\fontsize{10}{14.6}\selectfont torch.cuda 能用，因为实现注册在 CUDA 分发键上。\par}\endgroup
\item \begingroup\sloppy{\fontsize{10}{14.6}\selectfont 六层版本必须成套：引擎要求主次补丁号一致、后缀放宽，检查失败要锁存。\par}\endgroup
\end{itemize}
\meta{数字}{目标版本串 0.28.1rc1.dev491+g462fdb097.das.462fdb0.dtk2604；构建 ROCM\_PATH 取代 CUDA\_HOME、sm\_90 变 gfx938；CI：Python 3.10.12/3.12.12、torch 2.11.0+、DTK 26.04/26.10}
\meta{易错}{版本检查只读元数据：通过检查不代表 PyTorch、DTK 和算子库的二进制接口已验证兼容。}
\phantomsection\addcontentsline{toc}{section}{150\quad 读 vllm-plugin-das：平台与运行时补丁}
\lec{150}{读 vllm-plugin-das：平台与运行时补丁}
\onel{接入面只有入口、平台类、补丁机制三处，难点在补丁的四道守卫。}
\begin{itemize}[leftmargin=1.35em,itemsep=2.2pt,topsep=1pt,parsep=0pt]
\item \begingroup\sloppy{\fontsize{10}{14.6}\selectfont 接入面只有三处：入口模块 138 行、平台类 1033 行、补丁机制 195 个文件。\par}\endgroup
\item \begingroup\sloppy{\fontsize{10}{14.6}\selectfont 入口顶层只设算子库环境变量和兼容性检查，不 import torch 与 vLLM，替换才来得及。\par}\endgroup
\item \begingroup\sloppy{\fontsize{10}{14.6}\selectfont 失败锁存：平台初始化错过一次，之后每次调用都抛同一个错误，不在半残注册表上重试。\par}\endgroup
\item \begingroup\sloppy{\fontsize{10}{14.6}\selectfont 补丁三形态：整模块替换 16 对、导入后回调（大多数）、配置改写。\par}\endgroup
\item \begingroup\sloppy{\fontsize{10}{14.6}\selectfont meta\_path 查找器四道守卫：精确模块名、禁重载已打补丁模块、禁重入、失败拒绝导入。\par}\endgroup
\item \begingroup\sloppy{\fontsize{10}{14.6}\selectfont 版本关卡只读元数据，PEP 440 解析器在函数内导入，python -S 下也能用。\par}\endgroup
\end{itemize}
\meta{公式}{$\text{find\_spec}(\text{fullname})\to$ 精确字典查表 $\to$ 整模块替换 / 挂导入后回调 $\to$ 失败锁存并拒绝导入}
\meta{数字}{818 个源文件、255134 行；tests 115629；补丁 195（platform 45 + worker 150）；整模块替换 16 对；平台类 1033 行}
\meta{易错}{supports\_fp8 只有一行：返回这块卡是不是 gfx938，精度支持按架构代号判断而不是型号。}
\phantomsection\addcontentsline{toc}{section}{151\quad DCU 上的矩阵乘：hipBLASLt 与小 M kernel}
\lec{151}{DCU 上的矩阵乘：hipBLASLt 与小 M kernel}
\onel{库在小 M 上退化，于是手写只沿 N 分块、波前沿 K 切的 MMAC kernel。}
\begin{itemize}[leftmargin=1.35em,itemsep=2.2pt,topsep=1pt,parsep=0pt]
\item \begingroup\sloppy{\fontsize{10}{14.6}\selectfont M $\leq$64 时 hipBLASLt 选 256$\times$256 宏块，约 135 $\mu$s 且与 N 无关。\par}\endgroup
\item \begingroup\sloppy{\fontsize{10}{14.6}\selectfont 覆盖 M$\leq$64、N 取 256/512/1024/2048、K 为 128 的倍数。\par}\endgroup
\item \begingroup\sloppy{\fontsize{10}{14.6}\selectfont 一块吃 16 行不补零，波前沿 K 分工、在 LDS 折半归约，并行度靠 N 分块与 K 波前。\par}\endgroup
\item \begingroup\sloppy{\fontsize{10}{14.6}\selectfont 每条通道 4 个 FP32 累加器：16$\times$16 的 256 个元素摊到 64 条通道，是硬件定的。\par}\endgroup
\item \begingroup\sloppy{\fontsize{10}{14.6}\selectfont split-K 变体被宏关掉：DTK 会给 gfx938 生成空的 code object。\par}\endgroup
\item \begingroup\sloppy{\fontsize{10}{14.6}\selectfont 三层回退：前提不满足走 torch.mm，失败只警告一次并永久降级。\par}\endgroup
\end{itemize}
\meta{公式}{tile 为 $16$ 行 $\times\ \text{kNPerBlock}$ 列；$\text{kWarps}$ 个波前沿 $K$ 切，最后在 LDS 里折半归约}
\meta{数字}{移植自 SGLang 的 DCU 分支；lm\_head 分片 32224$\times$4096：M=16 时 285 $\rightarrow$ 153 $\mu$s，M=64 时 190 $\mu$s；hipBLASLt 约 0.93 TB/s，专用 kernel 1.4--1.7 TB/s}
\meta{易错}{库的 135 $\mu$s 与 N 无关：时间和工作量脱钩，代价全在分块的固定开销上，为大矩阵调好的库在小 M 上就这样退化。}
\phantomsection\addcontentsline{toc}{section}{152\quad DCU 上的注意力、量化与通信}
\lec{152}{DCU 上的注意力、量化与通信}
\onel{注意力按优先级挑后端，量化按架构代号，通信按条件分支。}
\begin{itemize}[leftmargin=1.35em,itemsep=2.2pt,topsep=1pt,parsep=0pt]
\item \begingroup\sloppy{\fontsize{10}{14.6}\selectfont 注册五个后端：普通注意力、MLA 各两个版本，另加稀疏 MLA。\par}\endgroup
\item \begingroup\sloppy{\fontsize{10}{14.6}\selectfont 优先级只由三个布尔定：是否 MLA、是否稀疏、是否优先 Triton，可用性交给校验过滤。\par}\endgroup
\item \begingroup\sloppy{\fontsize{10}{14.6}\selectfont 候选全不可用就抛错并打印每个原因，不悄悄降级到能跑但很慢的实现。\par}\endgroup
\item \begingroup\sloppy{\fontsize{10}{14.6}\selectfont 块大小受后端、索引器层（16 token）和压缩比约束，插件用最小公倍数对齐并打日志。\par}\endgroup
\item \begingroup\sloppy{\fontsize{10}{14.6}\selectfont 量化启用 23 种、注释掉 4 种，被删的都和 4 位浮点有关；FP8 按 gfx938 判断。\par}\endgroup
\item \begingroup\sloppy{\fontsize{10}{14.6}\selectfont 通信沿用 nccl 名字和上游 CudaCommunicator；所有卡必须走同一条路径，否则死锁。\par}\endgroup
\end{itemize}
\meta{数字}{后端 5 个：FlashAttention、Triton、FlashMLA、Triton MLA、稀疏 MLA；量化启用 23 种、注释掉 4 种（mxfp4、bitsandbytes、modelopt\_fp4、gpt\_oss\_mxfp4）；NCCL All-Gather 每次约 55 $\mu$s}
\meta{易错}{自设的 --block-size 不生效不是 bug：块大小被注意力后端、KV 布局和压缩比共同约束，插件会改写它。}
\phantomsection\addcontentsline{toc}{section}{153\quad 模型上线 DCU 的完整清单}
\lec{153}{模型上线 DCU 的完整清单}
\onel{DCU 上线按六阶段 29 项逐层点亮，每项都要有能机械判定的验收标准。}
\begin{itemize}[leftmargin=1.35em,itemsep=2.2pt,topsep=1pt,parsep=0pt]
\item \begingroup\sloppy{\fontsize{10}{14.6}\selectfont 六个阶段共 29 项：环境 5、跑通 4、精度 5、性能 6、多卡 4、服务化 5。\par}\endgroup
\item \begingroup\sloppy{\fontsize{10}{14.6}\selectfont 环境关最先：工具看到全部卡、显存温度正常、连跑一小时不掉卡，否则后面结论都不可信。\par}\endgroup
\item \begingroup\sloppy{\fontsize{10}{14.6}\selectfont 跑通关：单卡生成 32 token 不报错；分发器层记账，CPU 回退计数必须为 0。\par}\endgroup
\item \begingroup\sloppy{\fontsize{10}{14.6}\selectfont 精度关：按位置取最差误差，贪心逐 token 与参考平台比，量化模型与同方法参考比。\par}\endgroup
\item \begingroup\sloppy{\fontsize{10}{14.6}\selectfont 性能关：数主机同步次数，decode 与 prefill 形状各测矩阵乘和注意力。\par}\endgroup
\item \begingroup\sloppy{\fontsize{10}{14.6}\selectfont 多卡关：量 $\alpha$、$\beta$ 与拓扑 P2P；两卡逐 token 对单卡，扩到目标卡数验证曲线。\par}\endgroup
\item \begingroup\sloppy{\fontsize{10}{14.6}\selectfont 服务化关：预热覆盖即时编译与图捕获，扫到达率记 TTFT/TPOT，冻结版本并验证回滚。\par}\endgroup
\end{itemize}
\meta{公式}{每步理论下限 $=\frac{\text{每步要读的字节}}{\text{带宽}}$；矩阵乘下限 $=\frac{2MNK}{\text{峰值算力}}$}
\meta{数字}{清单六阶段 29 项；CPU 回退算子慢 15 倍；主机同步 150 $\mu$s 对一次小 kernel 下发 3.6 $\mu$s；即时编译一次几十秒}
\meta{易错}{最容易跳过的三件事：算子回退检查、按精度分别验（FP16 下 RMSNorm 平方和溢出，生成一串句号）、记录取舍理由。}
\clearpage
\chapter*{附录\quad 153 讲一句话速查}
\addcontentsline{toc}{chapter}{附录\quad 153 讲一句话速查}
\markboth{附录：153 讲一句话速查}{}
{\footnotesize\sffamily\color{grayt}每讲只留一句结论，用来快速定位\par}\vspace{4pt}
\begin{multicols}{2}
\raggedright
\noindent{\sffamily\bfseries 1}\;{\small 把事物变成一串有序的数，事物之间的关系就变成可计算的长度与距离。}\par\vspace{2.2pt}
\noindent{\sffamily\bfseries 2}\;{\small 点积等于两长度乘夹角余弦，除掉长度就是只看方向的余弦相似度。}\par\vspace{2.2pt}
\noindent{\sffamily\bfseries 3}\;{\small 矩阵乘法 C=XW 的每一格都是 X 的一行与 W 的一列的点积，整张 C 就是点积表。}\par\vspace{2.2pt}
\noindent{\sffamily\bfseries 4}\;{\small 矩阵就是线性变换，W 的每一行给出一个坐标轴方向被送去的终点。}\par\vspace{2.2pt}
\noindent{\sffamily\bfseries 5}\;{\small 张量是任意维度的数的阵列，读代码先看形状往往就知道运算在做什么。}\par\vspace{2.2pt}
\noindent{\sffamily\bfseries 6}\;{\small softmax 先取指数再除以总和，把任意一组打分变成和为 1 的概率。}\par\vspace{2.2pt}
\noindent{\sffamily\bfseries 7}\;{\small 累积概率把 [0,1) 切成长度等于概率的段，均匀随机数 u 落哪段就选哪个词。}\par\vspace{2.2pt}
\noindent{\sffamily\bfseries 8}\;{\small 导数是局部放大倍数，多层嵌套靠链式法则把各段局部导数连乘。}\par\vspace{2.2pt}
\noindent{\sffamily\bfseries 9}\;{\small 旋转保长度、保夹角；各转 mφ、nφ 时点积只取决于 m-n，即 RoPE。}\par\vspace{2.2pt}
\noindent{\sffamily\bfseries 10}\;{\small 浮点数=符号+指数+尾数，尾数定精度、指数定范围；训练需另存 FP32 主权重。}\par\vspace{2.2pt}
\noindent{\sffamily\bfseries 11}\;{\small 计算量 ÷ 算力 = 理想时间；推理每 token 2N，训练 6N。}\par\vspace{2.2pt}
\noindent{\sffamily\bfseries 12}\;{\small 感知机=点积+偏置+阶跃函数，用一条直线分割平面，学不会线性不可分的异或。}\par\vspace{2.2pt}
\noindent{\sffamily\bfseries 13}\;{\small 神经元=点积+偏置+激活；n 个神经元并排打分就是 y=x$\cdot$W+b。}\par\vspace{2.2pt}
\noindent{\sffamily\bfseries 14}\;{\small 线性层叠多少层都等于一层；插入非线性激活才打破塌缩，两个 ReLU 能算异或。}\par\vspace{2.2pt}
\noindent{\sffamily\bfseries 15}\;{\small MLP 把线性层与激活层层串起来：深度让折线段数指数增长，宽度只能线性增长。}\par\vspace{2.2pt}
\noindent{\sffamily\bfseries 16}\;{\small 拟合数值用 MSE，预测概率用交叉熵 -ln p，它的指数就是困惑度。}\par\vspace{2.2pt}
\noindent{\sffamily\bfseries 17}\;{\small 对每个参数求偏导排成梯度，沿反方向走 η 倍，就能把损失降到最低。}\par\vspace{2.2pt}
\noindent{\sffamily\bfseries 18}\;{\small 一次前向加一次反向，用上游梯度乘局部导数，就能算出所有参数的梯度。}\par\vspace{2.2pt}
\noindent{\sffamily\bfseries 19}\;{\small Momentum 累积速度压震荡，Adam 用梯度平方为每个参数定步长。}\par\vspace{2.2pt}
\noindent{\sffamily\bfseries 20}\;{\small 训练循环五步：取批、前向、损失、反向、更新；batch 与学习率调度决定收敛。}\par\vspace{2.2pt}
\noindent{\sffamily\bfseries 21}\;{\small BPE 反复数相邻 token 对的频次、合并最高的一对，从语料学出子词词表。}\par\vspace{2.2pt}
\noindent{\sffamily\bfseries 22}\;{\small 编号查 V$\times$d 的嵌入表得到向量，等价于独热乘 E，即无偏置的线性层。}\par\vspace{2.2pt}
\noindent{\sffamily\bfseries 23}\;{\small 预训练只做条件概率 P(下一个 token | 前文)，每个位置一道交叉熵题。}\par\vspace{2.2pt}
\noindent{\sffamily\bfseries 24}\;{\small 每步前向得分布、选一个 token 接回输入，循环到结束标记。}\par\vspace{2.2pt}
\noindent{\sffamily\bfseries 25}\;{\small 注意力本身不看顺序，靠给每个位置一个向量加到词嵌入上来补回顺序。}\par\vspace{2.2pt}
\noindent{\sffamily\bfseries 26}\;{\small 每个词用 q 与所有 k 点积打分，softmax 成权重后对 v 加权求和。}\par\vspace{2.2pt}
\noindent{\sffamily\bfseries 27}\;{\small softmax 前把 j>i 的分数设成 -$\infty$，每个位置只能看自己和前面的词。}\par\vspace{2.2pt}
\noindent{\sffamily\bfseries 28}\;{\small 把 Q、K、V 按列切成 n\_h 段，各头独立做注意力，输出拼接后再乘 W\_o。}\par\vspace{2.2pt}
\noindent{\sffamily\bfseries 29}\;{\small 每层输出改成 x+F(x)，导数多一个 1，梯度沿捷径原样传回前面的层。}\par\vspace{2.2pt}
\noindent{\sffamily\bfseries 30}\;{\small 对每个词的向量减均值、除标准差，再乘 γ 加 $\beta$，统一到均值 0 方差 1。}\par\vspace{2.2pt}
\noindent{\sffamily\bfseries 31}\;{\small FFN 对每个词独立做升维$\rightarrow$ReLU$\rightarrow$降维，词间不交流，是参数与计算的大头。}\par\vspace{2.2pt}
\noindent{\sffamily\bfseries 32}\;{\small 一层 = 注意力子层 + FFN 子层，各带残差与归一化，形状不变可叠 N 层。}\par\vspace{2.2pt}
\noindent{\sffamily\bfseries 33}\;{\small 输出层用 d$\times$V 的 lm\_head 加 softmax，给出词表概率。}\par\vspace{2.2pt}
\noindent{\sffamily\bfseries 34}\;{\small 一次前向=查表、位置编码、N 个 Block、lm\_head、softmax。}\par\vspace{2.2pt}
\noindent{\sffamily\bfseries 35}\;{\small 每 token 推理$\approx$2N、训练$\approx$6N FLOPs，参数按结构逐项相加。}\par\vspace{2.2pt}
\noindent{\sffamily\bfseries 36}\;{\small 归一化挪到子层入口成 Pre-Norm，并换成 RMSNorm。}\par\vspace{2.2pt}
\noindent{\sffamily\bfseries 37}\;{\small RoPE 在注意力里按位置旋转 q、k，使 q$\cdot$k 只取决于位置差。}\par\vspace{2.2pt}
\noindent{\sffamily\bfseries 38}\;{\small 门控 FFN 把升维拆成门 g 与内容 u，SiLU(g)$\odot$u 后降维。}\par\vspace{2.2pt}
\noindent{\sffamily\bfseries 39}\;{\small MQA 全部共用一组 K、V，GQA 分组共用，缓存随组数缩小。}\par\vspace{2.2pt}
\noindent{\sffamily\bfseries 40}\;{\small MLA 不缓存 K、V，每 token 只存 512+64=576 个数。}\par\vspace{2.2pt}
\noindent{\sffamily\bfseries 41}\;{\small FFN 拆成 E 个专家，每 token 只挑 k 个算，总参数与计算量解耦。}\par\vspace{2.2pt}
\noindent{\sffamily\bfseries 42}\;{\small 位置靠 YaRN 分治；显存计算靠滑动窗口加 Attention Sink。}\par\vspace{2.2pt}
\noindent{\sffamily\bfseries 43}\;{\small 骨架没变，两条主线：每层更好用、推理更省，KV 压到 68.6 KiB。}\par\vspace{2.2pt}
\noindent{\sffamily\bfseries 44}\;{\small 线性注意力把历史压进 d$\times$d 状态、每步与 n 无关，混合架构补精确召回。}\par\vspace{2.2pt}
\noindent{\sffamily\bfseries 45}\;{\small 网页为主按配比混合，过抽取去重质量安全过滤，数据质量决定能力。}\par\vspace{2.2pt}
\noindent{\sffamily\bfseries 46}\;{\small 损失随 N、D 按幂律降；C$\approx$6ND 时两者一起放大，最优在 U 形底。}\par\vspace{2.2pt}
\noindent{\sffamily\bfseries 47}\;{\small 用提示---回答对继续预测下一 token，模板标角色，掩码只对回答算损失。}\par\vspace{2.2pt}
\noindent{\sffamily\bfseries 48}\;{\small 二选一偏好数据：RLHF 训奖励模型再 PPO，DPO 合成相对参考模型的损失。}\par\vspace{2.2pt}
\noindent{\sffamily\bfseries 49}\;{\small 奖励来自可自动验证的结果，用 RL 让模型自己学会先想再答。}\par\vspace{2.2pt}
\noindent{\sffamily\bfseries 50}\;{\small 冻结 W，只训低秩旁路 A$\cdot$B，可训练参数约 0.52\%。}\par\vspace{2.2pt}
\noindent{\sffamily\bfseries 51}\;{\small 困惑度、基准测试、开放式对战三类尺子各管一段；评测本身是大批量推理。}\par\vspace{2.2pt}
\noindent{\sffamily\bfseries 52}\;{\small 讲义清洗成 499,220 字符，训一个从字符出发、可退回字节的 BPE。}\par\vspace{2.2pt}
\noindent{\sffamily\bfseries 53}\;{\small 长序列按 stride 开窗口，x 前 T 个、y 错一位；两次查表相加。}\par\vspace{2.2pt}
\noindent{\sffamily\bfseries 54}\;{\small 四个 d$\times$d 投影，拆头成 [B,H,T,h\_d]，批量矩阵乘一次算完所有头。}\par\vspace{2.2pt}
\noindent{\sffamily\bfseries 55}\;{\small 按 Pre-Norm 补齐 LN、FFN、残差，组装 530 万参数的 GPT。}\par\vspace{2.2pt}
\noindent{\sffamily\bfseries 56}\;{\small loss = -ln p，起点约 ln V；数据太少时验证 loss 先降后升。}\par\vspace{2.2pt}
\noindent{\sffamily\bfseries 57}\;{\small 存档存配置+参数+优化器状态并留验证最优；解码加温度、Top-k、Top-p。}\par\vspace{2.2pt}
\noindent{\sffamily\bfseries 58}\;{\small GPT-2 权重改名、切 qkv、转置后装进同一个 GPT 类。}\par\vspace{2.2pt}
\noindent{\sffamily\bfseries 59}\;{\small SFT 只把指令位置的 y 设 -100，只在回答与结束符上算 loss。}\par\vspace{2.2pt}
\noindent{\sffamily\bfseries 60}\;{\small 算力来自 132 个 SM 各 4 个 Tensor Core 的矩阵乘加。}\par\vspace{2.2pt}
\noindent{\sffamily\bfseries 61}\;{\small H100 从 HBM 每读 1 字节能做 295 次运算，所以要在片上复用数据。}\par\vspace{2.2pt}
\noindent{\sffamily\bfseries 62}\;{\small 速度上限 min(P, I$\times$B)；H100 屋脊点 295 FLOPs/字节。}\par\vspace{2.2pt}
\noindent{\sffamily\bfseries 63}\;{\small 精度减半则算力翻倍：BF16 989、FP8 1979 TFLOPS。}\par\vspace{2.2pt}
\noindent{\sffamily\bfseries 64}\;{\small 机内 NVLink 450 GB/s、机间 50 GB/s，机内快约 9 倍。}\par\vspace{2.2pt}
\noindent{\sffamily\bfseries 65}\;{\small 超节点把 scale-up 域扩到几十上百卡，All-to-All 受益最大}\par\vspace{2.2pt}
\noindent{\sffamily\bfseries 66}\;{\small 读时间线四步：GPU 空等、按类别分账、找最耗时 kernel、和屋顶比}\par\vspace{2.2pt}
\noindent{\sffamily\bfseries 67}\;{\small 读内存标定带宽、大矩阵乘标定算力，得到真实屋顶再预测 decode 上限}\par\vspace{2.2pt}
\noindent{\sffamily\bfseries 68}\;{\small prefill 算力受限、decode 带宽受限，两种瓶颈决定两套优化}\par\vspace{2.2pt}
\noindent{\sffamily\bfseries 69}\;{\small 只存每层的 k、v，decode 只算新 token，计算量降约 750 倍}\par\vspace{2.2pt}
\noindent{\sffamily\bfseries 70}\;{\small TTFT 看 prefill、TPOT 看 decode；吞吐靠并发，上限看显存}\par\vspace{2.2pt}
\noindent{\sffamily\bfseries 71}\;{\small 权重、KV Cache、激活、开销四项算清，剩下的显存全给 KV Cache}\par\vspace{2.2pt}
\noindent{\sffamily\bfseries 72}\;{\small 每百万 token 成本 = 租金 ÷ 吞吐，输出远比输入贵、靠并发降}\par\vspace{2.2pt}
\noindent{\sffamily\bfseries 73}\;{\small 采样前用温度调尖平、Top-k/Top-p/Min-p 砍长尾，再归一化抽样。}\par\vspace{2.2pt}
\noindent{\sffamily\bfseries 74}\;{\small 调度粒度从「一批请求」缩到「一步迭代」：走一个补一个，批几乎总是满的。}\par\vspace{2.2pt}
\noindent{\sffamily\bfseries 75}\;{\small 把 KV Cache 切成固定块按需分配，块表记录映射，浪费只剩最后一块的空位。}\par\vspace{2.2pt}
\noindent{\sffamily\bfseries 76}\;{\small 给每个装满的块算链式哈希当前缀指纹，命中的块直接复用，跳过这段 prefill。}\par\vspace{2.2pt}
\noindent{\sffamily\bfseries 77}\;{\small 每步设 token 预算，先保 decode，余下按块做 prefill。}\par\vspace{2.2pt}
\noindent{\sffamily\bfseries 78}\;{\small 按需分配块让并发最大，池子满时抢占最后加入的请求，用少量重算换高利用率。}\par\vspace{2.2pt}
\noindent{\sffamily\bfseries 79}\;{\small 不合并适配器：基座所有请求合算，小路各用各的 A、B，代价随不同适配器数涨。}\par\vspace{2.2pt}
\noindent{\sffamily\bfseries 80}\;{\small 每步把接上就不合法的 token logit 置 -$\infty$，输出必符合格式。}\par\vspace{2.2pt}
\noindent{\sffamily\bfseries 81}\;{\small 请求过前端、引擎核心、GPU 三进程九个模块，调度$\rightarrow$前向$\rightarrow$采样每步转一圈。}\par\vspace{2.2pt}
\noindent{\sffamily\bfseries 82}\;{\small 可信压测给的是吞吐-延迟曲线，以及在 SLO 下的最大吞吐 goodput。}\par\vspace{2.2pt}
\noindent{\sffamily\bfseries 83}\;{\small 把几个逐元素 kernel 合成一个，让中间结果留在片上，省掉显存往返。}\par\vspace{2.2pt}
\noindent{\sffamily\bfseries 84}\;{\small 分块算注意力配在线 softmax，S 不落显存，s$^2$ 级读写降到四分之一。}\par\vspace{2.2pt}
\noindent{\sffamily\bfseries 85}\;{\small decode 只有 1 个查询，按块表读散落 KV，再沿上下文切段用满 SM。}\par\vspace{2.2pt}
\noindent{\sffamily\bfseries 86}\;{\small CUDA 用线程、块、网格三级组织，一个块住一个 SM，块内靠共享内存合作。}\par\vspace{2.2pt}
\noindent{\sffamily\bfseries 87}\;{\small 三层分块喂 Tensor Core，块按波分给 SM，波次悬崖与块太少都浪费。}\par\vspace{2.2pt}
\noindent{\sffamily\bfseries 88}\;{\small 以一块数据为单位写 kernel，线程分工与共享内存交给编译器。}\par\vspace{2.2pt}
\noindent{\sffamily\bfseries 89}\;{\small CPU 发射跟不上 GPU 时，把整步 kernel 序列录成图回放即可。}\par\vspace{2.2pt}
\noindent{\sffamily\bfseries 90}\;{\small 整数量化用 scale 把实数映射到格点，误差不超过 scale÷2。}\par\vspace{2.2pt}
\noindent{\sffamily\bfseries 91}\;{\small 权重量化该最小化输出误差而非权重误差，激活大的通道那行最关键。}\par\vspace{2.2pt}
\noindent{\sffamily\bfseries 92}\;{\small 整数矩阵乘只允许激活每 token 一个 scale，得把难度匀给权重。}\par\vspace{2.2pt}
\noindent{\sffamily\bfseries 93}\;{\small 浮点格点在对数坐标上均匀，FP4 只有 8 个值，必须配细粒度分组缩放。}\par\vspace{2.2pt}
\noindent{\sffamily\bfseries 94}\;{\small KV 字节减半则同显存并发翻倍，满载 TPOT 几乎不变、吞吐翻倍。}\par\vspace{2.2pt}
\noindent{\sffamily\bfseries 95}\;{\small 注意力集中在开头几个汇聚点，驱逐必须保留它们，丢了损失暴涨。}\par\vspace{2.2pt}
\noindent{\sffamily\bfseries 96}\;{\small 小模型猜 k 个、大模型一次前向验证，输出与大模型自己生成完全一样。}\par\vspace{2.2pt}
\noindent{\sffamily\bfseries 97}\;{\small Medusa 多头跳猜、EAGLE 特征层逐步猜，树形注意力验证整棵树。}\par\vspace{2.2pt}
\noindent{\sffamily\bfseries 98}\;{\small 蒸馏拟合老师分布、剪枝去权重；只有规则稀疏才在硬件上真变快。}\par\vspace{2.2pt}
\noindent{\sffamily\bfseries 99}\;{\small 一卡不够有三因：放不下、decode 受带宽限、流量扛不住；前两个切分。}\par\vspace{2.2pt}
\noindent{\sffamily\bfseries 100}\;{\small 六种原语归为求和、拼接、交换；All-Reduce 是后两者之和。}\par\vspace{2.2pt}
\noindent{\sffamily\bfseries 101}\;{\small Ring 带宽达下限但步数多，Tree 只需约 2log$_2$p 步赢小消息。}\par\vspace{2.2pt}
\noindent{\sffamily\bfseries 102}\;{\small T = $\alpha$ + m/$\beta$，分界点 m* = $\alpha$$\cdot$$\beta$：小消息卡延迟，大消息卡带宽。}\par\vspace{2.2pt}
\noindent{\sffamily\bfseries 103}\;{\small 数据并行复制完整模型并 All-Reduce 梯度，ZeRO 依次切三块状态。}\par\vspace{2.2pt}
\noindent{\sffamily\bfseries 104}\;{\small BF16 算矩阵乘、FP32 记主参数，16 字节/参数；重计算多算 33\%。}\par\vspace{2.2pt}
\noindent{\sffamily\bfseries 105}\;{\small TP 把矩阵乘法拆到 t 张卡，每卡只读 1/t 权重。}\par\vspace{2.2pt}
\noindent{\sffamily\bfseries 106}\;{\small PP 按层切模型，段间只传激活，气泡占 (p-1)/(m+p-1)。}\par\vspace{2.2pt}
\noindent{\sffamily\bfseries 107}\;{\small SP 切逐 token 运算，CP 把注意力沿序列切开。}\par\vspace{2.2pt}
\noindent{\sffamily\bfseries 108}\;{\small EP 把专家分到各卡，每个 MoE 层两次 All-to-All。}\par\vspace{2.2pt}
\noindent{\sffamily\bfseries 109}\;{\small 通信越频繁越靠里：TP 在 NVLink，PP、DP 放在跨机网络。}\par\vspace{2.2pt}
\noindent{\sffamily\bfseries 110}\;{\small TP 用更多卡换低延迟，但每卡吞吐随卡数下降。}\par\vspace{2.2pt}
\noindent{\sffamily\bfseries 111}\;{\small 大规模 EP 下专家算得快，瓶颈在网络：每 token 每步约 10 MB。}\par\vspace{2.2pt}
\noindent{\sffamily\bfseries 112}\;{\small PD 分离把 prefill 与 decode 分到不同实例。}\par\vspace{2.2pt}
\noindent{\sffamily\bfseries 113}\;{\small 读回已算KV还是重算，看存储带宽是否高于分界带宽，据此分层缓存。}\par\vspace{2.2pt}
\noindent{\sffamily\bfseries 114}\;{\small 路由要兼顾前缀命中与负载均衡，只在负载不超上限的实例里选命中最多的。}\par\vspace{2.2pt}
\noindent{\sffamily\bfseries 115}\;{\small 超节点只改连法：放不进8卡的模型和MoE受益，放得进8卡的更慢。}\par\vspace{2.2pt}
\noindent{\sffamily\bfseries 116}\;{\small 先写SLO，再由TPOT上限与Little定律算单实例请求率，按余量算卡数。}\par\vspace{2.2pt}
\noindent{\sffamily\bfseries 117}\;{\small 引擎骨架是一个循环，按模型执行、状态、调度、内存管理、加速五部分拆开。}\par\vspace{2.2pt}
\noindent{\sffamily\bfseries 118}\;{\small v0用四个文件把1.5GB权重变成前向，59行算出Paris。}\par\vspace{2.2pt}
\noindent{\sffamily\bfseries 119}\;{\small 只缓存每层K、V，新token只算自己的qkv，生成从34s降到4.55s。}\par\vspace{2.2pt}
\noindent{\sffamily\bfseries 120}\;{\small 每步重定batch，新token拼一维、只有注意力分段，吞吐到1.9倍。}\par\vspace{2.2pt}
\noindent{\sffamily\bfseries 121}\;{\small 共享块池加块表取代按最大长度预留，块不够就抢占最晚请求重算。}\par\vspace{2.2pt}
\noindent{\sffamily\bfseries 122}\;{\small 装满且算完的块按链式哈希登记，新请求从头逐块命中即共用。}\par\vspace{2.2pt}
\noindent{\sffamily\bfseries 123}\;{\small 每通道一个 scale 的 INT8 量化，显存降到 1/4 但更慢。}\par\vspace{2.2pt}
\noindent{\sffamily\bfseries 124}\;{\small n-gram 找草稿、一次前向验证，贪心采样下输出与不投机相同。}\par\vspace{2.2pt}
\noindent{\sffamily\bfseries 125}\;{\small CPU 上 76\% 时间在线性层、24\% 在注意力，Python 调度可忽略。}\par\vspace{2.2pt}
\noindent{\sffamily\bfseries 126}\;{\small 拆成调度器、执行器、引擎三段，prefill 优先、一步只做一种。}\par\vspace{2.2pt}
\noindent{\sffamily\bfseries 127}\;{\small 执行侧补 GPU 写法：预热量显存分 KV、变长注意力、分档录图。}\par\vspace{2.2pt}
\noindent{\sffamily\bfseries 128}\;{\small 投机移植的关键是草稿单列、验证复用带前缀缓存的 prefill。}\par\vspace{2.2pt}
\noindent{\sffamily\bfseries 129}\;{\small V1 拆三类进程，调度器只维护「已算」「应有」两个数，每步补差额。}\par\vspace{2.2pt}
\noindent{\sffamily\bfseries 130}\;{\small 服务五层：卡间广播、副本间负载均衡、MoE 锁步、connector 接 PD。}\par\vspace{2.2pt}
\noindent{\sffamily\bfseries 131}\;{\small 每卡一个调度器进程，前缀缓存用基数树，调度与执行重叠即零开销调度。}\par\vspace{2.2pt}
\noindent{\sffamily\bfseries 132}\;{\small 选引擎看它押注在哪一层：kernel、量化、硬件覆盖面还是编译器。}\par\vspace{2.2pt}
\noindent{\sffamily\bfseries 133}\;{\small 带着小引擎的地图进去，只追一个请求的路径；定位功能从配置项搜起，用运行验证。}\par\vspace{2.2pt}
\noindent{\sffamily\bfseries 134}\;{\small 换卡只重做九层里接口下方的实现，nano-vllm 里只有 4.4\% 要改。}\par\vspace{2.2pt}
\noindent{\sffamily\bfseries 135}\;{\small 换卡要重算线程组宽度、指令形状、片上容量三组常数，Roofline 拐点不变。}\par\vspace{2.2pt}
\noindent{\sffamily\bfseries 136}\;{\small CUDA 锁定不在语言，而在算子库、kernel 里的硬件常数、框架外的工具链。}\par\vspace{2.2pt}
\noindent{\sffamily\bfseries 137}\;{\small 接新卡＝在算子分发表里填一列，要么沿用 CUDA 键，要么注册新设备。}\par\vspace{2.2pt}
\noindent{\sffamily\bfseries 138}\;{\small 图中断决定有没有图模式：去掉每层的 int 取值，27 张图降到 2 张。}\par\vspace{2.2pt}
\noindent{\sffamily\bfseries 139}\;{\small 精度目标不是逐位一致，而是每处误差都能归因；逐层按位置取最差才能定位。}\par\vspace{2.2pt}
\noindent{\sffamily\bfseries 140}\;{\small 在分发器记账：3932 次调用里只有 23 种、2221 次真要在设备上算。}\par\vspace{2.2pt}
\noindent{\sffamily\bfseries 141}\;{\small 五步：设备收口、逐 token 比对、算子盘点、逐精度对齐、先治同步点。}\par\vspace{2.2pt}
\noindent{\sffamily\bfseries 142}\;{\small 时间线两行有空隙就是下发瓶颈；decode 先治同步，再看访存，最后算力。}\par\vspace{2.2pt}
\noindent{\sffamily\bfseries 143}\;{\small decode M 只有几十，大分块补零太多，专用 kernel 快 1.9 倍。}\par\vspace{2.2pt}
\noindent{\sffamily\bfseries 144}\;{\small 同一个 RMSNorm，融合 kernel 快 4.6 到 10.6 倍。}\par\vspace{2.2pt}
\noindent{\sffamily\bfseries 145}\;{\small 量化只在反量化融进矩阵乘 kernel 时才省时间，否则只省显存容量。}\par\vspace{2.2pt}
\noindent{\sffamily\bfseries 146}\;{\small 通信库这层接口几乎不改，要重新量的是 $\alpha$、$\beta$、P2P 和拓扑。}\par\vspace{2.2pt}
\noindent{\sffamily\bfseries 147}\;{\small 接新硬件靠 entry point 插件：实现平台类与 worker，不改引擎源码。}\par\vspace{2.2pt}
\noindent{\sffamily\bfseries 148}\;{\small DCU 走兼容 GPGPU 路线，九层栈逐层有对应物，卡内是波前 64 加 MMAC。}\par\vspace{2.2pt}
\noindent{\sffamily\bfseries 149}\;{\small CUDA 代码按三类处理迁到 DCU，六层版本必须成套才能起来。}\par\vspace{2.2pt}
\noindent{\sffamily\bfseries 150}\;{\small 接入面只有入口、平台类、补丁机制三处，难点在补丁的四道守卫。}\par\vspace{2.2pt}
\noindent{\sffamily\bfseries 151}\;{\small 库在小 M 上退化，于是手写只沿 N 分块、波前沿 K 切的 MMAC kernel。}\par\vspace{2.2pt}
\noindent{\sffamily\bfseries 152}\;{\small 注意力按优先级挑后端，量化按架构代号，通信按条件分支。}\par\vspace{2.2pt}
\noindent{\sffamily\bfseries 153}\;{\small DCU 上线按六阶段 29 项逐层点亮，每项都要有能机械判定的验收标准。}\par\vspace{2.2pt}
\end{multicols}
\end{document}